% concatenated from neurips_2024.tex
\documentclass{article}


% if you need to pass options to natbib, use, e.g.:
\PassOptionsToPackage{numbers,sort,compress}{natbib}
% before loading neurips_2023

% ready for submission
 % \usepackage{neurips_2024}


% to compile a preprint version, e.g., for submission to arXiv, add add the
% [preprint] option:
\usepackage[final]{neurips_2024}


% to compile a camera-ready version, add the [final] option, e.g.:
%     \usepackage[final]{neurips_2024}


% to avoid loading the natbib package, add option nonatbib:
%    \usepackage[nonatbib]{neurips_2024}


\usepackage[utf8]{inputenc} % allow utf-8 input
\usepackage[T1]{fontenc}    % use 8-bit T1 fonts
%\usepackage{hyperref}       % hyperlinks
%\usepackage{url}            % simple URL typesetting
\usepackage{booktabs}       % professional-quality tables
\usepackage{amsfonts}       % blackboard math symbols
\usepackage{nicefrac}       % compact symbols for 1/2, etc.
\usepackage{microtype}      % microtypography
\usepackage{xcolor}         % colors

\usepackage{enumitem}
\usepackage{bm}
\usepackage{amsmath}
\usepackage{amsthm}
\usepackage{caption}
\usepackage{amssymb}
\usepackage{subcaption}
\usepackage{wrapfig}
\usepackage{import}

\usepackage[usenames,dvipsnames]{pstricks}
\usepackage{pstricks-add}
\usepackage{epsfig}
\usepackage{pst-grad} % For gradients
\usepackage{pst-plot} % For axes
\usepackage[space]{grffile} % For spaces in paths
\usepackage{etoolbox} % For spaces in paths
\makeatletter % For spaces in paths
\patchcmd\Gread@eps{\@inputcheck#1 }{\@inputcheck"#1"\relax}{}{}
\makeatother


\newtheorem{definition}{Definition}
\newtheorem{theorem}{Theorem}
\newtheorem{remark}{Remark}
\newtheorem{proposition}{Proposition}
\newtheorem{lemma}{Lemma}
\newtheorem{corollary}{Corollary}
\newtheorem{preposition}{Preposition}

\usepackage{tcolorbox}


\def\halp{\hat{\alpha}}

\def\enc{\mathbf{u}_\textrm{enc}}
\def\dec{\mathbf{u}_\textrm{dec}}

\def\encs{u_\textrm{enc}}
\def\decs{u_\textrm{dec}}
\def\encsbar{\bar{u}_\textrm{enc}}
\def\decsbar{\bar{u}_\textrm{dec}}


\def\encsj{u_{\textrm{enc}_j}}
\def\decsj{u_{\textrm{dec}_j}}

\def\enclamb{\lambda_\textrm{e}}
\def\declamb{\lambda_\textrm{d}}

\def\lenc{\mathcal{L}_\textrm{enc}}
\def\ldec{\mathcal{L}_\textrm{dec}}

\def\henc{\mathbf{\hat{u}}_\textrm{enc}}
\def\hdec{\mathbf{\hat{u}}_\textrm{dec}}

\def\hencs{\hat{u}_\textrm{enc}}
\def\hdecs{\hat{u}_\textrm{dec}}

\def\nprcc{\texttt{LeTCC}}
\def\func{\mathbf{f}}
\def\stset{\mathcal{F}}
\def\fhat{\mathbf{\hat{f}}}
\def\est{\fhat_{\bm{\alpha},\bm{\beta}}[\enc, \dec, \stset]}
\def\ests{\fhat_{\bm{\alpha},\bm{\beta}}[\encs, \decs]}


\newcommand\Abs[1]{\left\vert #1 \right \vert}
\newcommand{\norm}[1]{\left\lVert#1\right\rVert}
\newcommand{\lp}[1]{{L}^{#1}\left(\Omega; \mathbb{R}\right)}
\newcommand{\lpm}[1]{{L}^{#1}\left(\Omega; \mathbb{R}^M\right)}
\newcommand{\lploc}[1]{{L}_{\textrm{loc}}^{#1}\left(\Omega; \mathbb{R}\right)}
\newcommand{\lpmloc}[1]{{L}_{\textrm{loc}}^{#1}\left(\Omega; \mathbb{R}^M\right)}
\newcommand{\smp}{\mathbb{W}^{m,p}\left(\Omega; \mathbb{R}^M\right)}
\newcommand{\smpz}{\mathbb{W}_0^{m,p}\left(\Omega; \mathbb{R}^M\right)}
\newcommand{\smpeq}{\widetilde{\mathbb{W}}^{m,p}\left(\Omega; \mathbb{R}^M\right)}
\newcommand{\smploc}{\mathbb{W}_{\textrm{loc}}^{m,p}\left(\Omega; \mathbb{R}^M\right)}
\newcommand{\soblocm}[2]{\mathbb{W}_{\textrm{loc}}^{#1,#2}\left(\Omega; \mathbb{R}^M\right)}
\newcommand{\sobloc}[2]{\mathbb{W}_{\textrm{loc}}^{#1,#2}\left(\Omega; \mathbb{R}\right)}
\newcommand{\sob}[2]{\mathbb{W}^{#1,#2}\left(\Omega; \mathbb{R}\right)}
\newcommand{\sobz}[2]{\mathbb{W}_0^{#1,#2}\left(\Omega; \mathbb{R}\right)}
\newcommand{\sobeq}[2]{\widetilde{\mathbb{W}}^{#1,#2}\left(\Omega; \mathbb{R}\right)}
\newcommand{\sobm}[2]{\mathbb{W}^{#1,#2}\left(\Omega; \mathbb{R}^M\right)}
\newcommand{\sobmz}[2]{\mathbb{W}_0^{#1,#2}\left(\Omega; \mathbb{R}^M\right)}
\newcommand{\sobmeq}[2]{\widetilde{\mathbb{W}}^{#1,#2}\left(\Omega; \mathbb{R}^M\right)}
\newcommand{\hil}[1]{\mathcal{H}^{#1}\left(\Omega;\mathbb{R}\right)}
\newcommand{\hilz}[1]{\mathcal{H}_0^{#1}\left(\Omega;\mathbb{R}\right)}
\newcommand{\hilm}[1]{\mathcal{H}^{#1}\left(\Omega;\mathbb{R}^M\right)}
\newcommand{\hilmz}[1]{\mathcal{H}_0^{#1}\left(\Omega;\mathbb{R}^M\right)}
\newcommand{\hiltilde}[1]{\widetilde{\mathcal{H}}^{#1}\left(\Omega;\mathbb{R}\right)}
\newcommand{\hilmtilde}[1]{\widetilde{\mathcal{H}}^{#1}\left(\Omega;\mathbb{R}^M\right)}

\newcommand{\spline}{\mathbf{S}_{\lambda,n,m}}
% \newcommand{\splineCC}{\mathbf{S}_{\declamb,\stset,2}[\bm{f\circ\encs}]}
\newcommand{\splineCC}{\mathbf{S}_{\declamb,|\stset|,2}}
\newcommand{\lec}[2]{\stackrel{\text{#1}}{#2}}



\def\threedigits#1{%
  \number#1}

\newtcolorbox{mathbox}[1][]{colback=gray!20,  #1}
\tcbset{width=(\linewidth-4mm),
before=,after=\hfill,colframe=black!75!black,colback=black!5!gray}


\title{Coded Computing for Resilient Distributed Computing: A Learning-Theoretic Framework}


% The \author macro works with any number of authors. There are two commands
% used to separate the names and addresses of multiple authors: \And and \AND.
%
% Using \And between authors leaves it to LaTeX to determine where to break the
% lines. Using \AND forces a line break at that point. So, if LaTeX puts 3 of 4
% authors names on the first line, and the last on the second line, try using
% \AND instead of \And before the third author name.


\newcommand{\moh}[1]{{\color{red}#1}}
\interdisplaylinepenalty=2500 % As explained in bare_conf.tex



\author{%
  Parsa Moradi \\
  % Department of Electrical and Computer Engineering\\
  University of Minnesota\\
  \texttt{moradi@umn.edu} \\
  % examples of more authors
  \And
  Behrooz Tahmasebi \\
  MIT CSAIL \\
  % Address \\
  \texttt{bzt@mit.edu} \\
  \And
  Mohammad Ali Maddah-Ali \\
  University of Minnesota \\
  % Address \\
  \texttt{maddah@umn.edu} \\
  % \And
  % Coauthor \\
  % Affiliation \\
  % Address \\
  % \texttt{email} \\
  % \And
  % Coauthor \\
  % Affiliation \\
  % Address \\
  % \texttt{email} \\
}


\usepackage{xcolor}
%\usepackage[numbers]{natbib}

\definecolor{darkblue}{rgb}{0.0,0.0,0.65}
\definecolor{darkred}{rgb}{0.68,0.05,0.0}
\definecolor{darkgreen}{rgb}{0.0,0.29,0.29}
\definecolor{darkpurple}{rgb}{0.47,0.09,0.29}
\usepackage{hyperref}
\hypersetup{
   colorlinks = true,
   citecolor  = darkblue,
   linkcolor  = darkred,
   filecolor  = darkblue,
   urlcolor   = darkblue,
 }
\usepackage{url}
\usepackage{cleveref}



\begin{document}

\maketitle

\begin{abstract}
Coded computing has emerged as a promising framework for tackling significant challenges in large-scale distributed computing, including the presence of slow, faulty, or compromised servers. In this approach, each worker node processes a combination of the data, rather than the raw data itself. The final result then is decoded from the collective outputs of the worker nodes.  However, there is a significant gap between current coded computing approaches and the broader landscape of general distributed computing, particularly when it comes to machine learning workloads. To bridge this gap, we propose a novel foundation for coded computing, integrating the principles of learning theory, and developing a framework that seamlessly adapts with machine learning applications. 
In this framework, the objective is to find the encoder and decoder functions that minimize the loss function, defined as the mean squared error between the estimated and true values. Facilitating the search for the optimum decoding and functions, we show that the loss function can be upper-bounded by the summation of two terms: the generalization error of the decoding function and the training error of the encoding function. 
Focusing on 
the second-order Sobolev space, we then derive the optimal encoder and decoder. We show that in the proposed solution, the mean squared error of the estimation decays with the rate of $\mathcal{O}(S^3 N^{-3})$ and $\mathcal{O}(S^{\nicefrac{8}{5}}N^{\nicefrac{-3}{5}})$ in noiseless and noisy computation settings, respectively, where $N$ is the number of worker nodes with at most $S$ slow servers (stragglers). Finally, we evaluate the proposed scheme on inference tasks for various machine learning models and demonstrate that the proposed framework outperforms the state-of-the-art in terms of accuracy and rate of convergence. 
\end{abstract}


\section{Introduction}
% Coded Computing
The theory of \emph{coded computing} has been developed to improve the reliability and security of large-scale machine learning platforms, effectively tackling two major challenges: (1) the detrimental impact of \emph{slow workers (stragglers)} on overall computation efficiency, and (2) the threat of \emph{faulty or malicious workers} that can compromise data accuracy and integrity.
These challenges have been well-documented in the literature, including the seminal work \cite{dean2013tail} from Google. For instance, \citep{gupta2018oversketch} reported that in a sample set of $3000$ matrix multiplication jobs on AWS Lambda, while the median job time was $40$ seconds, approximately $5\%$ of worker nodes took $100$ seconds to respond, and two nodes took as long as $375$ seconds. Furthermore, coded computing has also been instrumental in addressing \emph{privacy concerns}, a crucial aspect of distributed computing systems~\citep{yu2019lagrange,SoAppxMPC,codedpri,jia2019capacity,aliasgari2019private,kim2019private, chang2019upload, TandonSecMatrix, yu2020straggler,d2020gasp}.


The concept of coded computing has been motivated by the success of coding in communication over unreliable channels, where instead of transmitting raw data, the transmitter sends a (linear) combination of the data, known as coded data. This redundancy in the coded data enables the receiver to recover the raw data even in the presence of errors or missing values. Similarly, coded computing includes three layers \citep{yu2017polynomial,yu2019lagrange,yu2020straggler} (see Figure~\ref{fig:convrate_slopes}(a)):
\begin{enumerate}[label={(\arabic*)}]
\setlength\itemsep{0.1em}
    \item \emph{The Encoding Layer} in which a master node sends a (linear) combination of data,  as coded data, to each worker node. 
    \item \emph{The Computing Layer},
    in which the worker nodes apply a predefined computation to their assigned coded data and send the results back to the master node.

     \item \emph{The Decoding Layer}, in which the master node recovers the final results from the computation results over coded data. In this layer, the decoder leverages the coded redundancy in the computation to recover the missing results of the stragglers and detect and correct the adversarial outputs.  
\end{enumerate}

The existing coded computing has largely built upon algebraic coding theory, drawing inspiration from the renowned Reed-Solomon code construction in communication \citep{reedSolomon}, with proven straggler and Byzantine resiliency \cite{SudanBook}. However, the coding in communication is designed for the exact recovery of the messages, 
 built on a foundation that is inconsistent with the computational requirements of machine learning.
Developing a code that preserves its specific construction while composing with computation is extremely challenging, leading to significant restrictions. Firstly, current methods are mainly restricted to specific computation functions, such as polynomials and matrix multiplication \cite{yu2019lagrange,speed,high,yu2017polynomial,yu2020straggler}. Secondly, rooted in algebraic error correction codes, existing approaches are tailored for finite field computations, leading to numerical instability when dealing with real-valued data \cite{gautschi1987lower,golub2013matrix}.  Furthermore, these methods are unsuitable for approximate, fixed-point, or floating-point computing, where exact computation is neither possible nor necessary, such as in machine learning inference or training tasks. Finally, these schemes typically have a recovery threshold, which is the minimum number of samples required to recover results from coded outputs of worker nodes \cite{yu2019lagrange,yu2017polynomial}. If the number of workers falls below this threshold, the recovery process fails entirely.

Several works have attempted to mitigate the aforementioned issues and transform the coded computing scheme into a more robust and adaptable one, applicable to a wide range of computation functions. These efforts include approximating non-polynomial functions with polynomial ones \cite{so2020scalable,codedpri}, refining the coding mechanism to enhance stability \cite{AnooshehRobust,RamamoorthyCir,RamamoorthyConv,soleymani2020analog,fahim2019numerically}, and leveraging approximation computing techniques to reduce the recovery threshold and increase recovery flexibility \cite{jahani2018codedsketch,overSketch,overSketchN,jahani2022berrut}. However, these attempts fail to bridge the existing gap between coded computing and general distributed computing systems. The root cause of these issues lies in the fact that they are grounded in coding theory, based on a foundation that is not compatible with the requirements of large-scale machine learning. Therefore, this paper aims to address the following objective:


\begin{mathbox}
{\bf Objective:} The main objective of this paper is to develop a new foundation for coded computing, not solely based on \emph{coding theory}, but also grounded in \emph{learning theory}, that seamlessly integrates with machine learning applications, offering a more natural and effective solution for \emph{general computing}.
\end{mathbox}

\begin{figure}[t]
     \centering
     \begin{subfigure}[t]{0.49\textwidth}
         \centering
         \includegraphics[width=\textwidth]{figs/coded_framework.png}
         \caption*{Figure 1(a): Coded Computing: Each worker node processes a combination of data (coded data). The decoder recovers the final results, even in the presence of missing outputs from some worker nodes.}\label{fig:coded_point_framework}
     \end{subfigure}
     \hfill
     \begin{subfigure}[t]{0.49\textwidth}
         \centering
         \includegraphics[width=\textwidth]{figs/conv_rate_det_s_32_k_30.png}
         \caption*{Figure\,1(b): The log-log plot of the expected error versus the number of workers ($N$) for the proposed framework ($\nprcc$) and the state-of-the-art \texttt{BACC} \citep{jahani2022berrut}. $\nprcc$ framework not only achieves a lower estimation error but also has a faster convergence rate.}
         %\label{fig:convrate_slopes}
     \end{subfigure}
        \captionlistentry{}
        \label{fig:convrate_slopes}
\end{figure}

In this paper, we establish a learning-theoretic foundation for coded computing, applicable to general computations. 
We adopt an end-to-end system perspective, that integrates an end-to-end loss function, to find the optimum encoding and decoding functions, focusing on straggler resiliency.
We show that the loss function is upper-bounded by the sum of two terms: one characterizing the \emph{generalization error} of the decoder function and the other capturing the \emph{training error} of the encoder function. Regularizing the decoder layer, we derive the optimal decoder in the Reproducing Kernel Hilbert space (RKHS) of second-order Sobolev functions.
This provides an explicit solution for the optimum decoder function and allows us to characterize the resulting loss of the decoding layer. The decoder loss appears as a regularizing term in optimizing the encoding function and represents the norm in another RKHS. Thus,  the optimum solution for the encoding function can be derived, too.  We address two noise-free and noisy computation settings, for which we derive the \emph{optimal encoder and decoder} and corresponding convergence rate.
We prove that the proposed framework exhibits a faster convergence rate compared to the state-of-the-art and the numerical evaluations support the theoretical derivations (see Figure~\ref{fig:convrate_slopes}(b)). \\

\textbf{Contributions:} The main contributions of this paper are:
\begin{itemize}%[label=(\roman*)]
   \item We develop a new foundation for coded computing by integrating it with learning theory, rather than relying solely on coding theory. We define the loss function as the mean square error of the computation estimation, averaged over all possible sets of at most $S$ stragglers (Section~\ref{sec:objective}). To be able to find the best encoding and decoding functions, we bound the loss function with the summation of two terms, one characterizing the generalization error of the decoder function and the other capturing the training error of the encoder function (Section~\ref{sec:framework}).
    
    \item Assuming that the encoder and decoder functions reside in the Hilbert space of second-order Sobolev functions, we use the theory of RKHSs to find the optimum encoding and decoding functions and characterize the 
    convergence rate for the expected loss in both noise-free and noisy computation regimes  (Section~\ref{sec:mainreults}).
    
    \item We have extensively evaluated the proposed scheme across different data points and computing functions including state-of-the-art deep neural networks and demonstrated that our proposed framework considerably outperforms the state-of-the-art in terms of recovery accuracy (Section~\ref{sec:exp_result}).
\end{itemize}

\section{Preliminaries and Problem Definition}
\subsection{Notations}
Throughout this paper, uppercase and lowercase bold letters denote matrices and vectors, respectively. Coded vectors and matrices are indicated by a $\sim$ sign, as in $\Tilde{\mathbf{x}}, \Tilde{\mathbf{A}}$. The set $\{1, 2, \dots, n\}$ is denoted as $[n]$ and symbol $|S|$ denotes the cardinality of the set $S$. Finally, we represent first, second and $k$-th order derivative of function $f$ as $f',f''$, and $f^{(k)}$, respectively.
\subsection{Problem Setting}
Consider a master node and a set of $N$ workers. The master node is tasked with computing $\{\func(x_k)\}_{k=1}^K$ using a cluster of $N$ worker nodes, given a set of $K$ data points $\{\mathbf{x}_k\}^K_{k=1},  \mathbf{x}_k\in \mathbb{R}^d$. Here, $\func: \mathbb{R}^d \to \mathbb{R}^m$ represents an arbitrary function, which could be a simple one-dimensional function or a complex deep neural network, and $K,d,m$ are integers. A naive approach would be to assign the computation of $\func(\mathbf{x}_k)$ to one worker node for $k \in [K]$. However, some worker nodes may act as stragglers, failing to complete their tasks within the required deadline. To mitigate this issue, the master node employs coding and sends $N$ coded data points to each worker node using an encoder function. Each coded data point is a combination of raw data points. Subsequently, each worker applies the function $\func(\cdot)$ to the received coded data and sends the result, coded results, back to the master node. The master node's goal is to approximately recover $\fhat(\mathbf{x}_k) \approx \func(\mathbf{x}_k)$ using a decoder function, even if some worker nodes appear to be stragglers. The redundancy in the coded data and corresponding coded results enables the master node to recover the desirable results, $\{\func(\mathbf{x}_k)\}^K_{k=1}$. 

% \import{./figs}{figure1.tex}
\begin{figure}[t]
\centering
\includegraphics[width=\textwidth]{figs/figure1.png}
\caption{$\nprcc$ framework.}
\label{fig:framework}
\end{figure}


\section{Proposed Framework: \nprcc}\label{sec:framework}
Here, we propose a novel straggler-resistant Learning-Theoretic Coded Computing ($\nprcc$) framework for general distributed computing. As depicted in Figure~\ref{fig:framework}, our framework comprises two encoding and decoding layers, with a computing layer sandwiched between them. The framework operates according to the following steps:
\begin{enumerate}[label={(\arabic*)}]
\setlength\itemsep{0.1em}
    \item \textbf{Encoding Layer:} The master node fits an encoder regression function $\enc: \mathbb{R} \to \mathbb{R}^d$ at points $\{(\alpha_k, \mathbf{x}_k)\}^K_{k=1}$ for fixed, distinct, and ordered values $\alpha_1 < \alpha_2 < \dots < \alpha_{K} \in \mathbb{R}$. Then, it computes the encoder function $\enc(\cdot)$ on another set of fixed, distinct, and ordered values $\{\beta_n\}_{n=1}^N$ where $\beta_1 < \beta_2 < \dots < \beta_{N} \in \mathbb{R}$. Subsequently, the master node sends the coded data points $\Tilde{\mathbf{x}}_n = \enc(\beta_n) \in \mathbb{R}^d$ to worker $n$ for $n \in [N]$. Note that each coded data point $\Tilde{\mathbf{x}}_n$ is a combination of all initial points $\{\mathbf{x}_k\}^K_{k=1}$.
    \item \textbf{Computing Layer:} Each worker node $n \in [N]$ computes $\func(\Tilde{\mathbf{x}}_n) = \func(\enc(\beta_n))$ on its assigned input and sends the result back to the master node.
    \item \textbf{Decoding Layer:} The master node receives the results $\{\func(\Tilde{\mathbf{x}}_v)\}_{v\in \stset}$ from the non-straggler worker nodes in the set $\stset$. Next, it fits a decoder regression function $\dec: \mathbb{R} \to \mathbb{R}^m$ at points ${(\beta_{v}, \func(\Tilde{\mathbf{x}}_v))}_{v\in \stset}= {(\beta_{v}, \func(\enc(\beta_{v})))}_{v\in \stset}$. Finally, using the function $\dec(\cdot)$, the master node computes $\hat{\func}(\mathbf{x}_k):=\dec(\alpha_k)$ as an approximation of $\func(\mathbf{x}_k)$ for $k \in [K]$. Recall that $\dec(\alpha_k) \approx \func(\enc(\alpha_k)) \approx \func(\mathbf{x}_k)$.
\end{enumerate}

As mentioned above, the master node selects and fixes the regression points, $\{\alpha_k\}^K_{k=1}$ and  $\{\beta_n\}^N_{n=1}$, which remain constant throughout the entire process. The encoder and decoder functions are the only components subject to optimization. 

Note that the computational efficiency of the encoding and decoding layers is crucial. This includes the fitting process of the encoder and decoder regression functions, as well as the computation of these regression functions at points  $\{\beta_v\}_{v\in \stset}$ and $\{\alpha_k\}^K_{k=1}$. If the master node's computation time is not substantially decreased compared to computing $\{\func(\mathbf{x}_k)\}^K_{k=1}$ by itself, then adopting this framework would not provide any benefits for the master node.

\subsection{Objective}\label{sec:objective}
We view the whole scheme as a unified predictive framework that provides an approximate estimation of the values $\left\{\func(\mathbf{x}_k)\right\}^K_{k=1}$. We denote the estimator function of the $\nprcc$ scheme as $\est(\cdot)$, where $\bm{\alpha} := [\alpha_1, \dots, \alpha_K]^T$, $\bm{\beta} := [\beta_1, \dots, \beta_N]^T$, and $\stset:=\{{i_1}, \dots, {i_{|\stset|}}\}$ represents the set of non-straggler worker nodes.

Let $F_{S,N}$ denote a distribution over the collection of subsets of $N$ workers with at most $S$ stragglers, i.e.,
\(
\{\mathcal{F} \subseteq [N] : |\mathcal{F}| \geq N - S\}
\). Also, suppose each worker node $n \in [N]$ computes the function $\func_n(x) = \func(x) + \bm{\epsilon}_n$,  where $\bm{\epsilon}_n$, $n \in [N]$  are independent zero-mean noise vectors with covariance $\sigma^2\mathbf{I}$.

This enables us to define the following loss function, which evaluates the framework's performance:
\begin{align}\label{eq:main_formula}
    \mathcal{R}(\fhat) := \mathop{\mathbb{E}}_{\bm{\epsilon}, \stset \sim 
    F_{S,N}} \left[\frac{1}{K} \sum^K_{k=1} \norm{\fhat(\mathbf{x}_k) - \func(\mathbf{x}_k)}^2_2\right]
=\mathop{\mathbb{E}}_{\bm{\epsilon}, \stset \sim 
F_{S,N}} \left[\frac{1}{K}
    \sum^K_{k=1} \norm{\dec(\alpha_k)- \func(\mathbf{x}_k)}^2_2\right],
\end{align}
where $\fhat(\mathbf{x}) := \est(\mathbf{x})$ to simplify the notation, $\norm{\cdot}_2$ represents  the $\ell_2$-norm, and $\bm{\epsilon} = [\bm{\epsilon}_1,\dots, \bm{\epsilon}_N]^T$. Our objective is to find $\enc(.)$ and $\dec(.)$ that minimize the objective function~\eqref{eq:main_formula}, which is  very challenging, given that $\fhat(.)$ is a composition of $\enc(.)$ and $\dec(.)$ and the computation in the middle. Here, we take an important step to decompose these two, to gain a deeper understanding of interactions.
Adding and subtracting $ \func(\enc(\alpha_k))$ and utilizing inequality of arithmetic and geometric means (AM-GM), one can obtain an upper bound for \eqref{eq:main_formula}:
\begin{alignat}{2}
    \label{eq:decompose}
   \mathcal{R}(\fhat) &= &&~\mathop{\mathbb{E}}_{\bm{\epsilon}, \stset \sim F_{S,N}} \left[\frac{1}{K}
     \sum^K_{k=1} \norm{(\dec(\alpha_k)- \func(\enc(\alpha_k)))+ (\func(\enc(\alpha_k)) - \func(\mathbf{x}_k))}^2_2\right] \nonumber \\
    &\leqslant &&\underbrace{\mathop{\mathbb{E}}_{\bm{\epsilon}, \stset \sim F_{S,N}} \left[\frac{2}{K} \sum^K_{k=1} \norm{\dec \left(\alpha_k\right) - \func \left(\enc\left(\alpha_k\right)\right)}^2_2\right]}_{\ldec(\fhat)} +  \underbrace{\frac{2}{K} \sum^K_{k=1} \norm{\func(\enc(\alpha_k)) - \func(\mathbf{x}_k)}^2_2}_{\lenc(\fhat)}.
\end{alignat}
The right-hand side of \eqref{eq:decompose} comprises two terms, which uncover an interesting interplay between the encoder and decoder regression functions. Let us elaborate on what each term corresponds to.
\begin{itemize}
    \item $\ldec(\fhat)$ -- \textbf{The expected generalization error of the decoder regression:} Recall that the master node fits a decoder regression function, $\dec(\cdot)$, at a set of points denoted as $\left\{(\beta_v, \func(\enc(\beta_v)))\right\}_{v \in \stset}$. $\ldec$ represents the $\ell_2$-norm of the decoder regression function's error on a distinct set of points $\{\alpha_k\}^K_{k=1}$, which are \emph{different} from its training data $\{ \beta_v\}_{v \in \stset}$. Consequently, this term provides an unbiased estimate of the decoder's generalization error. 

    Given that the decoder regression function develops to estimate $\func(\enc(\cdot))$, the generalization error of the decoder regression is inherently tied to the properties of $\func(\enc(\cdot))$. This, in turn, is influenced by characteristics of both the $\func(\cdot)$ and $\enc(\cdot)$ functions, making the $\ldec(\fhat)$ a complex interplay of these two functions.
    \item $\lenc(\fhat)$ -- \textbf{A proxy to the training error of the encoder regression:} Remember that the encoder regression is fitted at points $\{(\alpha_k, \mathbf{x}_k)\}^K_{k=1}$. Consequently, the training error is calculated as $\frac{1}{K} \sum^K_{k=1} \norm{\enc(\alpha_k) - \mathbf{x}_k}^2_2$. Therefore, $\lenc$ represents the encoder training error magnified by the effect of computing function $\func(\cdot)$. Specifically, if $\func(\cdot)$ is $q$-Lipschitz, then $\lenc(\fhat)$ can be upper bounded by:
    \begin{align}\label{eq:lipschitz_enc}
     \frac{2}{K} \sum^K_{k=1} \norm{\func(\enc(\alpha_k)) - \func(\mathbf{x}_k)}^2_2 \leqslant \frac{2q^2}{K} \sum^K_{k=1} \norm{\enc(\alpha_k) - \mathbf{x}_k}^2_2.
    \end{align}
\end{itemize}

\section{Main Results}\label{sec:mainreults}
In this section, we examine the proposed framework from a theoretical standpoint. We provide a comprehensive explanation of the design process for the decoder and encoder functions and subsequently analyze the convergence rate. For simplicity, we present the results for a one-dimensional function $f:\mathbb{R} \to \mathbb{R}$. These results are generalizable to the case where $f:\mathbb{R} \to \mathbb{R}^m$, as discussed in Appendix~\ref{app:hdim_f}.

Suppose the regression points, $\{\alpha_k\}^K_{k=1}, \{\beta_n\}^N_{n=1}$, are confined to the interval $\Omega:=(-1, 1)$ and $\encs, \decs \in \hiltilde{2}$, where $\hiltilde{2}$ is the reproducing kernel Hilbert space (RKHS) of second-order Sobolev functions on the interval $\Omega$ induced with the norm $\norm{f}_{\hiltilde{2}}^2 := \int_{\Omega} (f''(t))^2\,dt +  f(-1)^2 + f'(-1)^2$ which is an equivalent norm on Sobolev space introduced by \citep{kimeldorf1971some} (see \eqref{eq:sob_norm_eq} in Appendix~\ref{app:sobolev}). The definition and properties of Sobolev spaces, along with their reproducing kernels and norms, are reviewed in Appendix~\ref{app:sobolev}.

{\bf Decoder Design:} Since $\ldec(\fhat)$  in the decomposition~\eqref{eq:decompose} characterizes the generalization error of the decoder function, we propose a regularized objective function for the decoder:
\begin{align}\label{eq:decoder_opt}
    \decs^\star=\underset{u \in \hiltilde{2}}{\operatorname{argmin}} \frac{1}{|\stset|} \sum_{v \in \stset}\left(u\left(\beta_v\right)-f\left(\encs\left(\beta_v\right)\right)\right)^2+\declamb \int_{\Omega} \left(u''(t)\right)^2\,dt.
\end{align}
The first term in \eqref{eq:decoder_opt} corresponds to the mean squared error, while the second term characterizes the smoothness of the decoder function on the interval $\Omega$. Equation~\eqref {eq:decoder_opt} represents a Kernel Ridge Regression problem (KRR). It can be shown that the solution of \eqref{eq:decoder_opt} has the following form \citep{duchon1977splines, wahba1990spline}:
\begin{align}\label{eq:dec_sol_form}
     d_0 + d_1t + \sum^{|\stset|}_{v=1} c_v\phi_0(t,\beta_{i_v}),
\end{align}
where $d_0, d_1 \in \mathbb{R}$, $\phi_0(\cdot, \cdot)$ is the kernel function of $\hilz{2}$ (see Definition~\ref{def:sobz} and \eqref{eq:sob:kernel} in Appendix~\ref{app:sobolev}), and $\mathbf{c}=[c_1,\dots,c_{|\stset|}]^T \in \mathbb{R}^{|\stset|}$. Substituting \eqref{eq:dec_sol_form} into the main objective \eqref{eq:decoder_opt}, the coefficient vectors $\mathbf{c}$ and $\mathbf{d}:=[d_0, d_1]^T$ can be efficiently computed by optimizing a quadratic equation \citep{wahba1975smoothing}. This solution is known as the \emph{second-order smoothing spline} function. The theoretical properties of smoothing splines are reviewed in Appendix~\ref{app:smoothspline}.

Let us define the following variables, which represent the maximum and minimum distances between consecutive data points in the decoder layer, $\{\beta_n\}_{n=1}^N$:
\begin{gather}\label{eq:max_dist_st}
\Delta_\textrm{max}:=\underset{n\in \{0\} \cup [N]}{\max} \{\beta_{n+1}-\beta_n\}, \quad \Delta_\textrm{min}:=\underset{n\in [N-1]}{\min} \left\{\beta_{n+1}-\beta_n\right\},
\end{gather}
with $\beta_{0} := -1$ and $\beta_{N+1} := 1$. The following theorems provide crucial insights for designing the encoder function as well as deriving the convergence rates.
\begin{theorem}[Upper bound for noiseless computation,   $\sigma_0 = 0$] \label{th:conv_rate_noiseless}
Consider the $\nprcc$ framework with $N$ worker nodes and at most $S$ stragglers with $\declamb \leqslant N^{-4}$. Assume $\{\alpha_k\}^K_{k=1}$ are arbitrary and distinct points in $\Omega=(-1, 1)$ and there is  constant $B$ such that $\frac{\Delta_{\textrm{max}}}{\Delta_{\textrm{min}}} \leqslant B$. If $f(\cdot)$ is a $q$-Lipschitz continuous function, then:
\begin{align}\label{eq:th2}
    \mathcal{R}(\hat{f}) \leqslant C_1\left(\frac{S+1}{N}\right)^3  \cdot \norm{(f \circ \encs)''}^2_{\lp{2}} + \frac{2q^2}{K} \sum^K_{k=1} (\encs(\alpha_k) - x_k)^2 ,
\end{align}
where $C_1$ is a constant. 
\end{theorem}
The proof of Theorem \ref{th:conv_rate_noiseless} and the detailed expression for $C_1$ can be found in Appendix~\ref{app:proof_noiseless}.
\begin{theorem}[Upper bound for noisy computation]\label{th:conv_rate_noisy}
Consider the $\nprcc$ framework with $N$ worker nodes and at most $S$ stragglers and $\frac{1}{(N-S)^4}\leqslant \declamb \leqslant \lambda_0$ for constant $\lambda_0\in \mathbb{R}$. Assume each worker node computes %$f_{\epsilon_n}(x) = f(x) + \epsilon_n$ 
$f_n(x) = f(x) + \epsilon_n$ 
with $\mathbb{E}[\epsilon_n] = 0$ and $\mathbb{E}[\epsilon_n^2] \leqslant \sigma_0^2$.
Assume $\{\alpha_k\}^K_{k=1}$ are arbitrary and distinct points in $\Omega=(-1, 1)$ and suppose there is  constant $B$ such that $\frac{\Delta_{\textrm{max}}}{\Delta_{\textrm{min}}} \leqslant B$. Assume $f(\cdot)$ is a $q$-Lipschitz continuous function. Then,  
\begin{align}\label{eq:th1}
    \mathcal{R}(\hat{f}) \leqslant 4
    \left(\frac{\sigma^2_0}{N-S}\right)^{\frac{3}{5}} 
    %\cdot 
    \left(C_2\cdot C(\lambda_0)\cdot p_4(S) \cdot
    \norm{(f \circ \encs)''}^{2}_{\lp{2}}\right)^{\frac{2}{5}} + \frac{2q^2}{K} \sum^K_{k=1} (\encs(\alpha_k) - x_k)^2,
\end{align}
if
\begin{align}\label{eq:th2_opt_lambda}
    \frac{1}{(N-S)^\frac{1}{5}\lambda^\frac{1}{4}_0}\leqslant \left(\frac{C_2\cdot C(\lambda_0)\cdot p_4(S)\cdot\norm{(f\circ\encs)^{(2)}}^2_{\lp{2}}}{\sigma^2_0}\right)^\frac{1}{5} \leqslant (N-S)^\frac{4}{5}
    % \frac{\left(\frac{C_2}{4}\right)^\frac{-5}{2}\cdot \sigma_0^2} {p_4(S) \cdot (N-S)\cdot\norm{(f\circ\encs)^{(2)}}^2_{\lp{2}}} \leqslant \lambda_0^\frac{5}{4} \cdot \frac{\max \left\{4, \lambda_0\right\}^\frac{1}{2}}{\max\left\{\lambda_0, C_3 \right\}^\frac{1}{4}},
\end{align}
% \begin{align}\label{eq:th2_opt_lambda}
%     \frac{\left(\frac{C_2}{4}\right)^\frac{-5}{2}\cdot \sigma_0^2} {p_4(S) \cdot (N-S)\cdot\norm{(f\circ\encs)^{(2)}}^2_{\lp{2}}} \leqslant \lambda_0^\frac{5}{4} \cdot \frac{\max \left\{4, \lambda_0\right\}^\frac{1}{2}}{\max\left\{\lambda_0, C_3 \right\}^\frac{1}{4}},
% \end{align}
where $C_2$ is a constant, $C(\lambda_0) = \mathcal{O}(\lambda_0^{\frac{1}{2}})$ is an increasing function of $\lambda_0$, and $p_4(S)$ is a degree-4 polynomial in $S$ with positive constant coefficients.
\end{theorem}

The proof and expressions for $C_2$ and $C(\lambda_0)$ are provided in Appendix~\ref{sec:proof_noisy}.

% {\bf \textcolor{red} {Proof Idea of the Theorems} \ref{th:conv_rate_noiseless}, \ref{th:conv_rate_noisy}:}
% {\it Remark 1.} Theorems \ref{th:conv_rate_noiseless} and \ref{th:conv_rate_noisy} establish convergence rates of $\mathcal{O}(N^{-3})$ and $\mathcal{O}(N^{-\frac{3}{5}})$ for noiseless and noisy computation settings, respectively. Notably, these rates can be achieved using an overfitted encoder, such as a degree-$K$ polynomial that interpolates the points $\{x_k\}^K_{k=1}$ which belongs to the second-order Sobolev space $\mathcal{H}^{2,2}(\Omega)$.\\

% Some results from approximation theory suggest using Chebyshev points of the first kind, $\alpha_k = \cos\left(\frac{(2k-1)\pi}{2K}\right)$, as regression points, as employed in \citep{jahani2022berrut,soleymani2022approxifer}. To address such cases, we present the following corollary.
% \begin{corollary}
% For $\alpha_k = \cos(\frac{(2k-1)\pi}{2K})$, we have the same upper bounds as in Theorems~\ref{th:conv_rate_noisy} and \ref{th:conv_rate_noiseless} with different constants. 
% \end{corollary} 
% The proof can be found in
% Corollary~\ref{cor:dis_to_cont_cheb} in Appendix~\ref{app:theorem_proofs}.\\
% {\it Remark 1.} In \citep{jahani2022berrut}, it is recommended to employ $\alpha_k = \cos(\frac{(2k-1)\pi}{2K})$ and $\beta_n = \cos{\frac{(n-1)\pi}{N-1}}$, which correspond to the Chebyshev points of the first and second kind, respectively. This choice is motivated by the well-established stability of Chebyshev points in polynomial interpolation, as discussed in \citep{salzer1972lagrangian}. By utilizing these Chebyshev points for $\{\alpha_k\}^K_{k=1}$ and $\{\beta_n\}^N_{n=1}$ in $\nprcc$, one can show that this change solely affects the upper bounds \eqref{eq:th1} and \eqref{eq:th2} by a scaling factor that depends on the value of $K$ (see Corollary~\ref{cor:dis_to_cont_cheb} in Appendix~\ref{app:theorem_proofs} for the details).
{\bf Encoder Design:} The upper bounds established in Theorems~\ref{th:conv_rate_noiseless}, 
\ref{th:conv_rate_noisy} hold for all $\encs \in \hiltilde{2}$. However, they do not directly lead to a design for $\encs(\cdot)$. To address this, we present the following theorem which bounds the $\norm{f\circ\encs}^2_{\lp{2}}$, enabling us to construct $\encs(\cdot)$ without compromising the convergence rate.
% {\bf Encoder Design:} The upper bounds provided by Theorems \ref{th:conv_rate_noisy}, 
% \ref{th:conv_rate_noiseless} hold for all $\encs \in \hil{2}$. The two terms in the bounds correspond to the smoothness of the composed function $f\circ \encs(\cdot)$ and the training error of the encoder function, respectively. \moh{THEN WHAT? Connection LOST? }
\begin{theorem}\label{th:optimcal_encdoer} Consider a $\nprcc$ scheme. Assume computing function $f(\cdot)$ is $q$-Lipschitz continuous and
$\norm{f''}_{\lp{\infty}} \leqslant \nu$. Then:
\begin{align}\label{eq:optimcal_encdoer}
    \mathcal{R}(\hat{f}) \leqslant \frac{2q^2}{K} \sum^K_{k=1} (\encs(\alpha_k) - x_k)^2 + \enclamb \cdot   \psi\big(\norm{\encs}^2_{\hiltilde{2}}\big),
\end{align}
for some monotonically increasing function $\psi:\mathbb{R}^{+} \to \mathbb{R}^+$, where 
$\enclamb$ is depending on $(N, S, \sigma_0, q, \nu)$.
\end{theorem}
The proof can be found in Appendix~\ref{app:proof_optimal_enc}. By applying the representer theorem \citep{scholkopf2001generalized}, we can deduce that the optimal encoder $\encs(\cdot)$, which minimizes the right-hand side of \eqref{eq:optimcal_encdoer} takes the form $\encs (\cdot)= \sum^K_{k=1} z_k\phi(\alpha_k, \cdot)$, where $\mathbf{z} \in \mathbb{R}^K$, and $\phi$ is the kernel function of $\hiltilde{2}$, as discussed in Appendix~\ref{app:smoothspline} and in \eqref{eq:smothspline_kernel}. However, due to the non-linearity of $g(\cdot)$, calculating the values of the coefficients $\mathbf{z}$ is challenging. Nevertheless, we demonstrate that the coefficients can be efficiently derived under certain mild assumptions.
\begin{proposition}\label{prop:encoder_optimal} 
In the noiseless case, there exists $M \in \mathbb{R}$ that depends only on $\{\alpha_k\}_{k=1}^K$ and $\{x_k\}_{k=1}^K$, such that:
\begin{itemize}
\item[(i)] If $\norm{\encs}^2_{\hiltilde{2}} \leqslant M$, then:
    \begin{align}
        \mathcal{R}(\hat{f}) \leqslant \widetilde{R}(\encs),
    \end{align}
    where $\widetilde{R}(u)$ is defined as follows:
\begin{align}\label{eq:encoder_linear_upbound}
    \widetilde{R}(u) := \frac{2q^2}{K} \sum^K_{k=1} (u(\alpha_k) - x_k)^2 + \enclamb \cdot (m_1 + m_2M) \left(M + \int_\Omega u''(t)^2\,dt\right),
\end{align}
and $m_1, m_2$ are constants.

\item[(ii)] If  $u^*(\cdot)$ is the  minimizer of \eqref{eq:encoder_linear_upbound}, then $\norm{u^*}^2_{\hiltilde{2}} \leqslant M$.
  \end{itemize}
\end{proposition}

See Appendix~\ref{app:proof_cor_optimal_enc} for the proof. Proposition~\ref{prop:encoder_optimal} states that, under mild assumptions there exists a \emph{smoothing spline} that minimizes the upper bound given in \eqref{eq:encoder_linear_upbound} without changing in convergence rate.

{\bf Convergence Rate:} Using Theorems \ref{th:conv_rate_noiseless}, \ref{th:conv_rate_noisy}, and \ref{th:optimcal_encdoer}, we can derive the convergence rate of the proposed scheme as stated in the following theorem.
\begin{theorem}[Convergence rate]\label{th:conv_rate}
For $\nprcc$ scheme with $N$ worker nodes and a maximum of $S$ stragglers, 
$
 \mathcal{R}(\hat{f}) \leqslant \mathcal{O}
 \left(S^{\frac{8}{5}}N^{-\frac{3}{5}}
 \right)
$ for the noisy computation, and 
$
 \mathcal{R}(\hat{f}) \leqslant \mathcal{O}\left(S^{3}N^{-3}\right)
$ for the noiseless setting.
\end{theorem} 
Refer to Appendix~\ref{app:proof_convrate}
for the proof. Notably, the convergence rate yields from Theorem~\ref{th:conv_rate} surpasses the state-of-the-art Berrut coded computing approach upper bound \citep{jahani2022berrut} (Figure~\ref{fig:convrate_slopes}), both with respect to $S$ and $N$ (see Appendix~\ref{app:comp_berrut_conv_rate} for detailed comparison).

% Theorem~\ref{th:conv_rate} states that for a number of stragglers maximum $S$, the error converges at a rate of $\mathcal{O}(N^{-3})$ in the noiseless setting, which strictly outperforms the state-of-the-art Berrut coded computing approach \citep{jahani2022berrut} (see Figure~\ref{fig:convrate_slopes}). Note that \citep{jahani2022berrut} does not analyze the noisy setting.
% Theorem~\ref{th:conv_rate} states that, for a fixed number of worker nodes $N$, the fewer stragglers present, the lower the approximation error will be. Moreover, the proposed scheme exhibits a significant improvement in the exponent compared to the state-of-the-art Berrut coded computing approach \citep{jahani2022berrut}, which achieves a convergence rate of  $\mathcal{O}\left(S^6N^{-2}\right)$ in the noiseless setting.

% {\it Remark 4.} (Computational complexity) The regression function in the smoothing spline method is a linear model. Prediction on new points, as well as deriving the fitted coefficients, can be computed in $\mathcal{O}(K.d)$ operations in the encoding layer and $\mathcal{O}(|\stset|.m)$ operations in the decoding layer using B-spline basis functions \citep{de1978practical,de2001calculation, hu1986complete, eilers1996flexible}. Consequently, the computational complexity of the proposed scheme is on par with Berrut rational interpolation \citep{jahani2022berrut}.
% \textcolor{red}{
% {\bf Discussion} In the noiseless setting, when the function $f$ is highly non-smooth, using $\declamb = 0$ can lead to overfitting if number of non-straggler workers $|\stset|$ is not sufficiently large, due to the complexity of $f$ which gets worse when it combines with $\encs$ to form the function $f(\encs(.))$. In this case, in contrast to what theoretical results suggest, a non-zero regularization term $\declamb$ plays a crucial role in achieving lower loss.\\
% To address these issues and build upon previous theorems, one can model the noise $\epsilon_n$ such that its variance adapts to the number of non-straggler workers, $\stset$. This adaptive noise modeling approach enables us to better capture the relationships between the number of workers, the complexity of $f$, and the regularization strength.
% \begin{preposition} with the same assumptions as Theorem \ref{th:conv_rate_noisy}, and assuming that $\mathbb{E}\left[\epsilon_n^{2}\right]=\sigma_{0}^{2} N^{-\theta}$ for each worker node $n \in [N]$, then for  $0\leqslant \theta \leqslant 4$ :
% \begin{align}
%  \mathcal{R}(\hat{f}) \leqslant \mathcal{O}\left(N^{-\frac{3}{5}(1+\theta)}\right),    
% \end{align}
% and for $\theta>4$ :
% \begin{align}
%  \mathcal{R}(\hat{f}) \leqslant \mathcal{O}\left(N^{-3}\right).  
% \end{align}
% \end{preposition}
% The proof can be found in Appendix~\ref{}. $\theta=0$ case is equivalent to the noisy computation setting.
% }


% \begin{theorem}\label{th:conv_rate_noise_free}
% Consider the $\nprcc$ scheme with $\{\alpha_i\}^K_{k=1} = \cos\left(\frac{(2k-1)\pi}{2K}\right)$ and $\{\beta_n\}^N_{n=1} = \cos\left(\frac{(n-1)\pi}{N-1}\right)$ and cubic smoothing spline as the decoding regression function which is fitted on $\beta_{i_1} < \beta_2 < \dots < \beta_{i_{|\stset|}}$ with smoothing parameter $\declamb$. Also, assume regression functions $\encs, \decs \in \mathbb{H}^{2, 2}(\Omega)$ with $||\encs||_{\infty,\Omega},||\decs||_{\infty,\Omega} \leqslantL$, and $f:\mathbb{R} \to \mathbb{R}$ a $q_1$-lipchitz function with $||f''||_{\infty,\Omega} \leqslantq_2$, then for every $\Tilde{\epsilon} > 0$ there exist constants $\Tilde{A}_{\Tilde{\epsilon}}, C_d, \Tilde{B}(L)$ such that:
%     \begin{align}
%         L_{\nprcc} &= \frac{1}{K} \sum^K_{k=1} (\decs(\alpha_k)- f(x_k))^2 \\ \nonumber
%         &\leqslant\frac{q_1^2}{K} \sum^K_{k=1} (\encs(\alpha_k) - x_k)^2 + \Tilde{A}_{\Tilde{\epsilon}}.C_d.\declamb(1 + \frac{\Delta^4}{\declamb}) (q_1^2 +q_2^2). \Tilde{g}\bigl(||\encs||^2_{\mathbb{W}^2(\Omega)}\bigr) + \frac{\Tilde{B}}{K} + \Tilde{\epsilon},
%     \end{align}
%     where $\Delta =\sup_{t\in\Omega} \inf_{s\in \stset} |\beta_{s} - t|$ and  $\Tilde{g}:\mathbb{R} \to \mathbb{R}$ is a monotonic function on $[0, \infty)$.
% \end{theorem}
% \subsection{High-dimensional Computing Function}\label{sec:highdim}
% \begin{align}
% \enc = \arg\min_{\mathbf{u} \in \mathbb{W}^2_2[a, b]} &\sum_{k=1}^{K} ||\mathbf{x}_k - \mathbf{u}(\alpha_k)||^2_2 \nonumber \\ &+ \lambda_{\textrm{enc}} \int^b_a ||\mathbf{u}^{\left({2}\right)}(t)||^2_2 \,dt  \label{eq:encoder_obj} \\
% \dec = \arg\min_{\mathbf{u} \in \mathbb{W}^2_2[a, b]} &\sum_{j\in \mathcal{S}} ||\mathbf{\Tilde{x}}_j - \mathbf{u}(\beta_j)||^2_2 \nonumber \\ &+ \lambda_{\textrm{dec}} \int^b_a ||\mathbf{u}^{\left({2}\right)}(t)||^2_2 \,dt \label{eq:decoding_obj}
% \end{align}

% As discussed, the generalization error of decoding regression depends on both its smoothness and the smoothness of encoding regression. We summarize these two smoothness factors as two distinct regularization terms, namely $\lambda_{\textrm{enc}}$ and $\lambda_{\textrm{dec}}$, which should be optimized jointly. Taking all into account, we propose the following pair of optimization statements to achieve encoding and decoding regression functions:
% \begin{align}
% \enc = \arg\min_{\mathbf{u} \in \mathbb{W}^2_2[a, b]} &\sum_{k=1}^{K} ||\mathbf{x}_k - \mathbf{u}(\alpha_k)||^2_2 \nonumber \\ &+ \lambda_{\textrm{enc}} \int^b_a ||\mathbf{u}^{\left({2}\right)}(t)||^2_2 \,dt  \label{eq:encoder_obj} \\
% \dec = \arg\min_{\mathbf{u} \in \mathbb{W}^2_2[a, b]} &\sum_{j\in \stset} ||\mathbf{\Tilde{x}}_j - \mathbf{u}(\beta_j)||^2_2 \nonumber \\ &+ \lambda_{\textrm{dec}} \int^b_a ||\mathbf{u}^{\left({2}\right)}(t)||^2_2 \,dt \label{eq:decoding_obj}
% \end{align}
% where $a<b$ are arbitrary real numbers and $\mathbb{W}^2_2[a,b]$ containing all functions $g(\cdot)$ for which $g(\cdot)$ and $g'(\cdot)$ are absolutely continuous and $\int_{a}^{b} (g''(t))^2 \,dt < \infty$. Both \eqref{eq:encoder_obj} and \eqref{eq:decoding_obj} consist of two terms. The first term is the sum of squared errors at the initial points, and the second term measures the smoothness of the function up to the second derivative. The parameters $\lambda_{\textrm{dec}}$ and $\lambda_{\textrm{enc}}$ control the weight of the smoothing term in each optimization problem: As the $\lambda_{\textrm{enc}} \to 0$ ($\lambda_{\textrm{dec}} \to 0$), the estimation becomes closer to having an exact value at the input points. On the other hand, as $\lambda_{\textrm{enc}} \to \infty$ ($\lambda_{\textrm{dec}} \to \infty$), the objective returns the smoothest function with the minimum square error, resembling the best quadratic polynomial with the minimum square error.

% Unlike Berrut schemes, which predominantly rely on (1) exact interpolation at input points and (2) numerically stable rational interpolation, we relax the exactness constraint. Instead, we put much more weight emphasis on the smoothness of the regression functions, as we have demonstrated its pivotal role in the end-to-end loss.

% \subsection{Smoothing Spline}
% The regression models proposed in \eqref{eq:encoder_obj} and \eqref{eq:decoding_obj} have been extensively studied in the literature. It can be demonstrated that the solution to the optimization problem is unique and equals to the {\bf natural cubic smoothing spline} with knots at $\{\alpha_k\}^K_{k=1}$ and $\{\beta_j\}_{j \in \stset}$, respectively \citep{green1993nonparametric, wahba1975smoothing, de2001calculation}. In the limit where $\lambda_{\textrm{enc}} \to 0$ ($\lambda_{\textrm{dec}} \to 0$), the optimal function becomes a natural spline precisely passing through the data points. Conversely, as $\lambda_{\textrm{enc}} \to \infty$ ($\lambda_{\textrm{dec}} \to \infty$), the encoding (decoding) regression function converges to the best polynomial of degree three.

% More precisely, given $n$ knots at $\{x_i\}^n_{i=1}$ and their corresponding function values $\{f(x_i)\}^n_{i=1}$, the smoothing spline model with smoothing parameter $\lambda$ estimates $\hat{f}(x) = \sum^n_{i=1} \theta_i \eta_i(x)$, where the set of functions $\{\eta_i(\cdot)\}^n_{i=1}$ are predefined basis functions, and $\{\theta_i\}^n_{i=1}$ are the model coefficients. The coefficients of the smoothing spline are computed from the following statement:
% \begin{align}
%     \hat{\mathbf{\theta}} = (\mathbf{M}^T\mathbf{M} + \lambda \mathbf{\Omega})^{-1}\mathbf{M}^T\mathbf{y},
% \end{align}
% where $\mathbf{y} = (f(x_1),\dots, f(x_n))$ and $\mathbf{M}, \mathbf{\Omega} \in \mathbb{R}^{n\times n}$ are defined as follow:
% \begin{align}
%     \mathbf{M}_{ij} = \eta_j(x_i) \quad \text{and} \quad \Omega_{ij} = \int^b_a \eta''_i(t)\eta''_j(t)\,dt
% \end{align}
% The regression function in the smoothing spline method is a linear model. Prediction on new points, as well as deriving the fitted coefficients, can be computed in $\mathcal{O}(K.d)$ operations in the encoding layer and $\mathcal{O}(|\stset|.m)$ operations in the decoding layer using B-spline basis functions \citep{de1978practical,de2001calculation, hu1986complete, eilers1996flexible}. Consequently, the computational complexity of the proposed scheme is on par with Berrut rational interpolation \citep{jahani2022berrut}.




\section{Experimental Results}\label{sec:exp_result}
In this section, we extensively evaluate the proposed scheme across various scenarios. Our assessments involve examining multiple deep neural networks as computing functions and exploring the impact of different numbers of stragglers on the scheme's efficiency. The experiments are run using PyTorch~\citep{paszke2019pytorch} in a single GPU machine. We evaluate the performance of the $\nprcc$ scheme in three different model architectures:
\begin{itemize}
    \item \textbf{Shallow model}: We choose LeNet5 \citep{lecun1998gradient} architecture as a known shallow network with approximately $6\times 10^4$ parameters, trained on the MNIST \citep{lecun2010mnist}.
    \item \textbf{Deep model with low-dimensional output}: In this scenario, we evaluate the proposed scheme when the function is a deep neural network trained on color images in CIFAR-10 \citep{krizhevsky2009learning} dataset. We use the recently introduced RepVGG \citep{ding2021repvgg} network with around $26$ million parameters which was trained on CIFAR-10\footnote{The pre-trained weights can be found  \href{https://github.com/chenyaofo/pytorch-cifar-models}{here}.}.
    \item \textbf{Deep model with high-dimensional output}: Finally, we demonstrate the performance of the $\nprcc$ scheme in a scenario where the input and output of the computing function are high-dimensional, and the function is a relatively large neural network. We consider the Vision Transformer (ViT) \citep{dosovitskiy2020image} as one of the state-of-the-art base neural networks in computer vision for our prediction model, with more than $80$ million parameters (in the base version). The network was trained and fine-tuned on the ImageNet-1K dataset \citep{deng2009imagenet}\footnote{We use the official PyTorch pre-trained ViT network from \href{https://pytorch.org/vision/main/models/generated/torchvision.models.vit_b_16.html}{here}.}.
\end{itemize}
 We use the output of the last softmax layer of each model as the output.
\setlength{\tabcolsep}{5pt}
\begin{table*}[t]
%\captionsetup{font=small}
\caption{Comparison of the proposed framework ($\nprcc$) and the state-of-the-art ($\texttt{BACC}$) in terms of the Root Mean Squared Error (RMSE) and the Relative Accuracy (RelAcc).}
\label{table:perf_mse_racc}
% \vskip 0.15in
\begin{center}
\begin{small}
\begin{sc}
\begin{tabular}{lcc|cc|cc}
    \hline  & \\[-1.5ex]
    ~ %& 
    & \multicolumn{2}{c}{LeNet5} & \multicolumn{2}{c}{RepVGG}  & \multicolumn{2}{c}{ViT} \\
    $(N, K, |\stset|)$
    &  \multicolumn{2}{c}{$(100, 20, 60)$}  & \multicolumn{2}{c}{$(60, 20, 20)$}  & \multicolumn{2}{c}{$(20, 8, 3)$} \\
    \cline{2-3}
    \cline{4-5}
    \cline{6-7}
     & \\[-1.5ex]
     \bf{Method} & \multicolumn{1}{c}{RMSE} & \multicolumn{1}{c}{RelAcc} & \multicolumn{1}{c}{RMSE} & \multicolumn{1}{c}{RelAcc} & \multicolumn{1}{c}{RMSE} & \multicolumn{1}{c}{RelAcc} \\
    \hline  & \\[-1.5ex]
    \texttt{BACC}  & 2.55$\pm$ 0.43   & 0.92$\pm$ 0.04  & 2.44$\pm$ 0.38 & 0.83$\pm$ 0.05 & 0.68$\pm$ 0.13 & 0.90$\pm$ 0.07  \\
    $\nprcc$    & \textbf{2.18$\pm$ 0.51}  &\textbf{0.94$\pm$ 0.04} & \textbf{2.04$\pm$ 0.42}& \textbf{0.87$\pm$ 0.05} & \textbf{0.62$\pm$ 0.11} & \textbf{0.94$\pm$ 0.06} \\
    \bottomrule
\end{tabular}
\end{sc}
\end{small}
\end{center}
% \vskip -0.1in
\end{table*}

\textbf{Hyper-parameters:} The entire encoding and decoding process is the same for different functions, as we adhere to a non-parametric approach. The sole hyper-parameters involved are the two smoothing parameters ($\lambda_{\textrm{enc}}, \lambda_{\textrm{dec}}$) which are determined using cross-validation and greed search over different values of the smoothing parameters.

\textbf{Baseline:} We compare $\nprcc$ with the Berrut approximate coded computing (\texttt{BACC}) introduced by \citep{jahani2022berrut} as the state-of-the-art coded computing scheme for general computing. The $\texttt{BACC}$ framework is used in \citep{jahani2022berrut} for training neural networks and in \citep{soleymani2022approxifer} for inference. Although Berrut coded computing \citep{jahani2022berrut} is the only existing coded computing scheme for general functions, we include a comparison of the proposed framework with the Lagrange coded computing scheme \citep{yu2019lagrange} for polynomial computation in Appendix~\ref{app:comp_lagrange}.


\textbf{Interpolation Points:} We choose Chebyshev points of the first and second kind, $\{\alpha_i\}^K_{k=1} = \cos\left(\frac{(2k-1)\pi}{2K}\right)$ and $\{\beta_n\}^N_{n=1} = \cos\left(\frac{(n-1)\pi}{N-1}\right)$, for fair comparison with \citep{jahani2022berrut}.

\textbf{Evaluation Metrics:} We employ two evaluation metrics to assess the performance of the proposed framework: Relative Accuracy (RelAcc) and Root Mean Squared Error (RMSE). RelAcc is defined as the ratio of the base model's prediction accuracy to the accuracy of the estimated model on the initial data points. RMSE, on the other hand, is our main loss defined in \eqref{eq:main_formula} which measures the empirical average of the root mean square difference over multiple batches of and non-straggler set $\stset$, providing an unbiased estimation of expected mean square error, $\mathbb{E}_{\mathbf{x}\sim \mathcal{X}, \stset}\left[\frac{1}{K}\sum^K_{k=1} \| \func(\mathbf{x}_k) - \fhat(\mathbf{x}_k)\|_2\right]$, for data distribution $\mathcal{X}$.

{\bf Performance Evaluation:}
Table~\ref{table:perf_mse_racc} presents both RMSE and RelAcc metrics side by side.
The results demonstrate that $\nprcc$ outperforms \texttt{BACC} across various architectures and configurations, with an average improvement of 15\%, 17\%, and 9\% in RMSE for LeNet, RepVGG, and ViT architectures, respectively, and a 2\%, 5\%, and 4\% enhancement in RelAcc.\\
In a subsequent analysis, we evaluate the performance of $\nprcc$ in comparison to \texttt{BACC} across a variety of straggler scenarios. For each number of stragglers, $S$, we randomly select $S$ workers to act as stragglers. Both schemes are then run with the same input data points and straggler configurations, and the process is repeated $20$ times. We record the average values of the RelAcc and RMSE metrics, along with their $95\%$ confidence intervals. Figures~\ref{fig:comp_all_better} and \ref{fig:comp_all} illustrate the performance of both schemes across three model architectures. As shown in both figures, the proposed scheme consistently outperforms \texttt{BACC} for nearly all straggler values. In Figure~\ref{fig:comp_all_better}, where $\frac{N}{K}$ is relatively small--indicating a system design without excessive redundancy, which is more practical--the proposed scheme demonstrates even greater improvements in both metrics.

\begin{figure}[t]
     \centering
     \begin{subfigure}[b]{0.32\textwidth}
         \centering
         \captionsetup{justification=centering} 
         \includegraphics[width=\textwidth]{figs/perf_comp/mse_vit_20_14_conf_1e-06_1e-05_95.png}
         \label{fig:comp_all_mse_vit}
     \end{subfigure}
     \hfill
     \begin{subfigure}[b]{0.32\textwidth}
         \centering
         \captionsetup{justification=centering} 
         \includegraphics[width=\textwidth]{figs/perf_comp/mse_repvgg_conf_60_45.png}
         \label{fig:comp_all_mse_cifar}
     \end{subfigure}
     \hfill
     \begin{subfigure}[b]{0.32\textwidth}
         \centering
         \captionsetup{justification=centering} 
         \includegraphics[width=\textwidth]{figs/perf_comp/mse_lenet_comf_100_60.png}
         \label{fig:comp_all_mse_mnist}
     \end{subfigure}
     \begin{subfigure}[b]{0.32\textwidth}
         \centering
         \captionsetup{justification=centering} 
         \includegraphics[width=\textwidth]{figs/perf_comp/racc_vit_20_14_conf_1e-06_1e-05_95.png}
         \label{fig:comp_all_racc_vit}
     \end{subfigure}
     \hfill
     \begin{subfigure}[b]{0.32\textwidth}
         \centering
         \captionsetup{justification=centering} 
         \includegraphics[width=\textwidth]{figs/perf_comp/racc_repvgg_conf_60_45_png.png}
         \label{fig:comp_all_racc_cifar}
     \end{subfigure}
     \hfill
     \begin{subfigure}[b]{0.32\textwidth}
         \centering
         \captionsetup{justification=centering} 
         \includegraphics[width=\textwidth]{figs/perf_comp/racc_lenet_comf_100_60.png}
         \label{fig:comp_all_racc_mnist}
     \end{subfigure}
     %\captionsetup{font=small}
     \caption{Performance comparison of $\nprcc$ and \texttt{BACC} with a $95\%$ confidence interval across a diverse range of stragglers for different models in a low-redundancy regime (smaller $\frac{N}{K}$).}
     \label{fig:comp_all_better}
\end{figure}

{\bf Computational Complexity:} The calculation and inference of smoothing spline coefficients can be performed linearly in the number of regression points by leveraging B-spline basis functions \citep{eilers1996flexible,de2001calculation,hu1986complete}. Consequently, the encoding and decoding processes in $\nprcc$, which involve evaluating new points and calculating the fitted coefficients, have computational complexities of $\mathcal{O}(K \cdot d)$ and $\mathcal{O}((N - S) \cdot m)$, respectively, where $d$ is the input dimension and $m$ is the output dimension of the computing function $\func(\cdot)$. This complexity is comparable to that of \texttt{BACC}, which has complexities of $\mathcal{O}(K)$ and $\mathcal{O}(N - S)$ for its encoding and decoding layers \citep{jahani2022berrut}. Table~\ref{tab:complexity} in Appendix~\ref{app:comp_berrut_complexity} presents a comparison of the total end-to-end processing time statistics for the $\nprcc$ and \texttt{BACC} schemes.


{\bf Sensitivity Analysis:} We additionally investigate the sensitivity of the proposed scheme’s performance to the value of the smoothing parameter, as well as the sensitivity of the optimal smoothing parameter to the number of stragglers (or workers). The results are presented in Appendix~\ref{app:sens_analysis}. As shown in Table~\ref{tab:sensitivity_straggler} and Figure~\ref{fig:sm_sens}, the optimal smoothing parameters and the scheme's performance exhibit low sensitivity to the number of stragglers (or worker nodes) and to the smoothing parameters, respectively.



{\bf Coded Points:}
We also compare the coded points $\{\mathbf{\Tilde{x}_n}\}_{n=1}^N$ sent to the workers in $\nprcc$ and \texttt{BACC} schemes. The results, shown in Figure~\ref{fig:coded_samples_all}, demonstrate that \texttt{BACC} coded samples exhibit high-frequency noise which causes the scheme to approximate the original prediction worse than $\nprcc$.
\section{Related Work}
Coded computing was initially introduced to tackle the challenges of distributed computation, particularly the existence of stragglers or slow workers, and also faulty or adversarial nodes.   Traditional approaches to deal with stragglers primarily rely on repetition \citep{zaharia2008improving, dean2013tail,suresh2015c3,shah2015redundant,gardner2015reducing,chaubey2015replicated}, where each task is assigned to multiple workers either proactively or reactively.  Recently, coded computing approaches have reduced the overhead of repetition  by leveraging coding theory and embedding redundancy in the worker's input data \citep{jahani2022berrut,yu2019lagrange,yu2017polynomial,jahani2018codedsketch,li2018near,das2019random,ramamoorthysurvey,karakus2017straggler,Grad_Dimakis}.  This technique, which mainly relies on theory of coding, has been developed 
for specific types of structured computations, such as polynomial computation and matrix multiplication \citep{yu2017polynomial, yu2019lagrange, ramamoorthy2021numerically,high,entangle,das2019random,d2020gasp,aliasgari2019private,fahimappx}. 
Recently, there have been attempts to generalize coded computing for general computations \citep{jahani2022berrut,kosaian2019parity,SoAppxMPC,codedpri}. Towards extending the application of coding computing to machine learning computation, \citet{kosaian2019parity} suggest training a neural network to predict coded outputs from coded data points. However, the scheme of \citet{kosaian2019parity} requires a complex training process and tolerates only one straggler. In another work,  \citet{jahani2022berrut} proposes \texttt{BACC}, a model-agnostic and numerically stable framework for general computations. They successfully employed \texttt{BACC} to train neural networks on a cluster of workers, while tolerating a larger number of stragglers. Building on the \texttt{BACC} framework, \citet{soleymani2022approxifer}
introduced \texttt{ApproxIFER} scheme, as a straggler resistance and Byzantine-robust prediction serving system. However, the scheme of~\texttt{BACC} uses a reasonable rational interpolation (Berrut interpolation \citep{berrut1988rational}), off the shelf, 
for encoding and decoding, without considering any end-to-end cost function to optimize. In contrast, we theoretically formalize a new foundation of coded computing grounded in learning theory, which can be naturally used for machine learning applications. 
\begin{figure}[t]
     \centering
     \begin{subfigure}[b]{0.32\textwidth}
         \centering
         \captionsetup{justification=centering} 
         \includegraphics[width=\textwidth]{figs/perf_comp/mse_vit_20_8_conf_1e-06_0.001_95.png}
         \label{fig:comp_all_mse_vit_highred}
     \end{subfigure}
     \hfill
     \begin{subfigure}[b]{0.32\textwidth}
         \centering
         \captionsetup{justification=centering} 
         \includegraphics[width=\textwidth]{figs/perf_comp/mse_repvgg_conf_60_20.png}
         \label{fig:comp_all_mse_cifar_highred}
     \end{subfigure}
     \hfill
     \begin{subfigure}[b]{0.32\textwidth}
         \centering
         \captionsetup{justification=centering} 
         \includegraphics[width=\textwidth]{figs/perf_comp/mse_lenet_comf_100_20.png}
         \label{fig:comp_all_mse_mnist_highred}
     \end{subfigure}
     \begin{subfigure}[b]{0.32\textwidth}
         \centering
         \captionsetup{justification=centering} 
         \includegraphics[width=\textwidth]{figs/perf_comp/racc_vit_20_8_conf_1e-06_0.001_95.png}
         \label{fig:comp_all_racc_vit_highred}
     \end{subfigure}
     \hfill
     \begin{subfigure}[b]{0.32\textwidth}
         \centering
         \captionsetup{justification=centering} 
         \includegraphics[width=\textwidth]{figs/perf_comp/racc_repvgg_conf_60_20_png.png}
         \label{fig:comp_all_racc_cifar_highred}
     \end{subfigure}
     \hfill
     \begin{subfigure}[b]{0.32\textwidth}
         \centering
         \captionsetup{justification=centering} 
         \includegraphics[width=\textwidth]{figs/perf_comp/racc_lenet_comf_100_20.png}
         \label{fig:comp_all_racc_mnist_highred}
     \end{subfigure}
     %\captionsetup{font=small}
     \caption{Performance comparison of $\nprcc$ and \texttt{BACC} with a $95\%$ confidence interval across a diverse range of stragglers for different models in a high-redundancy regime (larger $\frac{N}{K}$).}
     \label{fig:comp_all}
\end{figure}





% However, these approaches have mainly relied on coding theory, which does not align with the requirements of machine learning applications. In contrast with these works, we theoretically formalize a new foundation of coded computing grounded in learning theory, while leveraging the benefits of coding theory. 



% Coded computing was initially introduced to tackle the challenges of distributed computation for specific types of computations, such as polynomial computation and matrix multiplication \citep{yu2017polynomial, yu2019lagrange, ramamoorthy2021numerically,high,entangle,das2019random,d2020gasp,aliasgari2019private,fahimappx}. Recently, there have been attempts to generalize coded computing for general computations \citep{jahani2022berrut,kosaian2019parity,SoAppxMPC,codedpri}. However, these approaches have mainly relied on coding theory, which does not align with the requirements of machine learning applications. In contrast with these works, we theoretically formalize a new foundation of coded computing grounded in learning theory, while leveraging the benefits of coding theory.

% {\bf Straggler Resistant Distributed Computing.} Mitigating stragglers or slow workers, a significant challenge in distributed computing systems, has been extensively studied in recent years. Traditional approaches primarily rely on repetition \citep{zaharia2008improving, dean2013tail,suresh2015c3,shah2015redundant,gardner2015reducing,chaubey2015replicated}, where each task is assigned to multiple workers either proactively or reactively. 
% %In proactive approaches, a task is reassigned to another worker if the progress rate of the original worker falls below a certain threshold \citep{zaharia2008improving,dean2013tail}. In contrast, reactive methods assign each task to multiple workers from the outset. However, this approach has significant drawbacks: assigning the same task to multiple workers incurs excessive resource overhead on the system, and a large number of worker nodes are required to guarantee straggler resiliency. 
% Recently, coded computing approaches have reduced this overhead by leveraging coding theory and embedding redundancy in the worker's input data \citep{jahani2022berrut,yu2019lagrange,yu2017polynomial,jahani2018codedsketch,li2018near,das2019random,ramamoorthysurvey,karakus2017straggler,Grad_Dimakis}.

% {\bf Approximate Coded Computing.} Towards extending the application of coding computing to machine learning computation, \citet{kosaian2019parity} suggest training a neural network to predict coded outputs from coded data points. However, the scheme of \citet{kosaian2019parity} requires a complex training process and tolerates only one straggler. In another work,  \citet{jahani2022berrut} proposes \texttt{BACC}, a model-agnostic and numerically stable framework for general computations. They successfully employed \texttt{BACC} to train neural networks on a cluster of workers, while tolerating a larger number of stragglers. Building on the \texttt{BACC} framework, \citet{soleymani2022approxifer}
% introduced \texttt{ApproxIFER} scheme, as a straggler resistance and Byzantine-robust prediction serving system. However, the scheme of~\texttt{BACC} uses a reasonable rational interpolation (Berrut interpolation \citep{berrut1988rational}), off the shelf, 
% for encoding and decoding, without considering any end-to-end cost function to optimize.   



% Mitigating stragglers or slow workers is one of the main challenges of distributed computing systems, such as prediction serving systems, and has been extensively studied in recent years. Traditional approaches rely on repetition, as seen in works such as \citep{zaharia2008improving, dean2013tail}, where a worker's task is reassigned to another worker only if its progress rate in completing that task falls below a certain threshold. In recent years, coding-based techniques, termed coded computing, have emerged as a promising approach to enhancing the efficiency of distributed computing systems. These methods offer significant improvements in latency and overhead by mitigating the impact of stragglers \citep{lee2017speeding, kosaian2018learning, yu2019lagrange, narra2019distributed, wang2019erasurehead, wang2019fundamental}. In the coded computing approach, each node is given a combination of inputs, rather than the raw data, to process. The objective is to design the code in such a way that the results can be recovered from any set of available results on coded data while ignoring the results of stragglers.
 % Initially, coded computing approaches were designed to exploit the structure of computations, limiting their applications to highly structured tasks like matrix multiplication and polynomial computation. The other disadvantage is that the required number of workers becomes prohibitively large.
% Some effort has been made to extend those ideas to include non-polynomials by approximating the function with polynomials~\citep{so2020scalable}. Still, the scope of this approach is very limited. Subsequently, learning-based methods have been introduced to reduce the number of necessary workers required to mitigate stragglers \citep{kosaian2019parity, kosaian2020learning, dinh2021coded}. However, these models need to train a specific neural network for each computing function. In addition, that approach offers resistance to only a very small number of stragglers. More recently, a numerically stable, model-agnostic, and straggler-resilient method based on Berrut rational interpolation and approximate estimation of the original function at the input data points was introduced by \citep{jahani2022berrut}. In the scheme of~\citep{jahani2022berrut} the main effort was to choose a pole-free rational function that exactly interpolation at initial data points to have exact and numerically stable interpolation. The approach employed in this paper is based on assessing the end-to-end performance to derive design criteria that place value on the smoothness of the decoding and encoding functions in an intriguing manner.

\section{Conclusions and Future Work}
In this paper, we developed a new foundation for coded computing based on learning theory%with the loss function as the mean squared estimation error.
, contrasting with existing works that rely on coding theory and use metrics like minimum distance and recovery threshold for design.
This shift in foundations removes barriers to using coded computing for machine learning applications, allows us to design optimal encoding and decoding functions, and achieves convergence rates that outperform the state of the art. Moreover, the experimental evaluations validate the theoretical guarantees. While this work focuses on straggler mitigation, future work will extend our proposed scheme to achieve Byzantine robustness and privacy, offering promising avenues for further research.

\section{Acknowledgments}
This material is based upon work supported by the National
Science Foundation under Grant CIF-2348638. Behrooz Tahmasebi is supported by NSF Award CCF-2112665 (TILOS AI Institute) and NSF Award 2134108.


% In this paper, we have introduced a novel learning-theoretic foundation for coded computing and applied it to address the straggler mitigation issue. We have defined the end-to-end loss in the presence of stragglers and derived corresponding upper bounds and convergence rates for noisy and noiseless computations. Additionally, we design optimal encoder and decoder functions to achieve these upper bounds. Our experimental evaluation demonstrates significant improvements over the state-of-the-art scheme in various scenarios. Future work includes extending our proposed scheme to achieve Byzantine robustness and privacy, offering promising avenues for further research.

% In this paper, we have introduced $\nprcc$, a general and resilient coded computing framework for prediction serving systems. Our framework consists of two regression layers and a central computing layer, sandwiched between the two regression layers. In contrast to all existing coded computing schemes, where the primary focus is on exact interpolation, we prioritize the final accuracy of recovering original predictions as our main objective. We have stated that the overall objective reveals an underlying interplay between encoding and decoding regression, and we have proposed a technique to solve the mentioned coupled regression problem. Our experiments on MNIST, CIFAR-10, and ImageNet-1K datasets, using various model architectures such as LeNet, RepVGG, and ViT, show that $\nprcc$ and its model-agnostic version ($\nprcc^{Ag}$) outperform the state-of-the-art (\texttt{BACC}) in a wide range of straggler numbers, as well as different metrics, by up to 23\%.
\begin{thebibliography}{66}
\providecommand{\natexlab}[1]{#1}
\providecommand{\url}[1]{\texttt{#1}}
\expandafter\ifx\csname urlstyle\endcsname\relax
  \providecommand{\doi}[1]{doi: #1}\else
  \providecommand{\doi}{doi: \begingroup \urlstyle{rm}\Url}\fi

\bibitem[Dean and Barroso(2013)]{dean2013tail}
Jeffrey Dean and Luiz~Andr{\'e} Barroso.
\newblock The tail at scale.
\newblock \emph{Communications of the ACM}, 56\penalty0 (2):\penalty0 74--80, 2013.

\bibitem[Gupta et~al.(2018{\natexlab{a}})Gupta, Wang, Courtade, and Ramchandran]{gupta2018oversketch}
Vipul Gupta, Shusen Wang, Thomas Courtade, and Kannan Ramchandran.
\newblock {Oversketch}: Approximate matrix multiplication for the cloud.
\newblock In \emph{2018 IEEE International Conference on Big Data (Big Data)}, pages 298--304. IEEE, 2018{\natexlab{a}}.

\bibitem[Yu et~al.(2019)Yu, Li, Raviv, Kalan, Soltanolkotabi, and Avestimehr]{yu2019lagrange}
Qian Yu, Songze Li, Netanel Raviv, Seyed Mohammadreza~Mousavi Kalan, Mahdi Soltanolkotabi, and Salman~A Avestimehr.
\newblock Lagrange coded computing: Optimal design for resiliency, security, and privacy.
\newblock In \emph{The 22nd International Conference on Artificial Intelligence and Statistics}, pages 1215--1225. PMLR, 2019.

\bibitem[So et~al.(2020{\natexlab{a}})So, Guler, and Avestimehr]{SoAppxMPC}
Jinhyun So, Basak Guler, and Amir~Salman Avestimehr.
\newblock A scalable approach for privacy-preserving collaborative machine learning.
\newblock In \emph{Advances in Neural Information Processing Systems}, volume~33, pages 8054--8066, 2020{\natexlab{a}}.

\bibitem[So et~al.(2021)So, G{\"u}ler, and Avestimehr]{codedpri}
Jinhyun So, Ba{\c{s}}ak G{\"u}ler, and Amir~Salman Avestimehr.
\newblock {CodedPrivateML}: A fast and privacy-preserving framework for distributed machine learning.
\newblock \emph{IEEE Journal on Selected Areas in Information Theory}, 2\penalty0 (1):\penalty0 441--451, 2021.

\bibitem[Jia and Jafar(2021)]{jia2019capacity}
Zhuqing Jia and Syed~Ali Jafar.
\newblock On the capacity of secure distributed batch matrix multiplication.
\newblock \emph{IEEE Transactions on Information Theory}, 67\penalty0 (11):\penalty0 7420--7437, 2021.

\bibitem[Aliasgari et~al.(2020)Aliasgari, Simeone, and Kliewer]{aliasgari2019private}
Malihe Aliasgari, Osvaldo Simeone, and J{\"o}rg Kliewer.
\newblock Private and secure distributed matrix multiplication with flexible communication load.
\newblock \emph{IEEE Transactions on Information Forensics and Security}, 2020.

\bibitem[Kim and Lee(2019)]{kim2019private}
Minchul Kim and Jungwoo Lee.
\newblock Private secure coded computation.
\newblock In \emph{2019 IEEE International Symposium on Information Theory (ISIT)}, pages 1097--1101. IEEE, 2019.

\bibitem[Chang and Tandon(2019)]{chang2019upload}
Wei-Ting Chang and Ravi Tandon.
\newblock On the upload versus download cost for secure and private matrix multiplication.
\newblock In \emph{2019 IEEE Information Theory Workshop (ITW)}, pages 1--5. IEEE, 2019.

\bibitem[Chang and Tandon(2018)]{TandonSecMatrix}
Wei-Ting Chang and Ravi Tandon.
\newblock On the capacity of secure distributed matrix multiplication.
\newblock In \emph{2018 IEEE Global Communications Conference (GLOBECOM)}, pages 1--6, 2018.

\bibitem[Yu et~al.(2020{\natexlab{a}})Yu, Maddah-Ali, and Avestimehr]{yu2020straggler}
Qian Yu, Mohammad~Ali Maddah-Ali, and A~Salman Avestimehr.
\newblock Straggler mitigation in distributed matrix multiplication: Fundamental limits and optimal coding.
\newblock \emph{IEEE Transactions on Information Theory}, 66\penalty0 (3):\penalty0 1920--1933, 2020{\natexlab{a}}.

\bibitem[D’Oliveira et~al.(2020)D’Oliveira, El~Rouayheb, and Karpuk]{d2020gasp}
Rafael~G.L. D’Oliveira, Salim El~Rouayheb, and David Karpuk.
\newblock {GASP} codes for secure distributed matrix multiplication.
\newblock \emph{IEEE Transactions on Information Theory}, 66\penalty0 (7):\penalty0 4038--4050, 2020.

\bibitem[Yu et~al.(2017)Yu, Maddah-Ali, and Avestimehr]{yu2017polynomial}
Qian Yu, Mohammad Maddah-Ali, and Salman Avestimehr.
\newblock Polynomial codes: an optimal design for high-dimensional coded matrix multiplication.
\newblock \emph{Advances in Neural Information Processing Systems}, 30, 2017.

\bibitem[Reed and Solomon(1960)]{reedSolomon}
Irving~S Reed and Gustave Solomon.
\newblock Polynomial codes over certain finite fields.
\newblock \emph{Journal of the society for industrial and applied mathematics}, 8\penalty0 (2):\penalty0 300--304, 1960.

\bibitem[Guruswami et~al.(2022)Guruswami, Rudra, and Sudan]{SudanBook}
Venkatesan Guruswami, Atri Rudra, and Madhu Sudan.
\newblock \emph{Essential Coding Theory}.
\newblock Draft is Available, 2022.

\bibitem[Lee et~al.(2018)Lee, Lam, Pedarsani, Papailiopoulos, and Ramchandran]{speed}
K.~Lee, M.~Lam, R.~Pedarsani, D.~Papailiopoulos, and K.~Ramchandran.
\newblock Speeding up distributed machine learning using codes.
\newblock \emph{IEEE Transactions on Information Theory}, 64\penalty0 (3):\penalty0 1514--1529, 2018.

\bibitem[Lee et~al.(2017)Lee, Suh, and Ramchandran]{high}
K.~Lee, C.~Suh, and K.~Ramchandran.
\newblock High-dimensional coded matrix multiplication.
\newblock In \emph{Proceedings of IEEE International Symposium on Information Theory (ISIT)}, pages 2418--2422, 2017.

\bibitem[Gautschi and Inglese(1987)]{gautschi1987lower}
Walter Gautschi and Gabriele Inglese.
\newblock Lower bounds for the condition number of vandermonde matrices.
\newblock \emph{Numerische Mathematik}, 52\penalty0 (3):\penalty0 241--250, 1987.

\bibitem[Golub and Van~Loan(2013)]{golub2013matrix}
Gene~H. Golub and Charles~F. Van~Loan.
\newblock \emph{Matrix computations}.
\newblock JHU press, 2013.

\bibitem[So et~al.(2020{\natexlab{b}})So, Guler, and Avestimehr]{so2020scalable}
Jinhyun So, Basak Guler, and Salman Avestimehr.
\newblock A scalable approach for privacy-preserving collaborative machine learning.
\newblock \emph{Advances in Neural Information Processing Systems}, 33:\penalty0 8054--8066, 2020{\natexlab{b}}.

\bibitem[Subramaniam et~al.(2019)Subramaniam, Heidarzadeh, and Narayanan]{AnooshehRobust}
Adarsh~M. Subramaniam, Anoosheh Heidarzadeh, and Krishna~R. Narayanan.
\newblock Random {Khatri-Rao}-product codes for numerically-stable distributed matrix multiplication.
\newblock \emph{CoRR}, abs/1907.05965, 2019.

\bibitem[Ramamoorthy and Tang(2022)]{RamamoorthyCir}
Aditya Ramamoorthy and Li~Tang.
\newblock Numerically stable coded matrix computations via circulant and rotation matrix embeddings.
\newblock \emph{IEEE Transactions on Information Theory}, 68\penalty0 (4):\penalty0 2684--2703, 2022.

\bibitem[Das et~al.(2021)Das, Ramamoorthy, and Vaswani]{RamamoorthyConv}
Anindya~Bijoy Das, Aditya Ramamoorthy, and Namrata Vaswani.
\newblock Efficient and robust distributed matrix computations via convolutional coding.
\newblock \emph{IEEE Transactions on Information Theory}, 67\penalty0 (9):\penalty0 6266--6282, 2021.

\bibitem[Soleymani et~al.(2021)Soleymani, Mahdavifar, and Avestimehr]{soleymani2020analog}
Mahdi Soleymani, Hessam Mahdavifar, and Amir~Salman Avestimehr.
\newblock Analog {Lagrange} coded computing.
\newblock \emph{IEEE Journal on Selected Areas in Information Theory}, 2\penalty0 (1):\penalty0 283--295, 2021.

\bibitem[Fahim and Cadambe(2021{\natexlab{a}})]{fahim2019numerically}
Mohammad Fahim and Viveck~R Cadambe.
\newblock Numerically stable polynomially coded computing.
\newblock \emph{IEEE Transactions on Information Theory}, 67\penalty0 (5):\penalty0 2758--2785, 2021{\natexlab{a}}.

\bibitem[Jahani-Nezhad and Maddah-Ali(2021)]{jahani2018codedsketch}
Tayyebeh Jahani-Nezhad and Mohammad~Ali Maddah-Ali.
\newblock {CodedSketch}: A coding scheme for distributed computation of approximated matrix multiplication.
\newblock \emph{IEEE Transactions on Information Theory}, 67\penalty0 (6):\penalty0 4185--4196, 2021.

\bibitem[Gupta et~al.(2018{\natexlab{b}})Gupta, Wang, Courtade, and Ramchandran]{overSketch}
Vipul Gupta, Shusen Wang, Thomas Courtade, and Kannan Ramchandran.
\newblock {OverSketch}: Approximate matrix multiplication for the cloud.
\newblock In \emph{2018 IEEE International Conference on Big Data (Big Data)}, pages 298--304, 2018{\natexlab{b}}.

\bibitem[Gupta et~al.(2020)Gupta, Kadhe, Courtade, Mahoney, and Ramchandran]{overSketchN}
Vipul Gupta, Swanand Kadhe, Thomas Courtade, Michael~W. Mahoney, and Kannan Ramchandran.
\newblock Oversketched {Newton}: Fast convex optimization for serverless systems.
\newblock In \emph{2020 IEEE International Conference on Big Data (Big Data)}, pages 288--297, 2020.

\bibitem[Jahani-Nezhad and Maddah-Ali(2022)]{jahani2022berrut}
Tayyebeh Jahani-Nezhad and Mohammad~Ali Maddah-Ali.
\newblock Berrut approximated coded computing: Straggler resistance beyond polynomial computing.
\newblock \emph{IEEE Transactions on Pattern Analysis and Machine Intelligence}, 45\penalty0 (1):\penalty0 111--122, 2022.

\bibitem[Kimeldorf and Wahba(1971)]{kimeldorf1971some}
George Kimeldorf and Grace Wahba.
\newblock Some results on tchebycheffian spline functions.
\newblock \emph{Journal of mathematical analysis and applications}, 33\penalty0 (1):\penalty0 82--95, 1971.

\bibitem[Duchon(1977)]{duchon1977splines}
Jean Duchon.
\newblock Splines minimizing rotation-invariant semi-norms in sobolev spaces.
\newblock In \emph{Constructive Theory of Functions of Several Variables: Proceedings of a Conference Held at Oberwolfach April 25--May 1, 1976}, pages 85--100. Springer, 1977.

\bibitem[Wahba(1990)]{wahba1990spline}
Grace Wahba.
\newblock \emph{Spline models for observational data}.
\newblock SIAM, 1990.

\bibitem[Wahba(1975)]{wahba1975smoothing}
Grace Wahba.
\newblock Smoothing noisy data with spline functions.
\newblock \emph{Numerische mathematik}, 24\penalty0 (5):\penalty0 383--393, 1975.

\bibitem[Sch{\"o}lkopf et~al.(2001)Sch{\"o}lkopf, Herbrich, and Smola]{scholkopf2001generalized}
Bernhard Sch{\"o}lkopf, Ralf Herbrich, and Alex~J Smola.
\newblock A generalized representer theorem.
\newblock In \emph{International conference on computational learning theory}, pages 416--426. Springer, 2001.

\bibitem[Paszke et~al.(2019)Paszke, Gross, Massa, Lerer, Bradbury, Chanan, Killeen, Lin, Gimelshein, Antiga, et~al.]{paszke2019pytorch}
Adam Paszke, Sam Gross, Francisco Massa, Adam Lerer, James Bradbury, Gregory Chanan, Trevor Killeen, Zeming Lin, Natalia Gimelshein, Luca Antiga, et~al.
\newblock Pytorch: An imperative style, high-performance deep learning library.
\newblock \emph{Advances in neural information processing systems}, 32, 2019.

\bibitem[LeCun et~al.(1998)LeCun, Bottou, Bengio, and Haffner]{lecun1998gradient}
Yann LeCun, L{\'e}on Bottou, Yoshua Bengio, and Patrick Haffner.
\newblock Gradient-based learning applied to document recognition.
\newblock \emph{Proceedings of the IEEE}, 86\penalty0 (11):\penalty0 2278--2324, 1998.

\bibitem[LeCun et~al.(2010)LeCun, Cortes, Burges, et~al.]{lecun2010mnist}
Yann LeCun, Corinna Cortes, Chris Burges, et~al.
\newblock {MNIST} handwritten digit database, 2010.

\bibitem[Krizhevsky et~al.(2009)Krizhevsky, Hinton, et~al.]{krizhevsky2009learning}
Alex Krizhevsky, Geoffrey Hinton, et~al.
\newblock Learning multiple layers of features from tiny images.
\newblock 2009.

\bibitem[Ding et~al.(2021)Ding, Zhang, Ma, Han, Ding, and Sun]{ding2021repvgg}
Xiaohan Ding, Xiangyu Zhang, Ningning Ma, Jungong Han, Guiguang Ding, and Jian Sun.
\newblock Repvgg: Making vgg-style convnets great again.
\newblock In \emph{Proceedings of the IEEE/CVF conference on computer vision and pattern recognition}, pages 13733--13742, 2021.

\bibitem[Dosovitskiy et~al.(2020)Dosovitskiy, Beyer, Kolesnikov, Weissenborn, Zhai, Unterthiner, Dehghani, Minderer, Heigold, Gelly, et~al.]{dosovitskiy2020image}
Alexey Dosovitskiy, Lucas Beyer, Alexander Kolesnikov, Dirk Weissenborn, Xiaohua Zhai, Thomas Unterthiner, Mostafa Dehghani, Matthias Minderer, Georg Heigold, Sylvain Gelly, et~al.
\newblock An image is worth 16x16 words: Transformers for image recognition at scale.
\newblock \emph{arXiv preprint arXiv:2010.11929}, 2020.

\bibitem[Deng et~al.(2009)Deng, Dong, Socher, Li, Li, and Fei-Fei]{deng2009imagenet}
Jia Deng, Wei Dong, Richard Socher, Li-Jia Li, Kai Li, and Li~Fei-Fei.
\newblock {Imagenet}: A large-scale hierarchical image database.
\newblock In \emph{2009 IEEE conference on computer vision and pattern recognition}, pages 248--255. Ieee, 2009.

\bibitem[Soleymani et~al.(2022)Soleymani, Ali, Mahdavifar, and Avestimehr]{soleymani2022approxifer}
Mahdi Soleymani, Ramy~E Ali, Hessam Mahdavifar, and A~Salman Avestimehr.
\newblock {ApproxIFER}: A model-agnostic approach to resilient and robust prediction serving systems.
\newblock In \emph{Proceedings of the AAAI Conference on Artificial Intelligence}, volume~36, pages 8342--8350, 2022.

\bibitem[Eilers and Marx(1996)]{eilers1996flexible}
Paul~HC Eilers and Brian~D Marx.
\newblock Flexible smoothing with b-splines and penalties.
\newblock \emph{Statistical science}, 11\penalty0 (2):\penalty0 89--121, 1996.

\bibitem[De~Boor(2001)]{de2001calculation}
Carl De~Boor.
\newblock Calculation of the smoothing spline with weighted roughness measure.
\newblock \emph{Mathematical Models and Methods in Applied Sciences}, 11\penalty0 (01):\penalty0 33--41, 2001.

\bibitem[Hu and Schumaker(1986)]{hu1986complete}
Chiu~Li Hu and Larry~L Schumaker.
\newblock Complete spline smoothing.
\newblock \emph{Numerische Mathematik}, 49:\penalty0 1--10, 1986.

\bibitem[Zaharia et~al.(2008)Zaharia, Konwinski, Joseph, Katz, and Stoica]{zaharia2008improving}
Matei Zaharia, Andy Konwinski, Anthony~D Joseph, Randy~H Katz, and Ion Stoica.
\newblock Improving mapreduce performance in heterogeneous environments.
\newblock In \emph{Osdi}, volume~8, page~7, 2008.

\bibitem[Suresh et~al.(2015)Suresh, Canini, Schmid, and Feldmann]{suresh2015c3}
Lalith Suresh, Marco Canini, Stefan Schmid, and Anja Feldmann.
\newblock {C3}: Cutting tail latency in cloud data stores via adaptive replica selection.
\newblock In \emph{12th USENIX Symposium on Networked Systems Design and Implementation (NSDI 15)}, pages 513--527, 2015.

\bibitem[Shah et~al.(2015)Shah, Lee, and Ramchandran]{shah2015redundant}
Nihar~B Shah, Kangwook Lee, and Kannan Ramchandran.
\newblock When do redundant requests reduce latency?
\newblock \emph{IEEE Transactions on Communications}, 64\penalty0 (2):\penalty0 715--722, 2015.

\bibitem[Gardner et~al.(2015)Gardner, Zbarsky, Doroudi, Harchol-Balter, and Hyytia]{gardner2015reducing}
Kristen Gardner, Samuel Zbarsky, Sherwin Doroudi, Mor Harchol-Balter, and Esa Hyytia.
\newblock Reducing latency via redundant requests: Exact analysis.
\newblock \emph{ACM SIGMETRICS Performance Evaluation Review}, 43\penalty0 (1):\penalty0 347--360, 2015.

\bibitem[Chaubey and Saule(2015)]{chaubey2015replicated}
Manmohan Chaubey and Erik Saule.
\newblock Replicated data placement for uncertain scheduling.
\newblock In \emph{2015 IEEE International Parallel and Distributed Processing Symposium Workshop}, pages 464--472. IEEE, 2015.

\bibitem[Li et~al.(2018)Li, Kalan, Avestimehr, and Soltanolkotabi]{li2018near}
Songze Li, Seyed Mohammadreza~Mousavi Kalan, Amir~Salman Avestimehr, and Mahdi Soltanolkotabi.
\newblock Near-optimal straggler mitigation for distributed gradient methods.
\newblock In \emph{2018 IEEE International Parallel and Distributed Processing Symposium Workshops (IPDPSW)}, pages 857--866. IEEE, 2018.

\bibitem[Das et~al.(2019)Das, Ramamoorthy, and Vaswani]{das2019random}
Anindya~B Das, Aditya Ramamoorthy, and Namrata Vaswani.
\newblock Random convolutional coding for robust and straggler resilient distributed matrix computation.
\newblock \emph{arXiv preprint arXiv:1907.08064}, 2019.

\bibitem[Ramamoorthy et~al.(2020)Ramamoorthy, Das, and Tang]{ramamoorthysurvey}
Aditya Ramamoorthy, Anindya~Bijoy Das, and Li~Tang.
\newblock Straggler-resistant distributed matrix computation via coding theory: Removing a bottleneck in large-scale data processing.
\newblock \emph{IEEE Signal Processing Magazine}, 37\penalty0 (3):\penalty0 136--145, 2020.

\bibitem[Karakus et~al.(2017)Karakus, Sun, Diggavi, and Yin]{karakus2017straggler}
Can Karakus, Yifan Sun, Suhas Diggavi, and Wotao Yin.
\newblock Straggler mitigation in distributed optimization through data encoding.
\newblock \emph{Advances in Neural Information Processing Systems}, 30, 2017.

\bibitem[Tandon et~al.(2017)Tandon, Lei, Dimakis, and Karampatziakis]{Grad_Dimakis}
R.~Tandon, Q.~Lei, A.~G. Dimakis, and N.~Karampatziakis.
\newblock Gradient coding: Avoiding stragglers in distributed learning.
\newblock In \emph{Proceedings of the 34th International Conference on Machine Learning (ICML)}, pages 3368--3376, Sydney, Australia, August 2017.

\bibitem[Ramamoorthy and Tang(2021)]{ramamoorthy2021numerically}
Aditya Ramamoorthy and Li~Tang.
\newblock Numerically stable coded matrix computations via circulant and rotation matrix embeddings.
\newblock \emph{IEEE Transactions on Information Theory}, 68\penalty0 (4):\penalty0 2684--2703, 2021.

\bibitem[Yu et~al.(2020{\natexlab{b}})Yu, Maddah-Ali, and Avestimehr]{entangle}
Qian Yu, Mohammad~Ali Maddah-Ali, and Amir~Salman Avestimehr.
\newblock Straggler mitigation in distributed matrix multiplication: Fundamental limits and optimal coding.
\newblock \emph{IEEE Transactions on Information Theory}, 66\penalty0 (3):\penalty0 1920--1933, 2020{\natexlab{b}}.

\bibitem[Fahim and Cadambe(2021{\natexlab{b}})]{fahimappx}
Mohammad Fahim and Viveck~R. Cadambe.
\newblock Numerically stable polynomially coded computing.
\newblock \emph{IEEE Transactions on Information Theory}, 67\penalty0 (5):\penalty0 2758--2785, 2021{\natexlab{b}}.

\bibitem[Kosaian et~al.(2019)Kosaian, Rashmi, and Venkataraman]{kosaian2019parity}
Jack Kosaian, KV~Rashmi, and Shivaram Venkataraman.
\newblock Parity models: A general framework for coding-based resilience in ml inference.
\newblock \emph{arXiv preprint arXiv:1905.00863}, 2019.

\bibitem[Berrut(1988)]{berrut1988rational}
J-P Berrut.
\newblock Rational functions for guaranteed and experimentally well-conditioned global interpolation.
\newblock \emph{Computers \& Mathematics with Applications}, 15\penalty0 (1):\penalty0 1--16, 1988.

\bibitem[Leoni(2024)]{leoni2024first}
Giovanni Leoni.
\newblock \emph{A first course in Sobolev spaces}, volume 181.
\newblock American Mathematical Society, 2024.

\bibitem[Adams and Fournier(2003)]{adams2003sobolev}
Robert~A Adams and John~JF Fournier.
\newblock \emph{Sobolev spaces}.
\newblock Elsevier, 2003.

\bibitem[Berlinet and Thomas-Agnan(2011)]{berlinet2011reproducing}
Alain Berlinet and Christine Thomas-Agnan.
\newblock \emph{Reproducing kernel Hilbert spaces in probability and statistics}.
\newblock Springer Science \& Business Media, 2011.

\bibitem[Ragozin(1983)]{ragozin1983error}
David~L Ragozin.
\newblock Error bounds for derivative estimates based on spline smoothing of exact or noisy data.
\newblock \emph{Journal of approximation theory}, 37\penalty0 (4):\penalty0 335--355, 1983.

\bibitem[Utreras(1988)]{utreras1988convergence}
Florencio~I Utreras.
\newblock Convergence rates for multivariate smoothing spline functions.
\newblock \emph{Journal of approximation theory}, 52\penalty0 (1):\penalty0 1--27, 1988.

\bibitem[Guggenheimer et~al.(1995)Guggenheimer, Edelman, and Johnson]{guggenheimer1995simple}
Heinrich~W Guggenheimer, Alan~S Edelman, and Charles~R Johnson.
\newblock A simple estimate of the condition number of a linear system.
\newblock \emph{The College Mathematics Journal}, 26\penalty0 (1):\penalty0 2--5, 1995.

\end{thebibliography}


%%%%%%%%%%%%%%%%%%%%%%%%%%%%%%%%%%%%%%%%%%%%%%%%%%%%%%%%%%%%%%%%%%%%%%%%%%%%%%%
%%%%%%%%%%%%%%%%%%%%%%%%%%%%%%%%%%%%%%%%%%%%%%%%%%%%%%%%%%%%%%%%%%%%%%%%%%%%%%%
% APPENDIX
%%%%%%%%%%%%%%%%%%%%%%%%%%%%%%%%%%%%%%%%%%%%%%%%%%%%%%%%%%%%%%%%%%%%%%%%%%%%%%%
%%%%%%%%%%%%%%%%%%%%%%%%%%%%%%%%%%%%%%%%%%%%%%%%%%%%%%%%%%%%%%%%%%%%%%%%%%%%%%%
\newpage
\appendix
% \onecolumn
\section{Preliminaries}
\subsection{Sobolev spaces and Sobolev norms}\label{app:sobolev}
Let $\Omega$ be an open interval in $\mathbb{R}$ and $M$ be a positive integer. We denote by $\lpm{p}$ the class of all measurable functions $g:\mathbb{R} \to \mathbb{R}^M$ that satisfy:
\begin{align}
\int_{\Omega}|g_j(t)|^p \,dt\ < \infty, \quad \forall j\in [M],
\end{align}
where $g(\cdot) = [g_1(\cdot),\dots,g_M(\cdot)]^T$. The space $\lpm{p}$ can be endowed with the following norm, known as the ${L}^p$ norm:
\begin{align}
    \|g\|_{\lpm{p}} := \left(\sum^M_{j=1}\int_{\Omega}|g_i(t)|^p \,dt\right)^{\frac{1}{p}},
\end{align}
for $1 \leqslant p < \infty$, and
\begin{align}
    \|g\|_{\lpm{\infty}} := \max_{j \in [M]} \sup_{t\in \Omega} |g_j(t)|.
\end{align}
for $p=\infty$. Additionally, a function $g: \Omega \to \mathbb{R}^M$ is in $\lpmloc{p}$ if it lies in ${L}^p(V;\mathbb{R}^M)$ for all compact subsets $V \subseteq \Omega$. 
\begin{definition}[\textbf{Sobolev Space}] 
The \emph{Sobolev space} $\smp$ is the space of all functions $g \in \lpm{p}$ such that all weak derivatives of order $i$, denoted by $g^{(i)}$, belong to $\lpm{p}$ for $i \in [m]$. This space is endowed with the norm:
\begin{align}\label{eq:sob_norm1}
    \norm{g}_{\smp} :=  \left(\norm{g}_{\lpm{p}}^p  + \sum^m_{i=1} \norm{g^{(i)}}_{\lpm{p}}^p \right)^\frac{1}{p},
\end{align}
for $1 \leqslant p < \infty$, and
\begin{align}\label{eq:sob_norm2}
    \|g\|_{\sobm{m}{\infty}} := \max \left \{ \norm{g}_{\lpm{\infty}},  \max_{i\in [m]} \norm{g^{(i)}}_{\lpm{\infty}}\right \},
\end{align}
for $p=\infty$.
\end{definition}
Similarly, $\smploc$ is defined as the space of all functions $g \in \lpmloc{p}$ with all weak derivatives of order $i$ belonging to $\lpmloc{p}$ for $i \in [m]$. $\norm{g}_{\smploc}$ and $\norm{g}_{\soblocm{m}{\infty}}$ are defined similar to \eqref{eq:sob_norm1} and \eqref{eq:sob_norm2} respectively, using $\lploc{p}$ instead of $\lp{p}$. 

\begin{definition}\label{def:sobz}(\textbf{Sobolev Space with compact support}): Denoted by $\smpz$ collection of functions $g$ defined on interval $\Omega=(a, b)$ such that $g(a)=\mathbf{0}, g'(a) = \mathbf{0}, \dots, g^{(m-1)}(a)=\mathbf{0}$ and $\norm{g^{(m)}}_{\lp{2}} < \infty$. This space can be endowed with the following norm:
\begin{align}
    \norm{g}_{\smp} :=  \norm{g^{(m)}}_{\lpm{p}},
\end{align}
for $1 \leqslant p < \infty$, and
\begin{align}\label{eq:sob_norm2_2}
    \|g\|_{\sobm{m}{\infty}} :=  \norm{g^{(m)}}_{\lpm{\infty}},
\end{align}
for $p=\infty$.
\end{definition}


The next theorem provides an upper bound for ${L}^p$ norm of functions in the Sobolev space, which plays a crucial role in the proof of the main theorems of the paper.
\begin{theorem}[Theorem 7.34, \citep{leoni2024first}] \label{th:interp_ineq} Let $\Omega \subseteq \mathbb{R}$ be an open interval and let $g \in \soblocm{1}{1}$. Assume $1 \leqslant p,q,r \leqslant \infty$ and $r \geqslant q$. Then:
    \begin{align}
        \norm{g}_{\lpm{r}} \leqslant \ell^{\frac{1}{r}-\frac{1}{q}}\norm{g}_{\lpm{q}} + \ell^{1 - \frac{1}{p}+\frac{1}{r}}\norm{g'}_{\lpm{p}},
    \end{align}
    for every $0 < \ell < \mathcal{L}^1(\Omega)$.
\end{theorem}
Note that $\mathcal{L}^1(\Omega)$ in Theorem~\ref{th:interp_ineq} is the length of the interval $\Omega$. 
\begin{corollary}\label{col:interp_ineq}
    Suppose $g \in \soblocm{1}{1}$ and $\Omega \subseteq \mathbb{R}$ be an open interval. If $\frac{\norm{g}_{\lpm{2}}}{\norm{g'}_{\lpm{2}}} < \mathcal{L}^1(\Omega)$, then:
    \begin{align}
        \norm{g}_{\lpm{\infty}} \leqslant 2\sqrt{\norm{g}_{\lpm{2}}\cdot\norm{g'}_{\lpm{2}}}.
    \end{align}
\end{corollary}
\begin{proof}
    Substituting $p,q=2$ and $r=\infty$ in Theorem~\ref{th:interp_ineq} and optimizing over $\ell$, one can derive the optimum value of $\ell$, denoted by $\ell^*$ as
    \begin{align}
    \ell^* = \frac{\norm{g}_{\lpm{2}}}{\norm{g'}_{\lpm{2}}}.
    \end{align}
    Since $\frac{\norm{g}_{\lpm{2}}}{\norm{g'}_{\lpm{2}}} < \mathcal{L}^1(\Omega)$,  the optimum value is in the valid interval mentioned in Theorem~\ref{th:interp_ineq}.
\end{proof}

\begin{corollary}[Corollary 7.36, \citep{leoni2024first}] \label{cor:interp_ineq_2}
     Let $\Omega=(a, b)$, let $1 \leq p, q, r \leq \infty$ be such that $1+1 / r \geq$ $1 / p$ and $r \geq q$ and let $g \in \soblocm{1}{1}$ with $g^{\prime} \in \lpm{p}$. Let $x_0 \in[a, b]$ be such that $\left|g\left(x_0\right)\right|=\min _{[a, b]}|g|$. Then
\begin{align}
\left\|g-g\left(x_0\right)\right\|_{\lpm{r}} \leq 8\|g\|_{\lpm{q}}^\alpha\left\|g^{\prime}\right\|_{\lpm{p}}^{1-\alpha}
\end{align}
where $\alpha:=0$ if $r=q$ and $1-1 / p+1 / r=0$ and otherwise
$$
\alpha:=\frac{1-1 / p+1 / r}{1-1 / p+1 / q}.
$$
\end{corollary}

\begin{corollary}\label{col:interp_ineq2}
Theorem~\ref{th:interp_ineq} and Corollary~\ref{cor:interp_ineq_2} hold true when $f \in \sobm{2}{2}$.
\end{corollary}
\begin{proof}
    We prove the above corollary by showing that $\sobm{2}{2} \subseteq \sobm{1}{1} \subseteq \soblocm{1}{1}$.
    Using Cauchy-Schwartz inequality, one can show:
    \begin{align}\label{eq:col:interp_ineq2_1}
        \norm{f}_{\lpm{1}} = \sum_{j=1}^M \int_{\Omega} |f_j(t)|\,dt &\leqslant \sum_{j=1}^M \left(\int_{\Omega} 1^2\,dt\cdot\int_{\Omega} |f_j(t)|^2\,dt\right)^{\frac{1}{2}} \nonumber \\ &\leqslant \left(\mathcal{L}^1(\Omega)\right)^\frac{1}{2}\cdot \sum_{j=1}^M \left(\int_{\Omega} |f_j(t)|^2\,dt\right)^{\frac{1}{2}} \lec{(a)}{<} \infty,
    \end{align}
    where (a) follows from the bounded length of $\Omega$ and $f \in \sob{2}{2}$. Similarly:
    \begin{align}\label{eq:col:interp_ineq2_2}
        \norm{f'}_{\lpm{1}} = \sum_{j=1}^M \int_{\Omega} |f_j'(t)|\,dt &\leqslant \sum_{j=1}^M \left(\int_{\Omega} 1^2\,dt\cdot\int_{\Omega} |f_j'(t)|^2\,dt\right)^{\frac{1}{2}} \nonumber \\ &\leqslant \left(\mathcal{L}^1(\Omega)\right)^\frac{1}{2}\cdot \sum_{j=1}^M \left(\int_{\Omega} |f_j'(t)|^2\,dt\right)^{\frac{1}{2}} \lec{(a)}{<} \infty.
    \end{align}
    Equations \eqref{eq:col:interp_ineq2_1} and \eqref{eq:col:interp_ineq2_2} prove that $\sobm{2}{2} \subseteq \sobm{1}{1}$. Now, suppose $f\in\sobm{1}{1}$. For every closed subset $[c,d] \subset \Omega$ we have:
    \begin{align}
        \sum_{j=1}^M \left(\int^d_c |f_j(t)|\,dt\right) \lec{(a)}{\leqslant} \sum_{j=1}^M \left(\int_\Omega |f_j(t)|\,dt\right) \lec{(b)}{\leqslant} \infty, \\
        \sum_{j=1}^M \left(\int^d_c |f_j'(t)|\,dt\right) \lec{(a)}{\leqslant} \sum_{j=1}^M \left(\int_\Omega |f_j'(t)|\,dt\right) \lec{(b)}{\leqslant} \infty,
    \end{align}
    where (a) is due to $[c,d] \subset \Omega$ and (b) is because of $f\in\sobm{1}{1}$. Therefore, $f \in \soblocm{1}{1}$.
\end{proof}
{\bf Equivalent norms.} There have been various norms defined on Sobolev spaces in the literature that are equivalent to  \eqref{eq:sob_norm1} (see \citep{adams2003sobolev}, \citep[Ch. 7]{berlinet2011reproducing}, and \citep[Sec. 10.2]{wahba1990spline}). Note that two norms $\norm{\cdot}_{W_1}, \norm{\cdot}_{W_2}$ are equivalent if there exist positive constants $\eta_1, \eta_2$ such that
$
\eta_1\cdot\norm{g}_{W_2} \leqslant \norm{g}_{W_1} \leqslant \eta_2\cdot\norm{g}_{W_2}
$. The equivalent norm in which we are interested is the one introduced in \citep{kimeldorf1971some}. Let $\Omega = (a,b) \subset \mathbb{R}$. We define $\smpeq$ as the Sobolev space endowed with the following norm:
\begin{align}\label{eq:sob_norm_eq}
    \norm{g}_{\smpeq} := \left(\sum^M_{j=1} \left(g_j(a)^p +\sum^{m-1}_{i=1} \left(g_j^{(i)}(a)\right)^p \right) + \norm{g^{(m)}}_{\lpm{p}}^p \right)^\frac{1}{p}.
\end{align}
The following lemma derives the equivalence  constants ($\eta_1,\eta_2$) for the norms $\norm{\cdot}_{\sob{2}{2}}$ and $\norm{\cdot}_{\sobeq{2}{2}}$.
\begin{lemma}\label{lem:norm-equiv}
    Let $\Omega = (a,b)$ be an arbitrary open interval in $\mathbb{R}$. Then for every $g \in \sob{2}{2}$:
    \begin{align}
        \label{lem:norm-equiv1}
        &\norm{g}^2_{\sob{2}{2}}\leqslant \left[2(b-a)\max\{1, (b-a)\} \left (2\max\{1, (b-a)\}^2 + 1 \right) +1\right] \cdot\norm{g}^2_{\sobeq{2}{2}} \\ \label{lem:norm-equiv2}
        &\norm{g}^2_{\sobeq{2}{2}}\leqslant \left(\frac{4}{(b-a)}\max\{1, (b-a)\}^2 + 1\right)\cdot \norm{g}^2_{\sob{2}{2}}.
    \end{align}
\end{lemma}
\begin{proof}
    By expanding $g$ around the point $a$ and utilizing the integral remainder form of Taylor's expansion, for every $x \in (a,b)$ we have:
    \begin{align}
        g(x) = g(a) + \int_a^x g'(t)\,dt.
    \end{align}
    Therefore: \begin{align}\label{eq:semi_norm_neq1}
        g(x)^2 &= g(a)^2 + \left(\int_a^x g'(t)\,dt \right)^2 + 2g(a)\cdot\int_a^x g'(t)\,dt \nonumber \\ 
        &\lec{(a)}{\leqslant} 2g(a)^2 + 2\left(\int_a^x g'(t)\,dt \right)^2 \nonumber \\ 
        &\lec{(b)}{\leqslant} 2g(a)^2 + 2(x-a)\int_a^x g'(t)^2\,dt \nonumber \\  
        &\lec{(c)}{\leqslant}  2g(a)^2 + 2(b-a)\int_a^b g'(t)^2\,dt,
    \end{align}
    where (a) follows from the $ 2g(a)\cdot\int_a^x g'(t)\,dt \leqslant g(a)^2 + \left(\int_a^x g'(t)\,dt \right)^2$, (b) is based on Cauchy-Schwartz inequality, and (c) is due to $x \leqslant b$. Following the same steps as \eqref{eq:semi_norm_neq1}, we have:
\begin{align}\label{eq:semi_norm_neq2}
        g'(x)^2 \leqslant 2g'(a)^2 + 2(b-a)\int_a^b g''(t)^2\,dt.
    \end{align}
    Integrating both sides of \eqref{eq:semi_norm_neq2} we have:
\begin{align}\label{eq:semi_norm_neq3}
        \int_a^b g'(y)^2\,dy &\leqslant 2\int_a^b g'(a)^2\,dy + 2(b-a)\int_a^b \left(\int_a^b g''(t)^2\,dt\right) \,dy \nonumber \\ 
        &= 2(b-a).g'(a)^2 + 2(b-a)^2 \int_a^b g''(t)^2\,dt \nonumber \\
        &\lec{(a)}{\leqslant} 2(b-a)\cdot \max\{1, (b-a)\} \cdot \left(g(a)^2 + g'(a)^2 + \int_a^b g''(t)^2\,dt \right)\nonumber \\
        &= 2(b-a)\cdot\max\{1, (b-a)\}\cdot\norm{g}^2_{\sobeq{2}{2}},
    \end{align}
    where (a) is based on adding the positive term $(b-a)\cdot g(a)^2$ to the right-hand side. Based on the \eqref{eq:semi_norm_neq1} and \eqref{eq:semi_norm_neq3} we have:
\begin{align}\label{eq:semi_norm_neq4}
        \int_a^b g(x)^2\,dx &\lec{(a)}{\leqslant} 2\int_a^b g(a)^2\,dx + 2(b-a)\int_a^b \left(\int_a^b g'(t)^2\,dt\right)\,dx \nonumber \\ 
        &= 2(b-a).g(a)^2 + 2(b-a)^2 \int_a^b g'(t)^2\,dt \nonumber \\ 
        &\lec{(b)}{\leqslant} 2(b-a)\cdot g(a)^2 + 4(b-a)^3\cdot g'(a)^2 + 4(b-a)^4 \int_a^b g''(t)^2\,dt \nonumber \\ 
        &\leqslant 4(b-a)\cdot\max\{1, (b-a)^3\}\cdot\norm{g}^2_{\sobeq{2}{2}},
    \end{align}
    where (a) follows by integrating both sides of \eqref{eq:semi_norm_neq1} and (b) is due to \eqref{eq:semi_norm_neq2}. Using \eqref{eq:semi_norm_neq4} and \eqref{eq:semi_norm_neq3} and the fact that $\int_a^b g''(t)^2\, dt \leqslant \|g\|^2_{\sob{2}{2}}$, we have:
    \begin{align}
        \norm{g}^2_{\sob{2}{2}} &= \int_a^b g''(t)^2\, dt +  \int_a^b g'(t)^2\, dt + \int_a^b g(t)^2\, dt \nonumber \\  &\leqslant \left[2(b-a)\max\{1, (b-a)\} \left (2\max\{1, (b-a)\}^2 + 1 \right) +1\right] \cdot\norm{g}^2_{\sobeq{2}{2}}.
    \end{align}
    For the other side, using the same steps as in \eqref{eq:semi_norm_neq1}, we have:
    \begin{align}\label{eq:semi_norm_neq5}
        g(a)^2 &= g(x)^2 + \biggl(\int_a^x g'(t)\,dt \biggr)^2 - 2g(a).\int_a^x g'(t)\,dt \nonumber \\ 
        &\lec{(a)}{\leqslant} 2g(x)^2 + 2\biggl(\int_a^x g'(t)\,dt \biggr)^2 \nonumber \\
        &\lec{(b)}{\leqslant} 2g(x)^2 + 2(x-a)\int_a^x g'(t)^2\,dt \nonumber \\  
        &\lec{(c)}{\leqslant}  2g(x)^2 + 2(b-a)\int_a^b g'(t)^2\,dt,
    \end{align}
    where (a) is because of $-2g(a)\cdot\int_a^x g'(t)\,dt \leqslant g(a)^2 + \left(\int_a^x g'(t)\,dt \right)^2$, (b) follows from Cauchy-Schwartz inequality, and (c) is due to $x \leqslant b$. Integrating both sides of \eqref{eq:semi_norm_neq5}, we have
    \begin{align}\label{eq:semi_norm_neq6}
         (b-a)\cdot g(a)^2 = \int_a^b g(a)^2\,dx  &\leqslant 2\int_a^b g(x)^2\,dx + 2(b-a)\int_a^b \left(\int_a^b g'(t)^2\,dt \right)\,dx \nonumber \\ 
        &\lec{(a)}{=} 2\int_a^b g(t)^2\,dt + 2(b-a)^2 \int_a^b g'(t)^2\,dt \nonumber \\
        &\lec{(b)}{\leqslant} 2\max\{1, (b-a)\}^2\left(\int_a^b g(t)^2\,dt + \int_a^b g'(t)^2\,dt + \int_a^b g''(t)^2\,dt\right) \nonumber \\
        &\leqslant 2\max\{1, (b-a)\}^2 \norm{g}^2_{\sob{2}{2}},
    \end{align}
    where (a) follows by a change of variable from $x$ to $t$ and (b) follows by adding the positive term $2\max\{1, (b-a)\}^2 \cdot \int_a^b g''(t)^2\,dt$ to the right side. Thus, \eqref{eq:semi_norm_neq6} directly results in
    \begin{align}\label{eq:semi_norm_neq7}
        g(a)^2 \leqslant \frac{2}{(b-a)}\max\{1, (b-a)\}^2 \norm{g}^2_{\sob{2}{2}}.
    \end{align}
    Following same steps as \eqref{eq:semi_norm_neq6} and \eqref{eq:semi_norm_neq7} results in
    \begin{align}
        g'(a)^2 \leqslant \frac{2}{(b-a)}\max\{1, (b-a)\}^2 \norm{g}^2_{\sob{2}{2}}.
    \end{align}
    Considering the fact that $\int_a^b g''(t)^2\,dt \leqslant \norm{g}^2_{\sob{2}{2}}$, we can complete the proof:
    \begin{align}
        \norm{g}^2_{\sobeq{2}{2}} &= g(a)^2 + g'(a)^2 + \int_a^b g''(t)^2\,dt  \nonumber \\ &\leqslant \left(\frac{4}{(b-a)}\max\{1, (b-a)\}^2 + 1\right)\cdot\norm{g}^2_{\sob{2}{2}}.
    \end{align}
\end{proof}
\begin{corollary}\label{cor:norm-equiv-1}
Based on Lemma~\ref{lem:norm-equiv}, for $\Omega=(-1, 1)$ we have 
\begin{align}
     \frac{1}{9}\cdot\norm{g}^2_{\sobeq{2}{2}} \leqslant \norm{g}^2_{\sob{2}{2}} \leqslant 73\cdot\norm{g}^2_{\sobeq{2}{2}}.
\end{align}   
\end{corollary}
Corollary~\ref{cor:norm-equiv-1} directly follows from Lemma~\ref{lem:norm-equiv} by substituting $(a,b)=(-1,1)$.
\begin{corollary}\label{cor:norm-equiv-2}
The result of Lemma~\ref{lem:norm-equiv} remains valid for multi-dimensional cases, where $g:\mathbb{R} \to \mathbb{R}^M$, for some $M>1$.
\end{corollary}
Corollary~\ref{cor:norm-equiv-2} directly follows from applying Lemma~\ref{lem:norm-equiv} to each component of the function $g(\cdot)$ and using the definition of vector-valued function norm:
$$
\norm{g}^2_{\sobm{2}{2}} = \sum^M_{j=1}\norm{g^{(j)}}^2_{\sob{2}{2}}, \quad \norm{g}^2_{\sobmeq{2}{2}} = \sum^M_{j=1}\norm{g^{(j)}}^2_{\sobeq{2}{2}}.
$$
\begin{proposition}\label{prop:sob_hilb} 
(\citep[Section 7.2]{leoni2024first}, \citep[Theorem 121]{berlinet2011reproducing},\citep{wahba1990spline}) For any open interval $\Omega \subseteq \mathbb{R}$ and $m, M \in \mathbb{N}$, 
\begin{align}
    \hilm{m} &:= \sobm{m}{2}, \nonumber \\ \hilmtilde{m} &:= \sobmeq{m}{2},\nonumber \\\hilmz{m} &:= \sobmz{m}{2},
\end{align}
are Reproducing Kernel Hilbert Spaces (RKHSs).
\end{proposition}
The full expression of the kernel function of $\hil{m}$ and $\hiltilde{m}$ and other equivalent norms of Sobolev spaces can be found in \citep[Section 4]{berlinet2011reproducing}. For $\hiltilde{m}$ the kernel function is as follows:
\begin{align}\label{eq:sob:kernel}
    R(t, s) = \sum_{j=0}^{m-1} \frac{t^j s^j}{j!^2}+\int_\Omega \frac{(t-x)_{+}^{m-1}(s-x)_{+}^{m-1}}{(m-1)!^2}\, dx,
\end{align}
where $(\cdot)_{+}$ is positive part function.
\subsection{Smoothing Splines}\label{app:smoothspline}
Consider the data model $y_i = f(t_i) + \epsilon_i$ for $i=1,\dots,n$, where $t_i \in \Omega=(a,b) \subset \mathbb{R}$, $\mathbb{E}[\epsilon_i] = 0$, and $\mathbb{E}[\epsilon_i^2] \leqslant \sigma_0^2$. 
Assuming $f\in \sobeq{m}{2}$, the solution to the following optimization problem is referred to as the smoothing spline:
\begin{align}\label{eq:sm_spline_obj}
    \spline[\mathbf{y}] :=\underset{g \in \sobeq{m}{2}}{\operatorname{argmin}} \frac{1}{n} \sum^n_{i=1}\left(g\left(t_i\right)-y_i\right)^2+\lambda \int_{\Omega} \left(g^{(m)}(t)\right)^2\,dt,
\end{align}
where $\mathbf{y} = [y_1,\dots, y_n]$. Based on Proposition~\ref{prop:sob_hilb}, $\hiltilde{m}:=\sobeq{m}{2}$ with the norm $\norm{\cdot}_{\sobeq{m}{2}}$ is a RKHS for some kernel function $\phi(\cdot,\cdot)$. Therefore for any $v \in \sobeq{m}{2}$, we have:
\begin{align}\label{eq:smothspline_kernel}
    v(t) = \langle v(\cdot), \phi(\cdot, t)\rangle_{\hiltilde{m}}.
\end{align}
It can be shown that $\phi(t,s) = R^P(t, s) + R^0(t,s)$ where $R^0(t, s)$ is kernel function of $\hilz{m}$ and $R^P(t,s)$ is a null space of $\hilz{m}$ which is the space of all polynomials with degree less than $m$.

The solution of \eqref{eq:sm_spline_obj} has the following form \citep{wahba1990spline,duchon1977splines}:
\begin{align}
    u^*(\cdot) = \sum^m_{i=1} d_i \zeta_i(\cdot) + \sum^n_{j=1} c_j \nu_j(\cdot),
\end{align}
where $\nu_j(\cdot) = R^0(\cdot, t_j)$ for $j \in [n]$, and $\left\{\zeta_i(\cdot)\right\}_{i=1}^m$ are the basis functions of the space of polynomials of degree at most $m-1$. Substituting $u^*$ into \eqref{eq:sm_spline_obj} and optimizing over $\mathbf{c} = [c_1,\dots, c_n]^T$ and $\mathbf{d} = [d_1,\dots, d_m]^T$, we obtain the following result \citep{wahba1990spline}:
\begin{align}\label{eq:spline_linear}
    \spline[\mathbf{y}](\mathbf{y}) = \mathbf{Q}\left(\mathbf{Q}^{\mathbf{T}} \mathbf{Q}+\lambda \mathbf{\Gamma}\right)^{-1} \mathbf{Q}^{\mathbf{T}} \mathbf{y},
\end{align}
where
\begin{align}
\mathbf{Q}_{n \times(n+m)}&=\left[\begin{array}{ll}
\mathbf{T}_{n \times m} & \boldsymbol{\Sigma}_{n \times n}
\end{array}\right], \nonumber \\
\mathbf{\Gamma}_{(n+m) \times(n+m)}&=\left[\begin{array}{cc}
\mathbf{0}_{m \times m} & \mathbf{0}_{m \times n} \nonumber  \\
\mathbf{0}_{n \times 1} & \boldsymbol{\Sigma}_{n \times n}
\end{array}\right], \nonumber \\
\boldsymbol{T}_{ij} &= \zeta_j(t_i),\nonumber \\
\boldsymbol{\Sigma}_{ij} &= R^0(t_i, t_j).
\end{align}
Equation \eqref{eq:spline_linear} states that the smoothing spline fitted on the data points $\mathbf{y}$ is a linear operator:
\begin{align}\label{eq:ssfunction_def}
\spline[\mathbf{y}](\mathbf{z}) := \mathbf{A_{\lambda}} \mathbf{z},
\end{align}
 for $\mathbf{z} \in \mathbb{R}^n$, where $\mathbf{A_\lambda}:=\mathbf{Q}\left(\mathbf{Q}^{\mathbf{T}} \mathbf{Q}+\lambda \mathbf{\Gamma}\right)^{-1} \mathbf{Q}^{\mathbf{T}}$. 
It can be shown that the $u^*(\cdot)$ is a natural spline \citep{wahba1975smoothing, wahba1990spline}. Thus, if $\{b_i(\cdot)\}_{i=1}^n$ is a basis function for an $m$-th order natural spline (such as truncated power or B-spline basis functions), we have:
\begin{align}\label{eq:spline_bspline_1}
     u^*(t) &= \sum^n_{i=1} \xi_ib_i(t),\\
     \boldsymbol{\xi} &= \left(\mathbf{N}^T\mathbf{N} + \lambda\Phi\right)^{-1}\mathbf{N}^T\mathbf{y},
\end{align}
where $\mathbf{N}_{ij} = b_i(t_j), \Phi_{ij} = \int_\Omega b''_i(t)b''_j(t)\,dt$ for $i, j \in [n]$, and $\boldsymbol{\xi} := [\xi_1,\dots,\xi_n]^T$.
% where $\mathbf{\Gamma}=(\mathbf{N}^T)^{-1}\Phi\mathbf{N}^{(-1)}$ and $\mathbf{N}, \mathbf{\Phi} \in \mathbb{R}^{n\times n}$ are defined as:
% \begin{align}
%     \mathbf{N}_{ij} = \phi(t_i, t_j) \quad\, \quad \Phi_{ij} = \langle\phi(\cdot, t_i), \phi(\cdot, t_j)\rangle_{\hiltilde{m}}.
% \end{align}

To characterize the estimation error of the smoothing spline, $|f - \spline(\mathbf{y})|$, we need to define two variables analogous to those in \eqref{eq:max_dist_st}, which quantify the minimum and maximum consecutive distance of the regression points $\{t_i\}^n_{i=1}$:
\begin{gather}\label{eq:max_dist}
\Delta_\textrm{max}:=\underset{i\in \{0\} \cup [n]}{\max} \left\{t_{i+1}-t_i\right\}, \quad \Delta_\textrm{min}:=\underset{i\in [n-1]}{\min} \left\{t_{i+1}-t_i\right\},
\end{gather}
where boundary points are defined as $(t_0, t_{n+1}):=(a,b)$. The following theorem offers an upper bound for the $j$-th derivative of the smoothing spline estimator error function in the absence of noise ($\sigma_0 = 0$).
\begin{theorem} \label{th:lit_spline_noiseless} (\citep[Theorem 4.10]{ragozin1983error})
Consider data model $y_i = f(t_i)$ with $\{t_i\}^n_{i=1}$ belong to $\Omega=[a, b]$ for $i \in [n]$. Let
\begin{align}
L=p_{2(m-1)}(\frac{\Delta_\textrm{max}}{\Delta_\textrm{min}})\cdot\frac{n \Delta_\textrm{max}}{b-a} \frac{\lambda}{2}+D(m)\cdot \left({\Delta_\textrm{max}}\right)^{2 m},
\end{align}
where $p_{d}(\cdot)$ is a degree $d$ polynomial with positive weights and $D(m)$ is a function of $m$. Then for each  $j\in \{0,1,\dots,m\}$ and any $f\in\sob{m}{2}$, there exist a function $H(m, j)$ such that:
\begin{align}\label{eq:noiseless_bound}
    \norm{\left(f - \spline(\mathbf{y})\right)^{(j)}}_{\lp{2}}^2 \leqslant H(m, j)\left(1+\left(\frac{L}{(b-a)^{2m}}\right)\right)^{\frac{j}{m}}\cdot L^{\frac{(m-j)}{m}}\cdot \norm{f^{(m)}}^2_{\lp{2}}.
\end{align}
\end{theorem}
Note that $\left(f - \spline(\mathbf{y})\right)^{(0)} := f - \spline(\mathbf{y})$. In the presence of noise, where $\sigma_0 > 0$, we can exploit the linearity of the smoothing spline operator and the mutual independence of the noise terms to conclude that:
\begin{align}\label{eq:smooth_spline_linear_eq}
   \mathbb{E}_{\bm{\epsilon}}\left[\norm{\left(f - \spline(\mathbf{y})\right)}_{\lp{2}}^2\right] &= \mathbb{E}_{\bm{\epsilon}}\left[\norm{\left(f - \spline[\mathbf{y}](\mathbf{f} + \bm{\epsilon})\right)}_{\lp{2}}^2\right] \nonumber \\ &=\norm{\left(f - \spline[\mathbf{y}](\mathbf{f})\right)}_{\lp{2}}^2 \nonumber \\ &+ \mathbb{E}_{\bm{\epsilon}}\left[\norm{\left(f - \spline[\mathbf{y}](\bm{\epsilon})\right)}_{\lp{2}}^2\right],
\end{align}
where $\mathbf{f} = [f(t_1),\dots,f(t_n)]^T$. The first term in \eqref{eq:smooth_spline_linear_eq} can be upper bounded using Theorem~\ref{th:lit_spline_noiseless}. The following theorem establishes an upper bound for the second term when $\frac{\Delta_\textrm{max}}{\Delta_\textrm{min}}$ is bounded:
\begin{theorem} \label{th:lit_spline_noisy} (\citep[Section 5]{utreras1988convergence})
Consider data model $y_i = f(t_i) + \epsilon_i$, where $\mathbb{E}[\epsilon_i] = 0$ and $\mathbb{E}[\epsilon_i^2]\leqslant \sigma_0^2$ for $i\in[n]$. Assume there exist a constant $B>0$ such that $\frac{\Delta_\textrm{max}}{\Delta_\textrm{min}} \leqslant B$. Then for each $j\in \{0,1,\dots,m\}$, there exist a constant $\lambda_0 >0$ and function $Q(m, j,\lambda_0)$ such that:
\begin{align}
\mathbb{E}_{\bm{\epsilon}}\left[\norm{\left(f - \spline[\mathbf{y}](\bm{\epsilon})\right)^{(j)}}_{\lp{2}}^2\right] \leqslant \frac{Q(m,j,\lambda_0)\cdot\sigma_0^2}{n}\lambda^{\frac{-(2j+1)}{2m}}
\end{align}
for $\lambda \leqslant \lambda_0$ and $n \lambda^{\frac{1}{2m}} \geqslant 1$.
\end{theorem}
Note that based on \citep{utreras1988convergence}, $Q(m,j,\lambda_0) = w(m, j)\lambda_0^\frac{1}{2m} + \Tilde{w}(m, j)$. 




\section{Proof of Theorems}\label{app:theorem_proofs}
Recall from \eqref{eq:decompose} that $\mathcal{R}(\hat{f}) \leqslant \lenc(\hat{f}) + \ldec(\hat{f})$, where 
\begin{gather}
    \ldec(\hat{f}) = \mathbb{E}_{\epsilon, \stset \sim F_{S,N}} \left[\frac{2}{K} \sum^K_{k=1} \left(\decs(\alpha_k) - f(\encs(\alpha_k))\right)^2 \right], \label{eq:ldec}\\ \lenc(\hat{f}) = \frac{2}{K} \sum^K_{k=1} \left(f(\encs(\alpha_k)) - f(x_k)\right)^2.
\end{gather}
We begin by deriving a general intermediate bound for $\ldec$ and $\lenc$ which will be a key component in the proofs of Theorems~\ref{th:conv_rate_noisy} and \ref{th:conv_rate_noiseless}. The subsequent subsections will then provide the remaining details to complete the proofs of both theorems.

\begin{lemma}\label{lem:lt}
Let $f:\mathbb{R} \to \mathbb{R}$ be a $q$-Lipschitz continuous function. Then:
\begin{align}
    \lenc \leqslant\frac{2q^2}{K} \sum^K_{k=1} (\encs(\alpha_k) - x_k)^2.
\end{align}
\end{lemma}
\begin{proof}
Using Lipschitz property, we have:
\begin{align}
     \lenc &= \frac{2}{K} \sum^K_{k=1} \left(f(\encs(\alpha_k))- f(x_k)\right)^2 \nonumber \\ 
     &= \frac{2}{K} \sum^K_{k=1} \left(|f(\encs(\alpha_k))- f(x_k)|\right)^2 \nonumber \\ 
     &\leqslant \frac{2}{K} \sum^K_{k=1} \left(q\cdot|\encs(\alpha_k)- x_k|\right)^2 \nonumber \\
     &= \frac{2q^2}{K} \sum^K_{k=1}  (\encs(\alpha_k)- x_k)^2.
\end{align}
\end{proof}
As previously mentioned,  $\ldec$ represents the expected generalization error of the decoder function. To leverage the results from Theorems~\ref{th:lit_spline_noiseless} and \ref{th:lit_spline_noisy} we must establish that the composition of $f$ with the encoder $\encs$ belongs to the Sobolev space $\sob{2}{2}$.
\begin{lemma}\label{lem:fou}
Let $f:\mathbb{R} \to \mathbb{R}$ be a $q$-Lipschitz continuous function with $\norm{f''}_{\lp{\infty}} \leqslant \nu$ and $\Omega \subset \mathbb{R}$ be an open interval. If $\encs \in \sob{2}{2}$ then $f\circ \encs \in \sob{2}{2}$.
\end{lemma}
\begin{proof}
Let us define $f_0(t):=f(t)-f(0)$. Thus $f_0(0) = 0$ and $f_0$ is $q$-Lipschitz. Using Lipschitz property
\begin{align}\label{eq:fou_1}
|f_0(\encs(t))|^2 = |f_0(\encs(t)) - f_0(0)|^2 \leqslant q^2\cdot\encs(t)^2.
\end{align}
Integrating both sides of \eqref{eq:fou_1}:
\begin{align}\label{eq:fou_2}
    \int_\Omega (f_0\circ \encs(t))^2\,dt \leqslant q^2 \cdot \int_\Omega \encs(t)^2\,dt \lec{(a)}{<} \infty,
\end{align}
where (a) follows by $\encs \in \sob{2}{2}$. Given that $f_0$ is $q$-Lipschitz, its derivative is bounded in the $\lp{\infty}$-norm, i.e., $\norm{f'_0}_{\lp{\infty}} \leqslant q$. Thus
\begin{align}\label{eq:fou_3}
    \int_\Omega \left((f_0\circ \encs(t)\right)')^2\,dt \lec{(a)}{=} \int_\Omega \left(f'_0( \encs(t))\right)^2\cdot \encs'(t)^2\,dt
    \leqslant q^2 \cdot \int_\Omega \encs'(t)^2\,dt \lec{(b)}{<} \infty,
\end{align}
where (a) follows by chain rule and (b) follows by $\encs \in \sob{2}{2}$. For the second derivative we have:
\begin{align}\label{eq:fou_4}
\int_\Omega\left[f_0(\encs(t))''\right]^2 \,dt &\lec{(a)}{=} \int_\Omega \left[ \encs''(t)\cdot f_0'(\encs(t)) + \encs'(t)^2 \cdot f_0''(\encs(t))\right]^2\,dt \nonumber \\ 
    &\lec{(b)}{\leqslant} \int_\Omega \left[\encs''(t)^2 + \encs'(t)^4\right] \left[f_0'(\encs(t))^2 + f_0''(\encs(t))^2\right]\,dt \nonumber \\
    &\lec{(c)}{\leqslant} (q^2 +\nu^2) \int_\Omega \left[\encs''(t)^2 + \encs'(t)^4\right]\,dt \nonumber\\ 
    &\lec{}{=} (q^2 + \nu^2) \left(\norm{\encs''(t)}^2_{\lp{2}} + \norm{\encs'(t)}^4_{\lp{4}} \right) \nonumber\\
    &\lec{(d)}{\leqslant} (q^2 + \nu^2) \left(\norm{\encs''(t)}^2_{\lp{2}} + 
    \left(\norm{\encs'(t)}_{\lp{2}} + \norm{\encs''(t)}_{\lp{2}}\right)^4\right) \nonumber \\ &\lec{(e)}{<} \infty,
    % &\lec{(e)}{\leqslant} (q^2 + \nu^2) \left(\norm{\encs''(t)}^2_{\lp{2}} + 
    % 4\left(\norm{\encs'(t)}^2_{\lp{2}} + \norm{\encs''(t)}^2_{\lp{2}}\right)^2\right)\lec{(e)}{<} \infty,
    % &\lec{(d)}{\leqslant} (q^2 +\nu^2) \left[\int_\Omega \encs''(t)^2\,dt +  \left(\int_\Omega \encs'(t)^2 \right)^2\,dt\right] \lec{(e)}{<} \infty,
\end{align}
where (a) follows from the chain rule, (b) is derived using the Cauchy-Schwartz inequality, (c) is due to $\norm{f_0''}_{\lp{\infty}} \leqslant \nu$, (d) follows from Theorem~\ref{th:interp_ineq} with $r=4, p=q=2, l=1$, and (e) is a result of $\encs \in \sob{2}{2}$. Equations \eqref{eq:fou_2}, \eqref{eq:fou_3}, and \eqref{eq:fou_2} demonstrate that $f_0 \circ \encs \in \sob{2}{2}$. Note that $\Omega$ is bounded, and every constant function belongs to $\sob{2}{2}$. Thus, we can conclude that $f_0 \circ \encs(t) + f(0) = f \circ \encs(t) \in \sob{2}{2}$.
\end{proof}
Let us define the function $h(t) := \decs(t) - f(\encs(t))$. Based on Lemma~\ref{lem:fou}, and given that $\decs$ and $f \circ \encs$ belong to the Sobolev space $\sob{2}{2}$, it follows that $h \in \sob{2}{2}$. The subsequent lemmas establish upper bounds for $\norm{h}_{\lp{\infty}}$ and $\norm{h'}_{\lp{\infty}}$, leveraging properties of functions in Sobolev spaces.
\begin{lemma}\label{lem:norm_infty_bound_1}
    If $\Omega=(-a, a)$, then:
\begin{align}
    \|h\|_{\lp{2}} \leqslant\left\|h^{\prime}\right\|_{\lp{2}} \cdot \sqrt{x_{0}^{2}+a^{2}} \leqslant\left\|h^{\prime}\right\|_{\lp{2}} \cdot \sqrt{2} a
\end{align}
\end{lemma}
\begin{proof}
Assume $\exists\, x_{0} \in \Omega: h\left(x_{0}\right)=0$. Therefore, $|h(x)|=\left|\int_{x_{0}}^{x} h^{\prime}(x)\,dx\right|$ for $x \in\left[x_{0}, a\right)$. Thus
\begin{align}
    |h(x)| =\left|\int_{x_{0}}^{x} h^{\prime}(x)\,dx\right| \lec{(a)}{\leqslant} \int_{x_{0}}^{x}\left|h^{\prime}(x)\right|\,dx \lec{(b)}{\leqslant} \left(\int_{x_{0}}^{x} 1^{2}\,dx\right)^{\frac{1}{2}}\left(\int_{x_{0}}^{x} |h^{\prime}(x)|^2\,dx\right)^{\frac{1}{2}},
\end{align}
where (a) and (b) are followed by the triangle and Cauchy-Schwartz inequalities respectively. Integrating the square of both sides over the interval $[x_0, a)$ yields:
\begin{align}\label{eq:norm_infty_bound_1}
 \int_{x_{0}}^{a}|h(x)|^{2}\,dx \lec{}{\leqslant} \int_{x_{0}}^{a}\left(z-x_{0}\right) \cdot \left( \int_{x_{0}}^{a}\left|h^{\prime}(x)\right|^{2} \,dx \right) \,dz \lec{(a)}{\leqslant} \int_{x_{0}}^{a}\left(z-x_{0}\right)\,dz \cdot \int_{-a}^{a}\left|h^{\prime}(x)\right|^{2}\,dx,
\end{align}
where (a) follows by $x_0 < a$. On the other side, we have the following for every $x \in\left(-a, x_{0}\right]$:
\begin{align}
& |h(x)|=\left|\int_{x}^{x_{0}} h^{\prime}(x)\,dx\right| \leqslant \int_{x}^{x_{0}}\left|h^{\prime}(x)\right|\,dx \leqslant\left(\int_{x}^{x_{0}} 1^{2}\,dx\right)^{\frac{1}{2}} \cdot\left(\int_{x}^{x_{0}}\left|h^{\prime}(x)\right|^{2}\,dx\right)^{\frac{1}{2}}.
\end{align}
Therefore, we have a similar inequality:
\begin{align}\label{eq:norm_infty_bound_2}
& \qquad \int_{-a}^{x_{0}}|h(x)|^{2}\,dx \leqslant \int_{-a}^{x_{0}}\left(x_{0}-x\right)\,dx \cdot \int_{-a}^{a}|h'(x)|^{2}\,dx.
\end{align}
Using \eqref{eq:norm_infty_bound_1} and \eqref{eq:norm_infty_bound_2} completes the proof:
\begin{align}
\|h\|_{L^{2}(\Omega)}^{2}&=\int_{-a}^{x_{0}}|h(x)|^{2}\,dx+\int_{x_{0}}^{a}|h(x)|^{2} \,dx \nonumber \\ 
&\leqslant\left\|h^{\prime}\right\|_{\lp{2}}^{2} \cdot \left(\int_{-a}^{x_{0}}\left(x_{0}-x\right)\,dx+\int_{x_{0}}^{a}\left(x-x_{0}\right)\,dx\right) \nonumber \\
&\leqslant\left\|h^{\prime}\right\|_{\lp{2}}^{2} \cdot\left(x_{0}\left(x_{0}+a\right)-\left(\frac{x_{0}^{2}-a^2}{2}\right)+\left(\frac{a^{2}-x_{0}^{2}}{2}\right)-x_{0}\left(a-x_{0}\right)\right) \\
&= \left\|h^{\prime}\right\|_{\lp{2}}^{2} \cdot\left(x_{0}^{2}+a^{2}\right) \leqslant\left\|h^{\prime}\right\|^2_{\lp{2}} 2a^2.
\end{align}
Thus, if $x_0$ exists, the proof is complete. In the next step, we prove the existence of  $x_0\in\Omega$ such that $h\left(x_{0}\right)=0$. Recall that $h(t) = \decs(t) - f(\encs(t))$ and $\decs(\cdot)$ is the solution of \eqref{eq:decoder_opt}. Assume there is no such $x_0$. Since $h \in \sob{2}{2}$, then $h(\cdot)$ is continuous. Therefore, if there exist $t_1,t_2 \in \Omega$ such that $h(t_1) < 0$ and $h(t_2) > 0$, then the intermediate value theorem states that there exists $x_0 \in (t_1, t_2)$ such that $h(x_0)=0$. Thus, $h(t) >0$ or $h(t)<0$ for all $t\in\Omega$. Without loss of generality, assume the first case where $h(t) > 0$ for all $t\in\Omega$. It means that $\decs(t) > f(\encs(t))$ for all $t\in\Omega$. Let us define $$\beta^* := \underset{\beta \in \stset}{\operatorname{argmin}}\, \decs(\beta) - f(\encs(\beta)).$$ 
Let $\decsbar(t) := \decs(t) - h(\beta^*)$. Note that $\int_\Omega \left(\decsbar''(t)\right)^2\,dt = \int_\Omega \left(\decs''(t)\right)^2\,dt$. Therefore, 
\begin{align}
    \sum_{v \in \stset}\left[\decs\left(\beta_v\right)-f\left(\encs\left(\beta_v\right)\right)\right]^2 &= \sum_{v \in \stset}\left[\decsbar \left(\beta_v\right) + h(\beta^*)-f\left(\encs\left(\beta_v\right)\right)\right]^2 \nonumber \\
    &= \sum_{v \in \stset}\left[\decsbar\left(\beta_v\right)-f\left(\encs\left(\beta_v\right)\right)\right]^2 + |\stset|\cdot h(\beta^*)^2 \nonumber \\
    &\quad\quad\quad + 2h(\beta^*) \sum_{v\in \stset}[\decsbar(\beta_v) - f(\encs(\beta_v))] \nonumber \\
    &\lec{(a)}{\geqslant} \sum_{v \in \stset} \left[\decsbar\left(\beta_v\right)-f\left(\encs\left(\beta_v\right)\right)\right]^2,
\end{align}
where (a) follows from $h(\beta^*) > 0$ and $\decsbar(\beta_v) > f(\encs(\beta_v))$ for all $v \in \stset$. This leads to a contradiction since it implies that $\decs$ is not the solution of the \eqref{eq:decoder_opt}. Therefore, our initial assumption must be wrong. Thus, there exists $x_0\in \Omega$ such that $h(x_0) = 0$.
\end{proof}
\begin{lemma}\label{lem:norm_infty_bound_2}
Let $\Omega=(-1,1)$. For $h(t) = \decs(t) - f(\encs(t))$ we have:
\begin{align}\label{eq:norm_infty_bound_2_1}
  \norm{h}_{\lp{\infty}} \leqslant 2\norm{h}^\frac{1}{2}_{\lp{2}}\cdot\norm{h'}^\frac{1}{2}_{\lp{2}} < \infty,
\end{align}
and
\begin{align}\label{eq:norm_infty_bound_2_2}
  \norm{h'}_{\lp{\infty}} \leqslant \norm{h'}_{\lp{2}} + \norm{h''}_{\lp{2}} < \infty.
\end{align}
\end{lemma}
\begin{proof}
    Using Lemma~\ref{lem:norm_infty_bound_1} one can conclude
    \begin{align}
        \frac{\norm{h}_{\lp{2}}}{\norm{h'}_{\lp{2}}} \leqslant \sqrt{2}.
    \end{align}
    Since $h \in \sob{2}{2}$ we can apply Corollary~\ref{col:interp_ineq} and Theorem~\ref{th:interp_ineq} with $r=\infty$ and $p,q=2$ to complete the proof of \eqref{eq:norm_infty_bound_2_1}. Furthermore, using Theorem~\ref{th:interp_ineq} with $r=\infty$ and $p,q=2$ and $\ell = 1$ completes the proof of \eqref{eq:norm_infty_bound_2_2}.
\end{proof}
% \begin{lemma}\label{lem:dis_to_cont}
%     Let $\Omega = (a,b)$ be an open bounded interval on $\mathbb{R}$, $\alpha_k = a + k\frac{(b-a)}{K+1}$ for $k \in [K]$. Then:
%     \begin{align}
%         \frac{1}{K}\sum^K_{k=1} h^2(\alpha_k) \leqslant \frac{K+1}{K(b-a)}\norm{h}^2_{\lp{2}}+ \frac{(b-a)}{K+1}\cdot \norm{hh'}_{\lp{\infty}}.
%     \end{align}
% \end{lemma}
% \begin{proof}
%     \begin{align}\label{lg_int_bound_2}
%     \frac{(b-a)}{K+1}\cdot\sum^K_{k=1} h^2(\alpha_k) &\leqslant \int_\Omega h^2(t) \,dt + \Abs{\frac{(b-a)}{K+1}\cdot\sum^K_{k=1} h^2(\alpha_k) - \int_\Omega h^2(t) \,dt}
%         \nonumber \\ 
%         &\lec{(a)}{=}  \int_\Omega h^2(t) \,dt + \Abs {\sum_{k=1}^K h^2(\alpha_k) (\alpha_{k} - \alpha_{k-1}) -\int_\Omega h^2(t) \,dt} \\ \nonumber
%         &\lec{}{=} \int_\Omega h^2(t) \,dt + \Abs {\sum_{k=1}^K \int_{\alpha_{k-1}}^{\alpha_k}h^2(\alpha_k) \,dt -\sum_{k=1}^K \int_{\alpha_{k-1}}^{\alpha_k} h^2(t) \,dt } \\ \nonumber
%         &=  \int_\Omega h^2(t) \,dt + \Abs {\sum_{k=1}^K \int_{\alpha_{k-1}}^{\alpha_k} (h^2(\alpha_k) - h^2(t) )\,dt
%         } \\ \nonumber
%         &\lec{(b)}{=}  \int_\Omega h^2(t) \,dt + \Abs { \sum_{k=1}^K \int_{\alpha_{k-1}}^{\alpha_k} (h^2)'(\eta_k) (\alpha_k - t)\,dt
%         }, \nonumber \\
%         &\lec{(c)}{\leqslant} \int_\Omega h^2(t) \,dt + 2\norm{hh'}_{\lp{\infty}} \sum_{k=1}^K \int_{\alpha_{k-1}}^{\alpha_k}(\alpha_k - t)\, \mathrm dt \nonumber \\ 
%         &=\int_\Omega h^2(t) \,dt + 2\norm{hh'}_{\lp{\infty}} \cdot\frac{(b-a)^2}{2K},
%     \end{align}
%     where (a) follows by $\alpha_k - \alpha_{k-1} = \frac{(b-a)}{K+1}$ with $\alpha_0 := a$, (b) follows by Lagrange’s mean value theorem with  $\eta_k \in (\alpha_{k-1}, \alpha_k)$, and (c) follows by $(h^2)'(\cdot) = 2h(\cdot)h'(\cdot)$. Thus:
%     $$
%     \frac{1}{K}\sum^K_{k=1} h^2(\alpha_k) \leqslant \frac{K+1}{K(b-a)}\norm{h}^2_{\lp{2}}+ \frac{(b-a)}{K+1}\cdot \norm{hh'}_{\lp{\infty}}.
%     $$
% \end{proof}

% \begin{corollary}\label{cor:dis_to_cont_cheb}
%     Let $\Omega = (-1,1)$ and $\alpha_k = \cos(\frac{(2k-1)\pi}{2K})$ be the Chebyshev points of the first kind for $k \in [K]$. Then:
%     \begin{align}
%         \frac{1}{K} \sum^K_{k=1} h^2(\alpha_k) \leqslant  \sin^{-1}(\cos(\frac{\pi}{K})) \cdot \frac{K+1}{2K}\norm{h}^2_{\lp{2}} + \frac{1}{K} \cdot \norm{hh'}_{\lp{\infty}}.
%     \end{align}
% \end{corollary}
% \begin{proof}
%     Let us reformulate the first kind of Chebyshev points based on equidistant points $\halp_k = -1 + \frac{2k}{K+1}$ for $k \in [K]$:
%     \begin{align}
%         \alpha_k = \cos(\frac{(2k-1)\pi}{2K}) = \cos\left(\gamma(\halp_k)\right),
%     \end{align}
%     where $\gamma(t) := \frac{(K+1)\pi}{2K}t + \frac{\pi}{2}$. Applying Lemma~\ref{lem:dis_to_cont} to function $\widetilde{h} := h\left(cos\left(\gamma\left(\halp_k\right)\right)\right)$, we have:
%     \begin{align}
%         \frac{1}{K}\sum^K_{k=1} h^2(\alpha_k) &= \frac{1}{K}\sum^K_{k=1} \widetilde{h}^2(\halp_k) \nonumber \\ &\lec{(a)}{\leqslant} \frac{K+1}{2K}\norm{\widetilde{h}}^2_{\lp{2}}+ \frac{2}{K+1} \cdot \frac{(K+1)\pi}{2K}\cdot  \norm{hh'}_{\lp{\infty}} \nonumber \\
%         &= \frac{K+1}{2K}\norm{\widetilde{h}}^2_{\lp{2}}+ \frac{1}{K} \cdot \norm{hh'}_{\lp{\infty}},
%     \end{align}
%     where (a) is due to the fact that $|\cos'(t)| \leqslant 1$. For $\norm{\widetilde{h}}^2_{\lp{2}}$ we have:
%     \begin{align}\label{eq:newvar_integral}
%         \norm{\widetilde{h}}^2_{\lp{2}} = 
%          \int_{-1}^{1} \biggl(\decs(\cos(\gamma(t))) -f(\encs(\cos(\gamma(t)))\biggr)^2\,dt.
%     \end{align}
%     We can define a new variable denoted by $t'=\cos(\gamma(t))$ and rewrite the LHS with the fact that $dt' = -dt.\frac{(K+1)\pi}{2K}\sin\left(\frac{(K+1)\pi}{2K}t + \frac{\pi}{2}\right) =  -dt.\frac{\pi}{2}\sqrt{1-t'^2}$:
%     \begin{align}\label{eq:newvar_integral}
%          \int_{-1}^{1} \left[\decs(\cos(\gamma(t))) -f(\encs(\cos(\gamma(t)))\right]^2\,dt = \int_{-\cos\left(\frac{\pi}{2K}\right)}^{\cos\left(\frac{\pi}{2K}\right)} \frac{\left(\decs(t')-f(\encs(t'))\right)^2}{\sqrt{1-t'^2}} \,dt'.
%     \end{align}
%      So, the integral range remains symmetric around zero. Thus, we have:
%     \begin{align}\label{eq:newvar_integral_2}
%           \int_{-\cos\left(\frac{\pi}{2K}\right)}^{\cos\left(\frac{\pi}{2K}\right)} \frac{h(t')}{\sqrt{1-t'^2}} \,dt'  &\leqslant \sin^{-1}(\cos(\frac{\pi}{K})) \cdot \int_{-\cos\left(\frac{\pi}{2K}\right)}^{\cos\left(\frac{\pi}{2K}\right)} h(t')^2 \nonumber \\
%    &\leqslant \sin^{-1}(\cos(\frac{\pi}{K})) \cdot \int_{-1}^{1} h(t')^2 \nonumber \\
%    &= \sin^{-1}(\cos(\frac{\pi}{K})) \cdot \norm{h}^2_{\lp{2}}.
%     \end{align}
% Therefore,
% \begin{align}
%      \frac{1}{K} \sum^K_{k=1} h^2(\alpha_k) \leqslant  \sin^{-1}(\cos(\frac{\pi}{K})) \cdot \frac{K+1}{2K}\norm{h}^2_{\lp{2}} + \frac{1}{K} \cdot  \norm{hh'}_{\lp{\infty}}.
% \end{align}  
% \end{proof}
% Using Lemma~\ref{lem:dis_to_cont} we have:
% \begin{align}\label{eq:ldec_step2}
%     \ldec \leqslant \frac{2(K+1)}{K(b-a)}\mathbb{E}_{\epsilon, \stset}\left[\norm{h}^2_{\lp{2}}\right]+ \frac{2(b-a)}{K+1}\cdot\mathbb{E}_{\epsilon, \stset}\left[ \norm{hh'}_{\lp{\infty}}\right].
% \end{align}
% In the subsequent lemmas, we offer an upper bound for the second term in \eqref{eq:ldec_step2}.
% \begin{lemma}\label{lem:norm_infty_bound_3}
% Let $\Omega = (-1,1)$. For $h(t) = \decs(t) - f(\encs(t))$, then $hh^{\prime} \in \sobloc{1}{1}$.
% \end{lemma}
% \begin{proof}
%     Since $\sob{1}{1} \subseteq \sobloc{1}{1}$, it is sufficient to show that $hh' \in \sob{1}{1}$. Therefore we have:
%     \begin{align}
%         \int_\Omega |h(t)h'(t)|\,dt &= \int_\Omega |h(t)||h'(t)|\,dt \lec{(a)}{\leqslant} \left(\int_\Omega |h(t)|^2\,dt\right)^\frac{1}{2}\left(\int_\Omega |h'(t)|^2\,dt\right)^\frac{1}{2} \lec{(b)}{<} \infty \\
%         \int_\Omega |(h(t)h'(t))^{\prime}|\,dt &= \int_\Omega |h(t)h''(t) + h'(t)^2|\,dt \nonumber \\ &\lec{(c)}{\leqslant}  \int_\Omega |h(t)h''(t)|\,dt + \int_\Omega h'(t)^2\,dt \nonumber \\ &\lec{(d)}{\leqslant}  \left(\int_\Omega |h(t)|^2\,dt\right)^\frac{1}{2} \left(\int_\Omega |h''(t)|^2\,dt\right)^\frac{1}{2} + \int_\Omega h'(t)^2\,dt \lec{(e)}{<} \infty,
%     \end{align}
%     where (a) and (d) follow by Cauchy-Schwartz inequality, (c) is because of triangle inequality, and (b), (d) is due to the fact that $h\in \sob{2}{2}$.
% \end{proof}
% \begin{lemma}\label{lem:norm_infty_bound_4}
% Under the same assumptions as Lemma~\ref{lem:norm_infty_bound_3}, we have $h(t)h'(t) \in \lp{2}$.
% \end{lemma}
% \begin{proof}
% Based on previous results we have:
% \begin{align}
%     \int_\Omega \left(|h(t)h''(t) + h'(t)^2|\,dt\right)^2 \lec{(a)}{\leqslant} &\int_\Omega \left(|h(t)h''(t)| + h'(t)^2\right)^2\,dt \nonumber \\
%     \lec{(b)}{\leqslant}  &2\int_\Omega |h(t)|^2|h''(t)|^2 + h'(t)^4\,dt \nonumber \\
%     \lec{(c)}{\leqslant} &2\mathcal{L}(\Omega)\cdot \norm{h}^2_{\lp{\infty}} \int_\Omega \left(|h''(t)|^2\right) + \nonumber \\ &2\mathcal{L}(\Omega)\cdot \norm{h'}^2_{\lp{\infty}} \int_\Omega \left(|h'(t)|^2\right) \lec{(d)}{<} \infty,
% \end{align}
% where (a) comes from triangle inequality, (b) follows by AM-GM inequality, (c) follows by $h(t) \leqslant \norm{h}_{\lp{\infty}}$ and  $h'(t) \leqslant \norm{h'}_{\lp{\infty}}$ for $t \in \Omega$, and finally (d) is because the fact that $h \in \sob{2}{2}$ and $\norm{h}_{\lp{\infty}}, \norm{h'}_{\lp{\infty}} < \infty$ using Lemma~\ref{lem:norm_infty_bound_2}.
% \end{proof}
Building upon Lemma~\ref{lem:norm_infty_bound_2} and starting from \eqref{eq:ldec}, we can derive an upper bound for $\ldec$:
\begin{align}\label{eq:ldec_step3}
    \ldec(\decs) &= \mathbb{E}_{\epsilon, \stset} \left[\frac{2}{K} \sum^K_{k=1} \left(\decs(\alpha_k) -f(\encs(\alpha_k))\right)^2_2\right]
    \nonumber \\ 
    &\lec{(a)}{=} \frac{2}{K} \sum^K_{k=1} \mathbb{E}_{\epsilon, \stset} \left[h(\alpha_k)^2
    \right] 
    \nonumber \\
    &\lec{(b)}{\leqslant} \frac{2}{K} \sum^K_{k=1} \mathbb{E}_{\epsilon, \stset} \left[\norm{h}^2_{\lp{\infty}}
    \right] 
    \nonumber \\
    &\lec{(c)}{\leqslant} 2 \mathbb{E}_{\epsilon, \stset} \left[\norm{h}_{\lp{2}}\cdot \norm{h'}_{\lp{2}}
    \right]
    \nonumber \\
    &\lec{(d)}{\leqslant} 2 \mathbb{E}_{\epsilon, \stset} \left[\norm{h}^2_{\lp{2}}\right]^\frac{1}{2}\cdot \mathbb{E}_{\epsilon, \stset} \left[ \norm{h'}^2_{\lp{2}}
    \right]^\frac{1}{2},
\end{align}
where (a) follows from the definition of $h(t)$, (b) is due to the fact that $h(t) \leqslant \norm{h}_{\lp{\infty}}$ for $t \in \Omega$, (c) follows by applying Lemma~\ref{lem:norm_infty_bound_2}, and (d) is a result of applying the Cauchy-Schwartz inequality.

\subsection{Proof of Theorem~\ref{th:conv_rate_noiseless}}\label{app:proof_noiseless}
\begin{proof}
As previously mentioned, $\decs(\cdot)$ is a second-order smoothing spline fitted on the data points $\left\{(\beta_{i_1}, f(\encs(\beta_{i_1})), \dots, (\beta_{i_{|\stset|}}, f(\encs(\beta_{i_{|\stset|}}))) \right\}$, where $\stset := \left\{\beta_{i_1}, \dots, \beta_{i_{|\stset|}} \right\}$ represents the set of non-straggler worker nodes, and $\bm{f \circ \encs} := \left\{f(\encs(\beta_{i_1})), \dots, f(\encs(\beta_{i_{|\stset|}})) \right\}$ is the corresponding set of computation results from these non-straggler workers. By the definition given in \eqref{eq:ssfunction_def}, $\splineCC(\cdot)$ denotes the smoothing spline operator for the decoder layer. Hence, we have the following:
\begin{alignat}{2}
    \mathbb{E}_{ \stset}\left[\norm{h}^2_{\lp{2}}\right] &= \mathbb{E}_{\stset}&&\left[\norm{f\circ\encs - \splineCC[\mathbf{f}]}^2_{\lp{2}}\right],
\end{alignat}
where $\mathbf{f}=\left\{f(\encs(\beta_{i_1})),\dots,f(\encs(\beta_{i_{|\stset|}}))\right\}$. Let us define the following variables analogous to those in \eqref{eq:max_dist_st}:
\begin{align}\label{eq:delta_max_stset}
\Delta^\stset_\textrm{max}:=\underset{f \in \left\{0,\dots,|\stset|\right\}}{\max} \left\{\beta_{i_{f+1}}-\beta_{i_f}\right\}, \quad \Delta^\stset_\textrm{min}:=\underset{f \in \left\{1,\dots,|\stset|-1\right\}}{\min} \left\{\beta_{i_{f+1}}-\beta_{i_f}\right\}.
\end{align}
Since there are at most $S$ stragglers among the worker nodes, we have $\Delta^\stset_\textrm{max} \leqslant (S+1) \cdot \Delta_\textrm{max}$ and $ \frac{\Delta^\stset_\textrm{max}}{\Delta^\stset_\textrm{min}} \leqslant (S+1)\cdot\frac{\Delta_\textrm{max}}{\Delta_\textrm{min}} \leqslant (S+1)B$. Additionally, because $\Delta_\textrm{min} \leqslant \frac{2}{N}$, there exists a constant $J$ such that $\Delta_\textrm{max} \leqslant \frac{J}{N}$, and consequently, $\Delta^\stset_\textrm{max} \leqslant \frac{J(S+1)}{N} \leqslant \frac{J(S+1)}{N-S}$. Otherwise, this would contradict the condition $\frac{\Delta_\textrm{max}}{\Delta_\textrm{min}} \leqslant B$.

Applying Theorem~\ref{th:lit_spline_noiseless} with $\Omega=(-1, 1), m=2$, we have
\begin{align}\label{eq:h_1}
    &\mathbb{E}_{ \stset}\left[\norm{h}^2_{\lp{2}}\right] \lec{}{\leqslant}H_0\norm{(f\circ\encs)^{(2)}}^2_{\lp{2}}\cdot\mathbb{E}_{\stset}\left[L\right],
\end{align}
and
\begin{align}\label{eq:hp_1}
\mathbb{E}_{\stset}\left[\norm{h'}^2_{\lp{2}}\right] &\lec{}{\leqslant}H_1\norm{(f\circ\encs)^{(2)}}^2_{\lp{2}}\cdot\mathbb{E}_{\stset}\left[L^{\frac{1}{2}}(1 + \frac{L}{16})^{\frac{1}{2}}\right],
\end{align}
where $L=p_2\left(\frac{\Delta^\stset_\textrm{max}}{\Delta^\stset_\textrm{min}}\right)\cdot\frac{(N-S) \Delta^\stset_\textrm{max}}{4} \declamb + D(2)\cdot \left({\Delta^\stset_\textrm{max}}\right)^{4}$ and $H_0, H_1:= H(2,0), H(2,1)$ as defined in  Theorem~\ref{th:lit_spline_noiseless}. Thus, we have:
\begin{align}\label{eq:bound_L}
    \mathbb{E}_{\stset}[L] &\leqslant \mathbb{E}_{\stset}\left[p_2\left(\frac{\Delta^\stset_\textrm{max}}{\Delta^\stset_\textrm{min}}\right)\cdot\frac{(N-S) \Delta^\stset_\textrm{max}}{4} \declamb + D\cdot \left({\Delta^\stset_\textrm{max}}\right)^{4} \right] \nonumber \\ 
    &\lec{(a)}{\leqslant} \mathbb{E}_{\stset}\left[p_2\left(\frac{\Delta^\stset_\textrm{max}}{\Delta^\stset_\textrm{min}}\right)\right]\cdot\frac{J(S+1)}{4} \declamb + D\cdot \mathbb{E}_{\stset}\left[\left({\Delta^\stset_\textrm{max}}\right)^{4}\right] 
    \\ &\lec{(b)}{\leqslant} \mathbb{E}_{\stset}\left[p_2\left(\frac{\Delta^\stset_\textrm{max}}{\Delta^\stset_\textrm{min}}\right)\right]\cdot\frac{J(S+1)}{4} \declamb + DJ^4\frac{(S+1)^4}{N^4},
\end{align}
where $D:=D(2)$ and (a) and (b) follow from  $\Delta^\stset_\textrm{max} \leqslant \frac{J(S+1)}{N-S}$ and $\Delta^\stset_\textrm{max} \leqslant \frac{J(S+1)}{N}$ respectively. Substituting in \eqref{eq:h_1} and \eqref{eq:hp_1}, we have:
\begin{alignat}{2}
    \label{eq:h_noiseless_bound_1}
    \mathbb{E}_{\stset}\left[\norm{h}^2_{\lp{2}}\right] &\leqslant H_0\norm{(f\circ\encs)^{(2)}}^2_{\lp{2}}\cdot&&\left(\mathbb{E}_{\stset}\left[p_2\left(\frac{\Delta^\stset_\textrm{max}}{\Delta^\stset_\textrm{min}}\right)\right]\cdot\frac{J(S+1)}{4} \declamb + DJ^4\frac{(S+1)^4}{N^4}\right),  \\ \label{eq:hp_noiseless_bound_1}
    \mathbb{E}_{\stset}\left[\norm{h'}^2_{\lp{2}}\right] &\leqslant H_1\norm{(f\circ\encs)^{(2)}}^2_{\lp{2}}\cdot &&\left(\mathbb{E}_{\stset}\left[p_2\left(\frac{\Delta^\stset_\textrm{max}}{\Delta^\stset_\textrm{min}}\right)\right]\cdot\frac{J(S+1)}{4} \declamb + DJ^4\frac{(S+1)^4}{N^4}\right)^{\frac{1}{2}} \nonumber \\ & &&\cdot\left(1 + \frac{\mathbb{E}_{\stset}\left[p_2\left(\frac{\Delta^\stset_\textrm{max}}{\Delta^\stset_\textrm{min}}\right)\right]\cdot\frac{J(S+1)}{4} \declamb + DJ^4\frac{(S+1)^4}{N^4}}{16}\right)^{\frac{1}{2}}.
\end{alignat}
Therefore, we can derive an upper bound for \eqref{eq:ldec_step3} based on the above inequalities. This upper bound holds for any $\declamb$. Since $\declamb \leqslant N^{-4}$ and $\frac{\Delta^\stset_\textrm{max}}{\Delta^\stset_\textrm{min}} \leqslant (S+1)B$, we have:
\begin{align}
\mathbb{E}_{\stset}\left[p_2\left(\frac{\Delta^\stset_\textrm{max}}{\Delta^\stset_\textrm{min}}\right)\right]\cdot\frac{J(S+1)}{4} \declamb + DJ^4\frac{(S+1)^4}{N^4} &\leqslant \Tilde{p}_3\left(S+1\right)N^{-4} + DJ^4\frac{(S+1)^4}{N^4} \nonumber \\
\lec{(a)}{\leqslant} \widetilde{D}\frac{(S+1)^4}{N^4},
\end{align}
where $\Tilde{p}_3$ is a degree-3 polynomial in $(S+1)$ with positive constant coefficients, and $\widetilde{D}$ is the sum of the coefficients of $\Tilde{p}_3$ and $DJ^4$. Therefore, we have:
\begin{alignat}{2} \label{eq:ldec_noiseless_bound_1}
    \ldec &\leqslant \norm{(f\circ\encs)^{(2)}}^2_{\lp{2}} &&\Biggl[\left(H_0\cdot r(S, N)\right)^\frac{1}{2}\left(H_1\cdot r(S, N)^\frac{1}{2}\left(1 + \frac{r(S, N)}{16}\right)^\frac{1}{2}\right)^\frac{1}{2}
    \Biggr]
    \nonumber \\
    &\lec{}{=} \norm{(f\circ\encs)^{(2)}}^2_{\lp{2}} &&\Biggl[H_0^\frac{1}{2}H_1^\frac{1}{2}\cdot r(S, N)^\frac{3}{4}\cdot \left(1 + \frac{r(S, N)}{16}\right)^\frac{1}{4}
    \Biggr],
\end{alignat}
where $r(S, N):=\widetilde{D}\frac{(S+1)^4}{N^4}$. Note that, since $S+1 \leqslant N$ then $1 + \frac{r(S, N)}{16} \leqslant 2\max(1, \widetilde{D})$. Defining $\eta:=2\max(1, \widetilde{D})$, we have:
\begin{align}
    \ldec &\lec{}{\leqslant} \norm{(f\circ\encs)^{(2)}}^2_{\lp{2}} \Biggl[H_0^\frac{1}{2}H_1^\frac{1}{2} \eta^\frac{1}{4} \cdot r(S, N)^\frac{3}{4}\Biggr],
\end{align}
Defining $C := H_0^\frac{1}{2}H_1^\frac{1}{2} \eta^\frac{1}{4}$ and applying Lemma~\ref{lem:lt}, completes the proof.
\end{proof}



% \subsection{Proof of Theorem~\ref{th:conv_rate_noiseless}}\label{app:proof_noiseless}
% \begin{proof}
% As previously mentioned, $\decs(\cdot)$ is a second-order smoothing spline fitted on the data points $\left\{(\beta_{i_1}, f(\encs(\beta_{i_1})), \dots, (\beta_{i_{|\stset|}}, f(\encs(\beta_{i_{|\stset|}}))) \right\}$, where $\stset := \left\{\beta_{i_1}, \dots, \beta_{i_{|\stset|}} \right\}$ represents the set of non-straggler worker nodes, and $\bm{f \circ \encs} := \left\{f(\encs(\beta_{i_1})), \dots, f(\encs(\beta_{i_{|\stset|}})) \right\}$ is the corresponding set of computation results from these non-straggler workers. By the definition given in \eqref{eq:ssfunction_def}, $\splineCC(\cdot)$ denotes the smoothing spline operator for the decoder layer. Hence, we have the following:
% \begin{alignat}{2}
%     \mathbb{E}_{ \stset}\left[\norm{h}^2_{\lp{2}}\right] &= \mathbb{E}_{\stset}&&\left[\norm{f\circ\encs - \splineCC[\mathbf{f}]}^2_{\lp{2}}\right],
% \end{alignat}
% where $\mathbf{f}=\left\{f(\encs(\beta_{i_1})),\dots,f(\encs(\beta_{i_{|\stset|}}))\right\}$. Let us define the following variables analogous to those in \eqref{eq:max_dist_st}:
% \begin{align}\label{eq:delta_max_stset}
% \Delta^\stset_\textrm{max}:=\underset{f \in \left\{0,\dots,|\stset|\right\}}{\max} \left\{\beta_{i_{f+1}}-\beta_{i_f}\right\}, \quad \Delta^\stset_\textrm{min}:=\underset{f \in \left\{1,\dots,|\stset|-1\right\}}{\min} \left\{\beta_{i_{f+1}}-\beta_{i_f}\right\}.
% \end{align}
% Since there are at most $S$ stragglers among worker nodes, $\Delta^\stset_\textrm{max} \leqslant (S+1) \cdot \Delta_\textrm{max}$ and $ \frac{\Delta^\stset_\textrm{max}}{\Delta^\stset_\textrm{min}} \leqslant (S+1)\cdot\frac{\Delta_\textrm{max}}
% {\Delta_\textrm{min}}$. Additionally because $\Delta_\textrm{max} = \mathcal{O}(\frac{1}{N})$, there exists a constant $J$ such that $\Delta_\textrm{max} \leqslant \frac{J}{N}$ and correspondingly $\Delta^\stset_\textrm{max} \leqslant \frac{J(S+1)}{N} \leqslant \frac{J(S+1)}{N-S}$.
% Applying Theorem~\ref{th:lit_spline_noiseless} with $\Omega=(-1, 1), m=2$, we have
% \begin{align}\label{eq:h_1}
%     &\mathbb{E}_{ \stset}\left[\norm{h}^2_{\lp{2}}\right] \lec{}{\leqslant}H_0\norm{(f\circ\encs)^{(2)}}^2_{\lp{2}}\cdot\mathbb{E}_{\stset}\left[L\right],
% \end{align}
% and
% \begin{align}\label{eq:hp_1}
% \mathbb{E}_{\stset}\left[\norm{h'}^2_{\lp{2}}\right] &\lec{}{\leqslant}H_1\norm{(f\circ\encs)^{(2)}}^2_{\lp{2}}\cdot\mathbb{E}_{\stset}\left[L^{\frac{1}{2}}(1 + \frac{L}{16})^{\frac{1}{2}}\right],
% \end{align}
% where $L=p_2\left(\frac{\Delta^\stset_\textrm{max}}{\Delta^\stset_\textrm{min}}\right)\cdot\frac{(N-S) \Delta^\stset_\textrm{max}}{4} \declamb + D(2)\cdot \left({\Delta^\stset_\textrm{max}}\right)^{4}$ and $H_0, H_1:= H(2,0), H(2,1)$ as defined in  Theorem~\ref{th:lit_spline_noiseless}. Thus, we have:
% \begin{align}\label{eq:bound_L}
%     \mathbb{E}_{\stset}[L] &\leqslant \mathbb{E}_{\stset}\left[p_2\left(\frac{\Delta^\stset_\textrm{max}}{\Delta^\stset_\textrm{min}}\right)\cdot\frac{(N-S) \Delta^\stset_\textrm{max}}{4} \declamb + D\cdot \left({\Delta^\stset_\textrm{max}}\right)^{4} \right] \nonumber \\ 
%     &\lec{(a)}{\leqslant} \mathbb{E}_{\stset}\left[p_2\left(\frac{\Delta^\stset_\textrm{max}}{\Delta^\stset_\textrm{min}}\right)\right]\cdot\frac{J(S+1)}{4} \declamb + D\cdot \mathbb{E}_{\stset}\left[\left({\Delta^\stset_\textrm{max}}\right)^{4}\right] 
%     \\ &\lec{(b)}{\leqslant} \mathbb{E}_{\stset}\left[p_2\left(\frac{\Delta^\stset_\textrm{max}}{\Delta^\stset_\textrm{min}}\right)\right]\cdot\frac{J(S+1)}{4} \declamb + DJ^4\frac{(S+1)^4}{N^4},
% \end{align}
% where $D:=D(2)$ and (a) and (b) follow from  $\Delta^\stset_\textrm{max} \leqslant \frac{J(S+1)}{N-S}$ and $\Delta^\stset_\textrm{max} \leqslant \frac{J(S+1)}{N}$ respectively. Substituting in \eqref{eq:h_1},\eqref{eq:hp_1}, and \eqref{eq:hp_1}, we have:
% \begin{alignat}{2}
%     \label{eq:h_noiseless_bound_1}
%     \mathbb{E}_{\stset}\left[\norm{h}^2_{\lp{2}}\right] &\leqslant H_0\norm{(f\circ\encs)^{(2)}}^2_{\lp{2}}\cdot&&\left(\mathbb{E}_{\stset}\left[p_2\left(\frac{\Delta^\stset_\textrm{max}}{\Delta^\stset_\textrm{min}}\right)\right]\cdot\frac{J(S+1)}{4} \declamb + DJ^4\frac{(S+1)^4}{N^4}\right),  \\ \label{eq:hp_noiseless_bound_1}
%     \mathbb{E}_{\stset}\left[\norm{h'}^2_{\lp{2}}\right] &\leqslant H_1\norm{(f\circ\encs)^{(2)}}^2_{\lp{2}}\cdot &&\left(\mathbb{E}_{\stset}\left[p_2\left(\frac{\Delta^\stset_\textrm{max}}{\Delta^\stset_\textrm{min}}\right)\right]\cdot\frac{J(S+1)}{4} \declamb + DJ^4\frac{(S+1)^4}{N^4}\right)^{\frac{1}{2}} \nonumber \\ & &&\cdot\left(1 + \frac{\mathbb{E}_{\stset}\left[p_2\left(\frac{\Delta^\stset_\textrm{max}}{\Delta^\stset_\textrm{min}}\right)\right]\cdot\frac{J(S+1)}{4} \declamb + DJ^4\frac{(S+1)^4}{N^4}}{16}\right)^{\frac{1}{2}}.
% \end{alignat}
% Therefore, we can derive an upper bound for \eqref{eq:ldec_step3} based on the above inequalities. This upper bound holds for any $\declamb$. To obtain an upper bound independent of $\declamb$, we minimize the right-hand side of this bound with respect to $\declamb$. Since all upper bounds in \eqref{eq:hp_noiseless_bound_1} and \eqref{eq:h_noiseless_bound_1} are increasing functions of $\declamb$, and the upper bound of $\ldec$ in \eqref{eq:ldec_step3} is an increasing function of $\mathbb{E}_{\stset}\left[\norm{h}^2_{\lp{2}}\right]$ and $\mathbb{E}_{\stset}\left[\norm{h'}^2_{\lp{2}}\right]$, setting $\declamb^* = 0$ minimizes the right-hand side of \eqref{eq:ldec_step3} with respect to $\declamb$. Substituting $\declamb = 0$, we obtain:
% \begin{alignat}{2}\label{eq:ldec_noiseless_bound_1}
%     \ldec &\leqslant \norm{(f\circ\encs)^{(2)}}^2_{\lp{2}} &&\Biggl[\left(H_0\cdot r(S, N)\right)^\frac{1}{2}\left(H_1\cdot r(S, N)^\frac{1}{2}\left(1 + \frac{r(S, N)}{16}\right)^\frac{1}{2}\right)^\frac{1}{2}
%     \Biggr]
%     \nonumber \\
%     &\lec{}{=} \norm{(f\circ\encs)^{(2)}}^2_{\lp{2}} &&\Biggl[H_0^\frac{1}{2}H_1^\frac{1}{2}\cdot r(S, N)^\frac{3}{4}\cdot \left(1 + \frac{r(S, N)}{16}\right)^\frac{1}{4}
%     \Biggr],
% \end{alignat}
% % \begin{alignat}{2}\label{eq:ldec_noiseless_bound_1}
% %     \ldec &\leqslant \norm{(f\circ\encs)^{(2)}}^2_{\lp{2}} &&\Biggl[\frac{H_0(K+1)}{K} \cdot r(S, N) + \frac{16H_0^{\frac{1}{4}}H_1^{\frac{1}{4}}}{K+1} \cdot r(S, N)^{\frac{3}{8}} \cdot \left(1 + \frac{r(S, N)}{16}\right)^{\frac{1}{8}}
% %     \nonumber \\ & 
% %     && \cdot \left(H_1r(S, N)^\frac{1}{2}\left(1 + \frac{r(S, N)}{16}\right)^\frac{1}{2} + H_2\left(1+\frac{r(S,N)}{16}\right)\right)^{\frac{1}{2}}
% %     \Biggr],
% % \end{alignat}
% where $r(S, N):=DJ^4\frac{(S+1)^4}{N^4}$. Note that, since $S+1 \leqslant N$ then $1 + \frac{r(S, N)}{16} \leqslant 2\max(1, DJ^4)$. Defining $\eta:=2\max(1, DJ^4)$, we have:
% \begin{align}
%     \ldec &\lec{}{\leqslant} \norm{(f\circ\encs)^{(2)}}^2_{\lp{2}} \Biggl[H_0^\frac{1}{2}H_1^\frac{1}{2} \eta^\frac{1}{4} \cdot r(S, N)^\frac{3}{4}\Biggr],
% \end{align}
% Defining $C := H_0^\frac{1}{2}H_1^\frac{1}{2} \eta^\frac{1}{4}$ and applying Lemma~\ref{lem:lt}, completes the proof.
% % \begin{align}
% %     \ldec &\leqslant \norm{(f\circ\encs)^{(2)}}^2_{\lp{2}} \Biggl[\frac{H_0(K+1)}{K} \cdot r(S, N) + \frac{16H_0^{\frac{1}{4}}H_1^{\frac{1}{4}}\eta^\frac{1}{8}}{K+1} &&\cdot r(S, N)^{\frac{3}{8}} \cdot 
% %     \nonumber \\ & 
% %     &&\left(H_1D^\frac{1}{2}J^2\eta^\frac{1}{2} + H_2\eta\right)^{\frac{1}{2}}
% %     \Biggr].
% % \end{align}
% % Defining $C_ := H_0$ and $C_2 := 16H_0^{\frac{1}{4}}H_1^{\frac{1}{4}}\eta^{\frac{1}{8}} \cdot \left(H_1D^{\frac{1}{2}}J^2\eta^{\frac{1}{2}} + H_2\eta\right)^{\frac{1}{2}}$, and applying Lemma~\ref{lem:lt}, we complete the proof.
% \end{proof}
% % and (a) follows from $1 + \frac{r(S, N)}{16} \leqslant \eta\cdot DJ^4$ for some $\eta > 0$. Therefore, defining $C_2:=\frac{(K+1)}{K} H_0$ and $C_1:=\frac{8\delta \eta D^{\frac{1}{4}}J}{K+1} H_0^{\frac{1}{2}}H_1^{\frac{1}{2}}$ as well as using Lemma~\ref{lem:lt} complete the proof.
\subsection{Proof of Theorem~\ref{th:conv_rate_noisy}}\label{sec:proof_noisy}
Using the decomposition \eqref{eq:smooth_spline_linear_eq} we have:
\begin{alignat}{2}
    \mathbb{E}_{\epsilon, \stset}\left[\norm{h}^2_{\lp{2}}\right] &= \mathbb{E}_{\stset}&&\left[\norm{f\circ\encs - \splineCC[\mathbf{f}]}^2_{\lp{2}}\right] \nonumber \\ & &&+  \mathbb{E}_{\epsilon, \stset}\left[\norm{f\circ\encs - \splineCC[\bm{\epsilon}]}^2_{\lp{2}}\right],
\end{alignat}
where $\mathbf{f}=\left\{f(\encs(\beta_{i_1})),\dots,f(\encs(\beta_{i_{|\stset|}}))\right\}$ and $\bm{\epsilon} = \left\{\epsilon_{i_1},\dots, \epsilon_{i_{|\stset|}}\right\}$. Same as \eqref{eq:delta_max_stset} we define the following variables:
\begin{align}
\Delta^\stset_\textrm{max}:=\underset{f \in \left\{0,\dots,|\stset|\right\}}{\max} \left\{\beta_{i_{f+1}}-\beta_{i_f}\right\}, \quad \Delta^\stset_\textrm{min}:=\underset{f \in \left\{1,\dots,|\stset|-1\right\}}{\min} \left\{\beta_{i_{f+1}}-\beta_{i_f}\right\}.
\end{align}
Again we have $\Delta^\stset_\textrm{max} \leqslant (S+1) \cdot \Delta_\textrm{max}$ and $ \frac{\Delta^\stset_\textrm{max}}{\Delta^\stset_\textrm{min}} \leqslant (S+1)\cdot\frac{\Delta_\textrm{max}}
{\Delta_\textrm{min}} \lec{(a)}{\leqslant} (S+1)B$, where (a) is because of the bounded condition that we have. Therefore, $\frac{\Delta^\stset_\textrm{max}}{\Delta^\stset_\textrm{min}}$ is bounded. Additionally, since $\Delta_\textrm{min} \leqslant \frac{2}{N}$, then $\frac{\Delta_\textrm{max}}{\Delta_\textrm{min}} \leqslant B$ implies that both $\Delta_\textrm{max} = \mathcal{O}(\frac{1}{N})$. Thus, there exists a constant $J$ such that $\Delta_\textrm{max} \leqslant \frac{J}{N}$. Therefore, we have $\Delta^\stset_\textrm{max} \leqslant \Delta_\textrm{max}\cdot (S+1) \leqslant \frac{J(S+1)}{N} \leqslant \frac{J(S+1)}{N-S}$.
Applying Theorems~\ref{th:lit_spline_noiseless} and \ref{th:lit_spline_noisy} with $\Omega=(-1, 1), m=2$, we have:
\begin{align}\label{eq:h_2}
    \mathbb{E}_{\epsilon, \stset}\left[\norm{h}^2_{\lp{2}}\right] &\lec{}{\leqslant}H_0\norm{(f\circ\encs)^{(2)}}^2_{\lp{2}}\cdot\mathbb{E}_{\stset}\left[L\right] + \frac{Q_0\sigma_0^2}{N-S}\declamb^{-\frac{1}{4}},
\end{align}
and 
\begin{align}\label{eq:hp_2}
    \mathbb{E}_{\epsilon, \stset}\left[\norm{h'}^2_{\lp{2}}\right] &\lec{}{\leqslant}H_1\norm{(f\circ\encs)^{(2)}}^2_{\lp{2}}\cdot\mathbb{E}_{\stset}\left[L^{\frac{1}{2}}(1 + \frac{L}{16})^{\frac{1}{2}}\right] + \frac{Q_1\sigma_0^2}{N-S}\declamb^{-\frac{3}{4}},
\end{align}
where $L=p_2\left(\frac{\Delta^\stset_\textrm{max}}{\Delta^\stset_\textrm{min}}\right)\cdot\frac{(N-S) \Delta^\stset_\textrm{max}}{4} \declamb + D(2)\cdot \left({\Delta^\stset_\textrm{max}}\right)^{4}$, $H_0, H_1 := H(2,0), H(2,1)$, and $Q_0(\lambda_0), Q_1(\lambda_0):=Q(2,0,\lambda_0), Q(2,1,\lambda_0)$ as defined in  Theorems~\ref{th:lit_spline_noiseless} and ~\ref{th:lit_spline_noisy}. 
Since $\frac{\Delta^\stset_\textrm{max}}{\Delta^\stset_\textrm{min}} \leqslant (S+1)B$, we have:
\begin{align}\label{eq:bound_L_2}
    \mathbb{E}_{\stset}[L] &\leqslant \mathbb{E}_{\stset}\left[\Tilde{p_2}(S+1)\cdot\frac{(N-S) \Delta^\stset_\textrm{max}}{4} \declamb + D\cdot \left({\Delta^\stset_\textrm{max}}\right)^{4} \right] 
    \nonumber \\ 
    &\lec{(a)}{\leqslant} \Tilde{p_2}(S+1)\cdot\frac{J(S+1)}{4} \declamb + DJ^4\frac{(S+1)^4}{(N-S)^4} 
    \nonumber \\ 
    &\lec{}{=} p_3(S+1)\cdot\declamb + DJ^4\frac{(S+1)^4}{(N-S)^4},
\end{align}
where $D:=D(2), \Tilde{p_2}(S+1) = p_2\left(B(S+1)\right)$, $p_3(S+1) := \Tilde{p_2}(S+1)\cdot\frac{J(S+1)}{4}$ is a degree three polynomial of $(S+1)$, and (a) follows from the  $\Delta^\stset_\textrm{max} \leqslant \frac{J(S+1)}{N-S}$. Substituting in \eqref{eq:h_2}, we have:
\begin{align}\label{eq:bound_h_2}
    \mathbb{E}_{\epsilon, \stset}\left[\norm{h}^2_{\lp{2}}\right] &\lec{}{\leqslant}H_0\norm{(f\circ\encs)^{(2)}}^2_{\lp{2}} 
    \left(p_3(S+1)\cdot\declamb + DJ^4\frac{(S+1)^4}{(N-S)^4}\right) + \frac{Q_0(\lambda_0)\sigma_0^2}{N-S}\declamb^{-\frac{1}{4}} \nonumber \\ 
    &\lec{(a)}{\leqslant} H_0\norm{(f\circ\encs)^{(2)}}^2_{\lp{2}}\declamb\cdot 
    \left(p_3(S+1) + DJ^4(S+1)^4\right) + \frac{Q_0(\lambda_0)\sigma_0^2}{N-S}\declamb^{-\frac{1}{4}} \nonumber \\ &\lec{(b)}{=} H_0\norm{(f\circ\encs)^{(2)}}^2_{\lp{2}}\declamb\cdot 
    p_4(S+1) + \frac{Q_0(\lambda_0)\sigma_0^2}{N-S}\declamb^{-\frac{1}{4}},
\end{align}
where (a) follows from the fact that $\declamb^{-1} (N-S)^{-4} \leqslant 1$, as assumed in Theorem~\ref{th:lit_spline_noisy} and (b) is by definition $p_4(S+1) := p_3(S+1) + DJ^4(S+1)^4$ is a degree four polynomial of $(S+1)$. An analogous upper bound can be derived for $\mathbb{E}_{\epsilon, \stset}\left[\norm{h'}^2_{\lp{2}}\right]$ as
\begin{alignat}{2}\label{eq:bound_hp_2}
    \mathbb{E}_{\epsilon, \stset}\left[\norm{h'}^2_{\lp{2}}\right] &\leqslant H_1\norm{(f\circ\encs)^{(2)}}^2_{\lp{2}}\cdot\declamb^{\frac{1}{2}}\cdot p_4(S+1)^{\frac{1}{2}} &&\cdot\left(1 + \declamb\frac{p_4(S+1)}{16}\right)^{\frac{1}{2}} \nonumber \\ 
    & &&+  \frac{Q_1(\lambda_0)\sigma_0^2}{N-S}\declamb^{-\frac{3}{4}} \nonumber \\ 
    &\lec{(a)}{\leq} H_1\norm{(f\circ\encs)^{(2)}}^2_{\lp{2}}\cdot\declamb^{\frac{1}{2}}\cdot p_4(S+1)^{\frac{1}{2}} &&\cdot\left(1 + \lambda_0\frac{p_4(S+1)}{16}\right)^{\frac{1}{2}} \nonumber \\ 
    & &&+ \frac{Q_1(\lambda_0)\sigma_0^2}{N-S}\declamb^{-\frac{3}{4}} \nonumber \\ 
    &\lec{(b)}{=} \norm{(f\circ\encs)^{(2)}}^2_{\lp{2}} \Tilde{\eta}\cdot\declamb^{\frac{1}{2}}\cdot p_4(S+1)^{\frac{1}{2}} &&+  \frac{Q_1(\lambda_0)\sigma_0^2}{N-S}\declamb^{-\frac{3}{4}},
\end{alignat}
% \begin{alignat}[2]\label{eq:bound_hp_2}
%     \mathbb{E}_{\epsilon, \stset}\left[\norm{h'}^2_{\lp{2}}\right] &\lec{}{\leqslant}H_1\norm{(f\circ\encs)^{(2)}}^2_{\lp{2}}\cdot\declamb^{\frac{1}{2}}\cdot p_4(S+1)^{\frac{1}{2}} &&\cdot\left(1 + \declamb\frac{p_4(S+1)}{16}\right)^{\frac{1}{2}} \nonumber \\ & &&+  \frac{Q_1\sigma_0^2}{N-S}\declamb^{-\frac{3}{4}} \nonumber \\ 
%     &\lec{(a)}{\leqslant} H_1\norm{(f\circ\encs)^{(2)}}^2_{\lp{2}}\cdot\declamb^{\frac{1}{2}}\cdot p_4(S+1)^{\frac{1}{2}} &&\cdot\left(1 + \lambda_0\frac{p_4(S+1)}{16}\right)^{\frac{1}{2}} \nonumber \\ & &&+ \frac{Q_1\sigma_0^2}{N-S}\declamb^{-\frac{3}{4}} 
%     \nonumber \\ &\lec{(b)}{=} \norm{(f\circ\encs)^{(2)}}^2_{\lp{2}} \Tilde{\eta}\cdot\declamb^{\frac{1}{2}}\cdot p_4(S+1)^{\frac{1}{2}} &&+  \frac{Q_1\sigma_0^2}{N-S}\declamb^{-\frac{3}{4}},
% \end{alignat}
where (a) follows from the definition $\declamb \leqslant \lambda_0$, (b) is derived from the definition $\Tilde{\eta} := H_1\left(1 + \lambda_0\frac{p_4(S+1)}{16}\right)^{\frac{1}{2}}$. By applying the upper bound for $\ldec$ from \eqref{eq:ldec_step3} and incorporating the results from \eqref{eq:bound_h_2} and \eqref{eq:bound_hp_2}, we can deduce the following:
\begin{alignat}{2}
    \ldec &\lec{(a)}{\leqslant} 2&&\left(\norm{(f\circ\encs)^{(2)}}^2_{\lp{2}}\cdot \mu_{0}(S)\declamb + \frac{Q_0(\lambda_0)\sigma_0^2}{N-S}\declamb^{-\frac{1}{4}} \right)^{\frac{1}{2}} \nonumber \\ & &&\cdot\left(\norm{(f\circ\encs)^{(2)}}^2_{\lp{2}}\cdot \mu_{1}(S)\declamb^{\frac{1}{2}} + \frac{Q_1(\lambda_0)\sigma_0^2}{N-S}\declamb^{-\frac{3}{4}} \right)^{\frac{1}{2}} \nonumber \\
    &\lec{(b)}{\leqslant} 2&& 
    \declamb^\frac{1}{4} \left(\norm{(f\circ\encs)^{(2)}}^2_{\lp{2}}\cdot\mu_{\textrm{max}}(S)\declamb^{\frac{1}{2}} +\frac{Q_{\textrm{max}}(\lambda_0)\sigma_0^2}{N-S}\declamb^{-\frac{3}{4}}\right) \nonumber \\
    &\lec{}{=} 2 &&\declamb^{\frac{3}{4}}\cdot \norm{(f\circ\encs)^{(2)}}^2_{\lp{2}}\cdot\mu_{\textrm{max}}(S) +\frac{Q_{\textrm{max}}(\lambda_0)\sigma_0^2}{N-S}\declamb^{-\frac{1}{2}},
\end{alignat}
where (a) follows by the definitions $\mu_0(S) := H_0\cdot p_4(S+1)$ and $\mu_1(S) := \Tilde{\eta}\cdot p_4(S+1)^{\frac{1}{2}}$, and (b) is due to $Q_0(\lambda_0), Q_1(\lambda_0) \leqslant Q_{\textrm{max}}(\lambda_0):=\max\{Q_0(\lambda_0), Q_1(\lambda_0)\}$ and $\mu_0(S), \mu_1(S) \leqslant \mu_{\textrm{max}}(S):=\max\{\mu_0(S), \mu_1(S)\}$. Therefore, we can conclude that:
\begin{align}\label{eq:dec_upper_final}
    \ldec \leqslant 2\declamb^{\frac{3}{4}}\cdot \mu_{\textrm{max}}(S) \cdot \norm{(f\circ\encs)^{(2)}}^2_{\lp{2}} + \declamb^{-\frac{1}{2}}\cdot\frac{Q_{\textrm{max}}(\lambda_0)\sigma_0^2}{N-S}.
\end{align}
Based on the definition of $\mu_{\textrm{max}}(S)$ we have:
\begin{align}
    \mu_{\textrm{max}}(S) &=\max\left\{H_0\cdot p_4(S), H_1\left(1 + \lambda_0\frac{p_4(S)}{16}\right)^{\frac{1}{2}}p_4(S)^\frac{1}{2}\right\} \nonumber \\ 
    &\leqslant H_{\textrm{max}}\cdot p_4(S)^\frac{1}{2} \cdot \max \left\{ p_4(S), 1 + \frac{\lambda_0p_4(S)}{16}\right\}^\frac{1}{2}, \nonumber \\ 
    &\leqslant H_{\textrm{max}}\cdot p_4(S)^\frac{1}{2} \cdot \left(\frac{1+p_4(S)}{16}\right)^\frac{1}{2} \max \left\{ 16, \lambda_0\right\}^\frac{1}{2}, \nonumber \\ 
    &= H_{\textrm{max}}\cdot \widetilde{p_4}(S)\max \left\{4, \lambda_0\right\}^\frac{1}{2}, 
\end{align}
where $H_{\textrm{max}}:=\max \{H_0, H_1\}$ and $\widetilde{p_4}(S):=p_4(S)^\frac{1}{2} \cdot \left(\frac{1+p_4(S)}{16}\right)^\frac{1}{2}$. Based on definition of $Q_{\textrm{max}}(\lambda_0)$ (mentioned in Theorem~\ref{th:lit_spline_noisy}), we have:
\begin{align}
    Q_{\textrm{max}}(\lambda_0) &= \max \{w_0\lambda_0^\frac{1}{4} + \Tilde{w}_0, w_1\lambda_0^\frac{1}{4} + \Tilde{w}_1\} \nonumber \\
    &\leqslant w_{\textrm{max}}\lambda_0^\frac{1}{4} + \Tilde{w}_{\textrm{max}},
    \nonumber \\
    &\leqslant 2w_{\textrm{max}}\max \left\{\lambda_0, \left(\frac{\Tilde{w}_{\textrm{max}}}{w_{\textrm{max}}}\right)^4 \right\}^\frac{1}{4}
\end{align}
where $w_{\textrm{max}}:=\max \{w_0, w_1\}$ and $\Tilde{w}_{\textrm{max}}:= \max \{\Tilde{w}_0, \Tilde{w}_1\}$. Therefore we have:
\begin{align}
    \ldec \leqslant 2\declamb^{\frac{3}{4}}\cdot H_{\textrm{max}}\cdot \widetilde{p_4}(S)\max \left\{4, \lambda_0\right\}^\frac{1}{2} \cdot \norm{(f\circ\encs)^{(2)}}^2_{\lp{2}} + \declamb^{-\frac{1}{2}}\frac{2\sigma_0^2 w_{\textrm{max}}\max\left\{\lambda_0^\frac{1}{4}, \frac{\Tilde{w}_{\textrm{max}}}{w_{\textrm{max}}} \right\}}{N-S}.
\end{align}
Since \eqref{eq:dec_upper_final} holds for all $\declamb$, by minimizing the right-hand side of \eqref{eq:dec_upper_final} with respect to $\declamb$, we obtain:
\begin{align}
    \declamb^* = \left(\frac{3H_{\textrm{max}}\cdot \widetilde{p_4}(S) \cdot (N-S)\cdot\norm{(f\circ\encs)^{(2)}}^2_{\lp{2}}}{4w_{\textrm{max}}\sigma_0^2} \right)^{-\frac{4}{5}} \left(\frac{\max \left\{4, \lambda_0\right\}^\frac{1}{2}}{\max\left\{\lambda_0^\frac{1}{4}, \frac{\Tilde{w}_{\textrm{max}}}{w_{\textrm{max}}} \right\}}\right)^{-\frac{4}{5}}
\end{align}
By substituting the expression for $\declamb^*$ into \eqref{eq:dec_upper_final}, we have:
\begin{align}
    \ldec \leqslant 4\left(\frac{3H_{\textrm{max}}}{4w_{\textrm{max}}}\right)^{\frac{2}{5}} \left(\frac{\sigma_0}{N-S}\right)^\frac{3}{5} \cdot \widetilde{p_4}(S)^\frac{2}{5}\cdot \norm{(f\circ\encs)^{(2)}}^\frac{4}{5}_{\lp{2}} \left(\frac{\max \left\{4, \lambda_0\right\}^\frac{1}{2}}{\max\left\{\lambda_0^\frac{1}{4}, \frac{\Tilde{w}_{\textrm{max}}}{w_{\textrm{max}}} \right\}}\right)^{\frac{2}{5}}.
\end{align}
Thus, defining $C_2 := \frac{3H_{\textrm{max}}}{4w_{\textrm{max}}}, C(\lambda_0):=\left(\frac{\max \left\{4, \lambda_0\right\}^\frac{1}{2}}{\max\left\{\lambda_0^\frac{1}{4}, \frac{\Tilde{w}_{\textrm{max}}}{w_{\textrm{max}}} \right\}}\right)$, and using previously driven upper bound for $\lenc$ in Lemma~\ref{lem:lt}  completes the proof.
% \subsection{Proof of Theorem~\ref{th:conv_rate_noisy}}\label{sec:proof_noisy}
% As previously mentioned, $\decs(\cdot)$ is a second-order smoothing spline fitted on
% $\left\{(\beta_{i_1},f_{\epsilon_{i_1}}(\encs(\beta_{i_1})),\dots,(\beta_{i_{|\stset|}},f_{\epsilon_{i_{|\stset|}}}(\encs(\beta_{i_{|\stset|}}))\right\}$, where $\stset := \left\{\beta_{i_1},\dots,\beta_{i_{|\stset|}} \right\}$ is a set of non-straggler worker nodes and $\bm{f\circ\encs} := \left\{f_{\epsilon_{i_1}}(\encs(\beta_{i_1})),\dots,f_{\epsilon_{i_{|\stset|}}}(\encs(\beta_{i_{|\stset|}}))\right\}$ is a set of noisy computation values form non-straggler workers. Let us define $\splineCC(\cdot)$ representing the smoothing spline operator of the decoder. Using the decomposition \eqref{eq:smooth_spline_linear_eq} we have:
% \begin{alignat}{2}
%     \mathbb{E}_{\epsilon, \stset}\left[\norm{h}^2_{\lp{2}}\right] &= \mathbb{E}_{\stset}&&\left[\norm{f\circ\encs - \splineCC(\mathbf{f})}^2_{\lp{2}}\right] \nonumber \\ & &&+  \mathbb{E}_{\epsilon, \stset}\left[\norm{f\circ\encs - \splineCC(\bm{\epsilon})}^2_{\lp{2}}\right],
% \end{alignat}
% where $\mathbf{f}=\left\{f(\encs(\beta_{i_1})),\dots,f(\encs(\beta_{i_{|\stset|}}))\right\}$ and $\bm{\epsilon} = \left\{\epsilon_{i_1},\dots, \epsilon_{i_{|\stset|}}\right\}$. Let us define the following variables analogous to those in \eqref{eq:max_dist_st}:
% \begin{align}
% \Delta^\stset_\textrm{max}:=\underset{f \in \left\{0,\dots,|\stset|\right\}}{\max} \left\{\beta_{i_{f+1}}-\beta_{i_f}\right\}, \quad \Delta^\stset_\textrm{min}:=\underset{f \in \left\{1,\dots,|\stset|-1\right\}}{\min} \left\{\beta_{i_{f+1}}-\beta_{i_f}\right\}.
% \end{align}
% Since there are at most $S$ stragglers among worker nodes, $\Delta^\stset_\textrm{max} \leqslant (S+1) \cdot \Delta_\textrm{max}$ and $ \frac{\Delta^\stset_\textrm{max}}{\Delta^\stset_\textrm{min}} \leqslant (S+1)\cdot\frac{\Delta_\textrm{max}}
% {\Delta_\textrm{min}} \leqslant (S+1)B$. Therefore, $\frac{\Delta^\stset_\textrm{max}}{\Delta^\stset_\textrm{min}}$ is bounded. Additionally, since $\Delta_\textrm{min} \leqslant \frac{2}{N}$ and $\Delta_\textrm{max} \geqslant \frac{2}{N}$, then $\frac{\Delta_\textrm{max}}{\Delta_\textrm{min}} \leqslant B$ implies that both $\Delta_\textrm{max}$ and $\Delta_\textrm{min}$ has the same order of $\mathcal{O}(\frac{1}{N})$. Thus, there exists a constant $J$ such that $\Delta_\textrm{max} \leqslant \frac{J}{N}$. Thus, we have $\Delta^\stset_\textrm{max} \leqslant \Delta_\textrm{max}\cdot (S+1) \leqslant \frac{J(S+1)}{N} \leqslant \frac{J(S+1)}{N-S}$.
% Applying Theorems~\ref{th:lit_spline_noisy} and \ref{th:lit_spline_noiseless} with $\Omega=(-1, 1), m=2$, we have
% \begin{align}\label{eq:h_1}
%     \mathbb{E}_{\epsilon, \stset}\left[\norm{h}^2_{\lp{2}}\right] &\lec{}{\leqslant}H_0\norm{(f\circ\encs)^{(2)}}^2_{\lp{2}}\cdot\mathbb{E}_{\stset}\left[L\right] + \frac{Q_0\sigma_0^2}{N-S}\declamb^{-\frac{1}{4}},
% \end{align}
% and 
% \begin{align}\label{eq:hp_1}
%     \mathbb{E}_{\epsilon, \stset}\left[\norm{h'}^2_{\lp{2}}\right] &\lec{}{\leqslant}H_1\norm{(f\circ\encs)^{(2)}}^2_{\lp{2}}\cdot\mathbb{E}_{\stset}\left[L^{\frac{1}{2}}(1 + \frac{L}{16})^{\frac{1}{2}}\right] + \frac{Q_1\sigma_0^2}{N-S}\declamb^{-\frac{3}{4}},
% \end{align}
% where $L=p_2\left(\frac{\Delta^\stset_\textrm{max}}{\Delta^\stset_\textrm{min}}\right)\cdot\frac{(N-S) \Delta^\stset_\textrm{max}}{4} \declamb + D(2)\cdot \left({\Delta^\stset_\textrm{max}}\right)^{4}$, $H_0, H_1 := H(2,0), H(2,1)$, and $Q_0, Q_1:=Q(2,0), Q(2,1)$ as defined in  Theorems~\ref{th:lit_spline_noiseless} and ~\ref{th:lit_spline_noisy}. 
% Since $\frac{\Delta^\stset_\textrm{max}}{\Delta^\stset_\textrm{min}} \leqslant (S+1)B$, we have:
% \begin{align}\label{eq:bound_L}
%     \mathbb{E}_{\stset}[L] &\leqslant \mathbb{E}_{\stset}\left[p_2(S+1)\cdot\frac{(N-S) \Delta^\stset_\textrm{max}}{4} \declamb + D\cdot \left({\Delta^\stset_\textrm{max}}\right)^{4} \right] \nonumber \\ &\lec{(a)}{\leqslant} p_2(S+1)\cdot\frac{J(S+1)}{4} \declamb + DJ^4\frac{(S+1)^4}{(N-S)^4} \nonumber \\ &\lec{}{=} p_3(S+1)\cdot\declamb + DJ^4\frac{(S+1)^4}{(N-S)^4},
% \end{align}
% where $D:=D(2)$, $p_3(S+1) := p_2(S+1)\cdot\frac{J(S+1)}{4}$ is a degree three polynomial of $(S+1)$, and (a) follows from the fact that $\Delta^\stset_\textrm{max} \leqslant \frac{J(S+1)}{N-S}$. Substituting in \eqref{eq:h_1}, we have:
% \begin{align}\label{eq:bound_h_2}
%     \mathbb{E}_{\epsilon, \stset}\left[\norm{h}^2_{\lp{2}}\right] &\lec{}{\leqslant}H_0\norm{(f\circ\encs)^{(2)}}^2_{\lp{2}} 
%     \left(p_3(S+1)\cdot\declamb + DJ^4\frac{(S+1)^4}{(N-S)^4}\right) + \frac{Q_0\sigma_0^2}{N-S}\declamb^{-\frac{1}{4}} \nonumber \\ 
%     &\lec{(a)}{\leqslant} H_0\norm{(f\circ\encs)^{(2)}}^2_{\lp{2}}\declamb\cdot 
%     \left(p_3(S+1) + DJ^4(S+1)^4\right) + \frac{Q_0\sigma_0^2}{N-S}\declamb^{-\frac{1}{4}} \nonumber \\ &\lec{(b)}{=} H_0\norm{(f\circ\encs)^{(2)}}^2_{\lp{2}}\declamb\cdot 
%     p_4(S+1) + \frac{Q_0\sigma_0^2}{N-S}\declamb^{-\frac{1}{4}},
% \end{align}
% where (a) follows from the fact that $\declamb^{-1} (N-S)^{-4} \leqslant 1$, as assumed in Theorem~\ref{th:lit_spline_noisy} and (b) is by definition $p_4(S+1) := p_3(S+1) + DJ^4(S+1)^4$ is a degree four polynomial of $(S+1)$. An analogous upper bound can be derived for $\mathbb{E}_{\epsilon, \stset}\left[\norm{h'}^2_{\lp{2}}\right]$ as
% \begin{align}\label{eq:bound_hp_2}
%     \mathbb{E}_{\epsilon, \stset}\left[\norm{h'}^2_{\lp{2}}\right] &\lec{}{\leqslant}H_1\cdot\declamb^{\frac{1}{2}}\cdot p_4(S+1)^{\frac{1}{2}}\cdot\left(1 + \declamb\frac{p_4(S+1)}{16}\right)^{\frac{1}{2}} +  \frac{Q_1\sigma_0^2}{N-S}\declamb^{-\frac{3}{4}} \nonumber \\ 
%     &\lec{(a)}{\leqslant} H_1\cdot\declamb^{\frac{1}{2}}\cdot p_4(S+1)^{\frac{1}{2}}\cdot\left(1 + \lambda_0\frac{p_4(S+1)}{16}\right)^{\frac{1}{2}} +  \frac{Q_1\sigma_0^2}{N-S}\declamb^{-\frac{3}{4}} 
%     \nonumber \\ &\lec{(b)}{=} \Bar{H_1}\cdot\declamb^{\frac{1}{2}}\cdot p_4(S+1)^{\frac{1}{2}} +  \frac{Q_1\sigma_0^2}{N-S}\declamb^{-\frac{3}{4}},
% \end{align}
% where (a) is by definition $\declamb \leqslant \lambda_0$ and (b) is by definition $\Bar{H_1} := H_1\left(1 + \lambda_0\frac{p_4(S+1)}{16}\right)^{\frac{1}{2}}$. Using the last upper bound for $\ldec$ in \eqref{eq:ldec_step3_2} and combining the results from \eqref{eq:bound_h_2} and \eqref{eq:bound_hp_2}, we can conclude that:
% \begin{alignat}{2}
%     \ldec &\lec{(a)}{\leqslant} \frac{(K+1)}{K} &&\mu_{0}(S)\declamb + \frac{Q_0\sigma_0^2}{N-S}\declamb^{-\frac{1}{4}} \nonumber \\& &&+ \frac{8C_0}{K+1}\cdot\left(\mu_{0}(S)\declamb + \frac{Q_0\sigma_0^2}{N-S}\declamb^{-\frac{1}{4}} \right)^{\frac{1}{2}}\cdot\left(\mu_{1}(S)\declamb^{\frac{1}{2}} + \frac{Q_1\sigma_0^2}{N-S}\declamb^{-\frac{3}{4}} \right)^{\frac{1}{2}} \nonumber \\
%     &\lec{(b)}{\leqslant} \frac{(K+1)}{K} &&\mu_{0}(S)\declamb + \frac{Q_0\sigma_0^2}{N-S}\declamb^{-\frac{1}{4}} + \frac{8C_0}{K+1}\declamb^{\frac{1}{4}}\left(\mu_{\textrm{max}}(S)\declamb^{\frac{1}{2}} +\frac{Q_{\textrm{max}}\sigma_0^2}{N-S}\declamb^{-\frac{3}{4}} \right) \nonumber \\
%     &\lec{}{=} \declamb^{\frac{3}{4}}\cdot&&\left(\frac{(K+1)}{K} \mu_{0}(S)\declamb^{\frac{1}{4}} + \frac{8C_0}{K+1}\mu_{\textrm{max}}(S)\right) + \declamb^{-\frac{1}{2}}\frac{Q_{\textrm{max}}\sigma_0^2}{N-S}\left(\frac{Q_0}{Q_{\textrm{max}}} \declamb^{\frac{1}{4}} + \frac{8C_0}{K+1}\right) \nonumber \\
%     &\lec{(c)}{\leqslant} \declamb^{\frac{3}{4}}\cdot&&\left(\frac{(K+1)}{K} \mu_{0}(S)\lambda_0^{\frac{1}{4}} + \frac{8C_0}{K+1}\mu_{\textrm{max}}(S)\right) + \declamb^{-\frac{1}{2}}\frac{Q_{\textrm{max}}\sigma_0^2}{N-S}\left(\frac{Q_0}{Q_{\textrm{max}}} \lambda_0^{\frac{1}{4}} + \frac{8C_0}{K+1}\right),
% \end{alignat}
% where (a) follows by the definitions $\mu_0(S) := H_0\norm{(f\circ\encs)^{(2)}}^2_{\lp{2}}\cdot p_4(S+1)$ and $\mu_1(S) := \Bar{H_1}\norm{(f\circ\encs)^{(2)}}^2_{\lp{2}}\cdot p_4(S+1)^{\frac{1}{2}}$, (b) is due to $Q_0, Q_1 \leqslant Q_{\textrm{max}}:=\max\{Q_0, Q_1\}$ and $\mu_0(S), \mu_1(S) \leqslant \mu_{\textrm{max}}(S):=\max\{\mu_0(S), \mu_1(S)\}$, and (c) follow from $\declamb \leqslant \lambda_0$. By defining $\widetilde{p_4}(S+1):=\frac{\left(\frac{(K+1)}{K} \mu_{0}(S)\lambda_0^{\frac{1}{4}} + \frac{8C_0}{K+1}\mu_{\textrm{max}}(S)\right)}{\norm{(f\circ\encs)^{(2)}}^2_{\lp{2}}}$ which is a polynomial of $(S+1)$ with maximum degree of $4$ and the coefficients depend on $(H_0, H_1, B, K)$ as well as $\widetilde{Q} = \left(\frac{Q_0}{Q_{\textrm{max}}} \lambda_0^{\frac{1}{4}} + \frac{8C_0}{K+1}\right)$, we can conclude that:
% \begin{align}\label{eq:dec_upper_final}
%     \ldec \leqslant \declamb^{\frac{3}{4}}\cdot \widetilde{p_4}(S+1) \cdot \norm{(f\circ\encs)^{(2)}}^2_{\lp{2}} + \declamb^{-\frac{1}{2}}\cdot\frac{\widetilde{Q} Q_{\textrm{max}}\sigma_0^2}{N-S}.
% \end{align}
% Optimizing over $\declamb$, the optimum $\declamb^*$ is:
% \begin{align}
%     \declamb^* = \left(\frac{\widetilde{p_4}(S+1) \cdot (N-S)\cdot\norm{(f\circ\encs)^{(2)}}^2_{\lp{2}}}{\widetilde{Q}Q_{\textrm{max}}\sigma_0^2} \right)^{-\frac{4}{5}}.
% \end{align}
% By substituting the expression for $\declamb^*$ into \eqref{eq:dec_upper_final}, defining $C_1 := 2\left(\widetilde{Q}Q_{\textrm{max}}\right)^{\frac{5}{3}}$  and using previously driven upper bound for $\lenc$ in Lemma~\ref{lem:lt}, the proof is completed.
% \subsection{Proof of Theorem~\ref{th:conv_rate_noiseless}}\label{app:proof_noiseless}
% In the noiseless setting, the expectation is just over $\stset$:
% \begin{align}
%     \mathbb{E}_{\stset}\left[\norm{h}^2_{\lp{2}}\right] = \mathbb{E}_{\stset}\left[\norm{f\circ\encs - \splineCC(\mathbf{f})}^2_{\lp{2}}\right].
% \end{align}
% Using the same steps as Appendix~\ref{app:theorem_proofs}, we have:
% \begin{alignat}{2}
%     \label{eq:h_noiseless_bound_1}
%     \mathbb{E}_{\stset}\left[\norm{h}^2_{\lp{2}}\right] &\leqslant H_0\norm{(f\circ\encs)^{(2)}}^2_{\lp{2}} 
%     &&\left(p_3(S+1)\cdot\declamb + DJ^4\frac{(S+1)^4}{(N-S)^4}\right), \\ \label{eq:hp_noiseless_bound_1}
%     \mathbb{E}_{\stset}\left[\norm{h'}^2_{\lp{2}}\right] &\leqslant H_1\norm{(f\circ\encs)^{(2)}}^2_{\lp{2}} 
%     &&\left(p_3(S+1)\cdot\declamb + DJ^4\frac{(S+1)^4}{(N-S)^4}\right)^{\frac{1}{2}} \nonumber \\ & &&\cdot\left(1 + \frac{p_3(S+1)\cdot\declamb + DJ^4\frac{(S+1)^4}{(N-S)^4}}{16}\right)^{\frac{1}{2}}.
% \end{alignat}
% Note that both \eqref{eq:h_noiseless_bound_1} and \eqref{eq:hp_noiseless_bound_1} are increasing function of $\declamb$ and the upper bound of $\ldec$ in \eqref{eq:ldec_step3_2} is increasing function of both $\mathbb{E}_{\stset}\left[\norm{h}^2_{\lp{2}}\right]$ and $\mathbb{E}_{\stset}\left[\norm{h'}^2_{\lp{2}}\right]$. Therefore, $\declamb^* = 0$ minimizes the right-hand side of \eqref{eq:ldec_step3_2} in the noiseless setting. Putting $\declamb =0$, we have:
% \begin{alignat}{2}
%     \ldec &\leqslant \norm{(f\circ\encs)^{(2)}}^2_{\lp{2}} \Biggl[\frac{(K+1)}{K} &&H_0\cdot r(S, N) + \frac{8C_0}{K+1} \cdot \left(H_0\cdot r(S, N)\right)^{\frac{1}{2}} \cdot \nonumber \\ & 
%     &&\left(H_1\cdot r(S, N)^{\frac{1}{2}}\cdot \left(1 + \frac{r(S, N)}{16}\right)^{\frac{1}{2}}\right)^{\frac{1}{2}}
%     \Biggr] \nonumber \\
%     &\lec{(a)}{\leqslant} \norm{(f\circ\encs)^{(2)}}^2_{\lp{2}} \cdot\Biggl[\frac{(K+1)}{K} && H_0\cdot r(S, N) + \frac{8C_0 \eta D^{\frac{1}{4}}J\cdot(S+1)}{K+1} \cdot H_0^{\frac{1}{2}}H_1^{\frac{1}{2}}\cdot r(S, N)^{\frac{3}{4}}\Biggr]
% \end{alignat}
% where $r(S, N):=DJ^4\frac{(S+1)^4}{(N-S)^4}$ and (a) follows from $1 + \frac{r(S, N)}{16} \leqslant \eta\cdot DJ^4(S+1)^4$ for some $\eta > 0$. Therefore, defining $C_3:=\frac{(K+1)}{K} H_0$ and $C_2:=\frac{8C_0 \eta D^{\frac{1}{4}}J}{K+1} H_0^{\frac{1}{2}}H_1^{\frac{1}{2}}$ as well as using Lemma~\ref{lem:lt} complete the proof.

\subsection{Proof of  Theorem~\ref{th:optimcal_encdoer}}\label{app:proof_optimal_enc}
The upper bounds presented in Theorems~\ref{th:conv_rate_noiseless} and \ref{th:conv_rate_noisy} depend on $\norm{(f\circ\encs)^{(2)}}^2_{\lp{2}}$, with exponents of $\frac{2}{5}$ and $1$, respectively. By applying the chain rule, we can demonstrate that:
\begin{align}\label{eq:optimal_encoder_1}
    \int_\Omega\left[f(\encs(t))''\right]^2 \,dt &\lec{(a)}{=} \int_\Omega \left[ \encs''(t)\cdot f'(\encs(t)) + \encs'(t)^2 \cdot f''(\encs(t))\right]^2\,dt \nonumber \\ 
    &\lec{(b)}{\leqslant} \int_\Omega \left[\encs''(t)^2 + \encs'(t)^4\right] \left[f'(\encs(t))^2 + f''(\encs(t))^2\right]\,dt \nonumber \\
    &\lec{(c)}{\leqslant} (q^2 +\nu^2) \int_\Omega \left[\encs''(t)^2 + \encs'(t)^4\right]\,dt \nonumber\\ 
    &\lec{}{=} (q^2 + \nu^2) \left(\norm{\encs''(t)}^2_{\lp{2}} + \norm{\encs'(t)}^4_{\lp{4}} \right) \nonumber\\
    &\lec{(d)}{\leqslant} (q^2 + \nu^2) \left(\norm{\encs''(t)}^2_{\lp{2}} + 
    \left(\norm{\encs'(t)}_{\lp{2}} + \norm{\encs''(t)}_{\lp{2}}\right)^4\right) \nonumber \\ 
    &\lec{(e)}{\leqslant} (q^2 + \nu^2) \left(\norm{\encs''(t)}^2_{\lp{2}} + 
    4\left(\norm{\encs'(t)}^2_{\lp{2}} + \norm{\encs''(t)}^2_{\lp{2}}\right)^2\right)
    \nonumber \\ 
    &\lec{(f)}{\leqslant} (q^2 + \nu^2) \left(\norm{\encs}^2_{\sob{2}{2}} + 
    4\norm{\encs}^4_{\sob{2}{2}}\right)
    \nonumber \\ 
    &\lec{(g)}{\leqslant} (q^2 + \nu^2) \left(73\norm{\encs}^2_{\sobeq{2}{2}} + 
    4\times73^2\norm{\encs}^4_{\sobeq{2}{2}}\right)
    \nonumber \\ 
    &\lec{(h)}{=} (q^2 + \nu^2) \cdot \psi\left(\norm{\encs}^2_{\sobeq{2}{2}}\right),
    \nonumber \\
    &\lec{(i)}{=} (q^2 + \nu^2) \cdot \psi \left(\norm{\encs}^2_{\hiltilde{2}}\right),
\end{align}
where (a) follows from the chain rule, (b) is derived using the Cauchy-Schwartz inequality, (c) is due to $\norm{f_0''}_{\lp{\infty}} \leqslant \nu$, (d) follows from Theorem~\ref{th:interp_ineq} with $r=4, p=q=2, l=1$, (e) follows from AM-GM inequality, (f) is a result of adding positive terms $\norm{\encs}^2_{\sob{2}{2}}$ and $\norm{\encs'}^2_{\sob{2}{2}}$ to the first term and $\norm{\encs}^2_{\sob{2}{2}}$ to the second term in the parenthesis, (g) is result of applying Corollary~\ref{cor:norm-equiv-1}, (h) is by defining $\psi(t):=73t + 4\times73^2t^2$, and (i) is because of Proposition~\ref{prop:sob_hilb}.

Combining \eqref{eq:optimal_encoder_1} with Theorems~\ref{th:conv_rate_noiseless} and \ref{th:conv_rate_noisy} we have
\begin{align}\label{eq:encoder_rep_noiseless}
    \mathcal{R}(\hat{f}) \leqslant \frac{2q^2}{K} \sum^K_{k=1} (\encs(\alpha_k) - x_k)^2 + \enclamb \cdot \psi\left(\norm{\encs}^2_{\hiltilde{2}}\right),
\end{align}
for the noiseless setting and
\begin{align}\label{eq:encoder_rep_noisy}
    \mathcal{R}(\hat{f}) \leqslant \frac{2q^2}{K} \sum^K_{k=1} (\encs(\alpha_k) - x_k)^2 + \widetilde{\enclamb} \cdot \psi\left(\norm{\encs}^2_{\hiltilde{2}}\right)^{\frac{2}{5}}
\end{align}
for the noisy setting, where the parameters $\enclamb$ and $\widetilde{\enclamb}$ are as follows:
\begin{align}\label{eq:optim_lambda_rate}
    \enclamb &= C_1\frac{(S+1)^3}{N^3}\cdot (q^2+\nu^2)
    \\
    \widetilde{\enclamb} &= 2C(\lambda_0)^\frac{3}{4} \left(\frac{\sigma^2_0}{N-S}\right)^{\frac{3}{5}} \cdot p_4(S)^{\frac{2}{5}} \cdot (q^2+\nu^2)^{\frac{2}{5}}.
\end{align}
Note that since $\psi(t)$ and $\gamma(t) := t^\frac{2}{5}$ are monotonically increasing in $\mathbb{R}^+$, its composition is monotonically increasing as well. Moreover, $\enclamb$ and $\widetilde{\enclamb}$  share the same exponent of $N$ as in Theorems~\ref{th:conv_rate_noiseless} and \ref{th:conv_rate_noisy}, respectively. Consequently, the provided upper bound does not compromise the convergence rate.

\subsection{Proof of Proposition~\ref{prop:encoder_optimal}}\label{app:proof_cor_optimal_enc}
For part (i), we know that for $t \leqslant M$, we have:
\begin{align}
    \psi(t)=73t + 4\times73^2t^2\leqslant t(73 + 4\times73^2t) \leqslant t(73 + 4\times73^2\cdot M) = t(m_1 + m_2M),
\end{align}
where $m_1:=73, m_2:=4\times73^2$. Therefore, if $\norm{\encs}^2_{\hiltilde{2}} \leqslant M$, then:
\begin{align}
    \psi(\norm{\encs}^2_{\hiltilde{2}}) &\leqslant (m_1 + m_2M)\norm{\encs}^2_{\hiltilde{2}} \nonumber \\
    &= (m_1 + m_2M)\left(\encs(-1)^2 + \encs'(-1)^2 + \int_\Omega |\encs''(t)|^2\,dt\right)
    \nonumber \\
    &\lec{(a)}{\leqslant} (m_1 + m_2M)\left(M + \int_\Omega |\encs''(t)|^2\,dt\right),
\end{align}
where (a) is because of $\norm{\encs}^2_{\hiltilde{2}} \leqslant M$. Thus, we have:
\begin{align}
    \mathcal{R}(\hat{f}) &\leqslant \frac{2q^2}{K} \sum^K_{k=1} (\encs(\alpha_k) - x_k)^2 + \enclamb \cdot \psi\left(\norm{\encs}^2_{\hiltilde{2}}\right) \nonumber \\
    &\leqslant \frac{2q^2}{K} \sum^K_{k=1} (\encs(\alpha_k) - x_k)^2 + \enclamb \cdot (m_1 + m_2M)\left(M + \int_\Omega |\encs''(t)|^2\,dt\right) \nonumber \\
    &\leqslant \widetilde{R}(\encs).
\end{align}
% it is sufficient to show that $\psi\big(\norm{\encs}^2_{\hiltilde{2}}\big) \leqslant T \left(8M^2 + \int_\Omega u''(t)^2\,dt\right)$. From the definition of $\norm{\cdot}_{\hiltilde{2}}$ we have:
% \begin{alignat}{2}\label{eq:enc_bounded_norm}
%     \norm{\encs}^2_{\hiltilde{2}} &= &&\encs(-1)^2 + \left[\encs'(-1)\right]^2 + \int_\Omega \left[\encs''(t)\right]^2\,dt \nonumber \\
%     &\lec{}{\leqslant} &&\norm{\encs}^2_{\lp{\infty}} + \norm{\encs'}^2_{\lp{\infty}} + \norm{\encs''}^2_{\lp{2}} \nonumber \\
%     &\lec{(a)}{\leqslant} &&\left(\norm{\encs}_{\lp{2}} + \norm{\encs'}_{\lp{2}}\right)^2 + \left(\norm{\encs'}_{\lp{2}} + \norm{\encs''}_{\lp{2}}\right)^2 \nonumber \\ & &&+ \norm{\encs''}^2_{\lp{2}} 
%     \nonumber \\ 
%     &\leqslant &&8M^2 +  \norm{\encs''}^2_{\lp{2}} \leqslant 9M^2,
% \end{alignat}
% where (a) follows from Theorem~\ref{th:interp_ineq} with $r=\infty, p=2, q=2$ and $\ell= 1$. From the definition of $\psi(t)$ for $t \leqslant 9M^2$ we have:
% \begin{align}
%     \psi(t):=73t + 4\times73^2t^2\leqslant t(73 + 4\times73^2t) \leqslant t(73 + 4\times73^2\cdot9M^2) = T\cdot t.
% \end{align}
% Therefore, based on \eqref{eq:enc_bounded_norm} we have
% \begin{align}\label{eq:optim_enc_final_noiseless}
%     \psi\left(\norm{\encs}^2_{\hiltilde{2}}\right) \leqslant T\cdot \norm{\encs}^2_{\hiltilde{2}} \leqslant T\cdot\left(8M^2 +  \norm{\encs''}^2_{\lp{2}}\right).
% \end{align}
% Combining the results from 
% ~\eqref{eq:enc_bounded_norm}, ~\eqref{eq:optim_enc_final_noiseless}, and \eqref{eq:encoder_rep_noiseless} we complete the proof.

For part (ii), let $\widetilde{\encs}(t)$ be a natural spline used as the encoder function fitted to the data points $\{(\alpha_k, x_k)\}_{k=1}^K$. Then, we have:
\begin{align}
    \enclamb (m_1 + m_2M)\left(M + \int_\Omega |(u^*)''(t)|^2\,dt\right) &\lec{}{\leqslant} \widetilde{R}(u^*) \nonumber \\ &\lec{(a)}{\leqslant} \widetilde{R}(\encs) \lec{(b)}{=} \enclamb (m_1 + m_2M)\left(M + \int_\Omega |\widetilde{\encs}''(t)|^2\,dt\right),
\end{align}
where (a) follows from the optimality of $u^*(\cdot)$, and (b) follows from $\widetilde{\encs}(\alpha_k) = x_k$ for $k \in [K]$. Therefore, $\int_\Omega |(u^*)''(t)|^2\,dt \leqslant \int_\Omega |\widetilde{\encs}''(t)|^2\,dt$. 

Since $u^*(\cdot)$ is smoothing spline, it has the representation in natural spline space (as mentioned in \eqref{eq:spline_bspline_1}):
\begin{align}
    u^*(t) = \sum^{K+4}_{k=1} \xi_kb_k(t),
\end{align}
where, $\{b_k(\cdot)\}^K_{k=1}$ is a basis functions of second order natural splines. Therefore, using Cauchy-Schwartz inequality, we have:
\begin{align} \label{eq:spline_l_inf}
    |u^*(t)|^2 \leqslant \left(\sum^{K+4}_{k=1} \xi^2 \right) \left(\sum^{K+4}_{k=1}|b_k(t)|^2\right),
\end{align}
and
\begin{align}\label{eq:spline_l_inf_p}
    |(u^*(t))'|^2 \leqslant \left(\sum^{K+4}_{k=1} \xi^2 \right) \left(\sum^{K+4}_{k=1}|b'_k(t)|^2\right).
\end{align}
Both \eqref{eq:spline_l_inf} and \eqref{eq:spline_l_inf_p} hold for all $t \in \Omega$. Thus:
\begin{align}\label{eq:spline_l_inf_2}
    |(u^*)'(-1)|^2 \leqslant \norm{u^*}^2_{\lp{\infty}} \leqslant \norm{\boldsymbol{\xi}}^2_2 \cdot \left(\sum^{K+4}_{k=1}\norm{b_k}^2_{\lp{\infty}}\right), \\ \label{eq:spline_l_inf_p_2}
    |(u^*)'(-1)|^2 \leqslant \norm{(u^*)'}^2_{\lp{\infty}} \leqslant \norm{\boldsymbol{\xi}}^2_2 \cdot \left(\sum^{K+4}_{k=1}\norm{b'_k}^2_{\lp{\infty}}\right),
\end{align}
where $\boldsymbol{\xi}:=[\xi_1,\dots,\xi_{K+4}]^T = \left(\mathbf{N}^T\mathbf{N} + \lambda\Phi\right)^{-1}\mathbf{N}^T\mathbf{x}$ with $\mathbf{N}_{ij} = b_i(\alpha_j), \Phi_{ij} = \int_\Omega b''_i(t)b''_j(t)\,dt$ for $i, j \in [{K+4}]$ as defined in \eqref{eq:spline_bspline_1}, and $\mathbf{x}:=[x_1,\dots,x_K]$. Noted that $\{\norm{b'_k}^2_{\lp{\infty}}\}_{k=1}^{K+4}$ and $\{\norm{b_k}^2_{\lp{\infty}}\}_{k=1}^{K+4}$ depend only on $\{\alpha_k\}_{k=1}^K$.
\begin{lemma}\label{lem:ridge_reg_weight_upperbound} If $\boldsymbol{\widetilde{\xi}}:=\left(\mathbf{N}^T\mathbf{N}\right)^{-1}\mathbf{N}^T\mathbf{x}$, then $\norm{\boldsymbol{\xi}}^2_2 \leqslant \kappa(\Phi)\cdot \norm{\boldsymbol{\widetilde{\xi}}}^2_2 <\frac{2}{|\operatorname{det} \Phi|}\left(\frac{\|\Phi\|_F^2}{K}\right)^{\frac{K}{2}} \norm{\boldsymbol{\widetilde{\xi}}}^2_2$, where $\kappa(\Phi)$ is condition number of $\Phi$.
\end{lemma}
\begin{proof}
    By defining $\widetilde{\mathbf{N}} := \mathbf{N}\Phi^{-\frac{1}{2}}$ and rearranging the expression for $\boldsymbol{\widetilde{\xi}}$, we obtain:
    \begin{align}\label{eq:ridge_w_1}
    \boldsymbol{\widetilde{\xi}} = \left(\mathbf{N}^T \mathbf{N}+\lambda \Phi\right)^{-1} \mathbf{N}^T \mathbf{x} &= \Phi^{-1 / 2}\left(\left[\mathbf{N} \Phi^{-1 / 2}\right]^T\left[\mathbf{N} \Phi^{-1 / 2}\right]+\lambda \mathbf{I}\right)^{-1}\left[\mathbf{N} \Phi^{-1 / 2}\right]^T \mathbf{x} \nonumber \\ &=\Phi^{-1 / 2}\left(\widetilde{\mathbf{N}}^T \widetilde{\mathbf{N}}+\lambda \mathbf{I}\right)^{-1} \widetilde{\mathbf{N}}^T \mathbf{x}.
    \end{align}
    Define $\mathbf{z}:=\left(\widetilde{\mathbf{N}}^T \widetilde{\mathbf{N}}\right)^{-1} \widetilde{\mathbf{N}}^T \mathbf{x}$. Thus, $\widetilde{\mathbf{N}}^T \mathbf{x} = \widetilde{\mathbf{N}}^T \widetilde{\mathbf{N}}\mathbf{z}$.
% Substituting in \eqref{eq:ridge_w_1}, $\boldsymbol{\widetilde{\xi}}= \Phi^{-1 / 2}\left(\widetilde{\mathbf{N}}^T \widetilde{\mathbf{N}}+\lambda I\right)^{-1}\widetilde{\mathbf{N}}^T \widetilde{\mathbf{N}}\mathbf{z}$.
Thus, by applying the Cauchy-Schwartz inequality, we have:
\begin{align}\label{eq:ridge_w_2}
    \norm{\left(\widetilde{\mathbf{N}}^T \widetilde{\mathbf{N}}+\lambda \mathbf{I}\right)^{-1}\widetilde{\mathbf{N}}^T \mathbf{x}}_2 =
    \norm{\left(\widetilde{\mathbf{N}}^T \widetilde{\mathbf{N}}+\lambda \mathbf{I}\right)^{-1}\widetilde{\mathbf{N}}^T \widetilde{\mathbf{N}}\mathbf{z}}_2 \leqslant \norm{\left(\widetilde{\mathbf{N}}^T \widetilde{\mathbf{N}}+\lambda \mathbf{I}\right)^{-1}\widetilde{\mathbf{N}}^T \widetilde{\mathbf{N}}}_2 \cdot \norm{\mathbf{z}}_2.
\end{align}
Let $\widetilde{\mathbf{N}} = \mathbf{UDV}^T$ be the singular value decomposition of $\widetilde{\mathbf{N}}$. Therefore, we have:
\begin{align}
    \norm{\left(\widetilde{\mathbf{N}}^T \widetilde{\mathbf{N}}+\lambda \mathbf{I}\right)^{-1} \widetilde{\mathbf{N}}^T \widetilde{\mathbf{N}}}_2 & =\norm{\left(\mathbf{V} \mathbf{D}^2 \mathbf{V}^T+\lambda \mathbf{I}\right)^{-1} \mathbf{V} \mathbf{D}^2 \mathbf{V}^T}_2 \nonumber \\ &= \norm{\mathbf{V}^{-T}\left(\mathbf{D}^2+\lambda \mathbf{I}\right)^{-1} \mathbf{V}^{-1} \mathbf{V} \mathbf{D}^2 \mathbf{V}^T}_2 \nonumber \\ &= \norm{\mathbf{V}^{-T}\left(\mathbf{D}^2+\lambda \mathbf{I}\right)^{-1} \mathbf{D}^2 \mathbf{V}^T}_2 
    \nonumber \\
    &\lec{(a)}{=} \norm{\left(\mathbf{D}^2+\lambda \mathbf{I}\right)^{-1} \mathbf{D}^2}_2
    \nonumber \\
    &\lec{}{=} \norm{\textrm{diag}\left(\frac{\lambda_1^2}{\lambda_1^2 + \lambda},\dots,\frac{\lambda_K^2}{\lambda_K^2 + \lambda}\right)}_2
    \nonumber \\
    &\lec{}{\leqslant} 1,
\end{align}
where (a) is because $\mathbf{V}$ is an unitary matrix and $\lambda_1,\dots,\lambda_K$ are eigenvalues of $\widetilde{\mathbf{N}}$. Continuing from \eqref{eq:ridge_w_1}, we obtain:
\begin{align}
    \norm{\Phi^{1/2}\boldsymbol{\widetilde{\xi}}}_2 &= \norm{\left(\widetilde{\mathbf{N}}^T \widetilde{\mathbf{N}}+\lambda \mathbf{I}\right)^{-1} \widetilde{\mathbf{N}}^T \widetilde{\mathbf{N}}}_2 \leqslant \norm{\mathbf{z}}_2.
\end{align}
Let us define $\mathbf{x}_0:=\left(\mathbf{N}^T \mathbf{N}\right)^{-1} \mathbf{N}^T \mathbf{x}$. Thus, we have:
\begin{align}
    \mathbf{z}&=\left(\widetilde{\mathbf{N}}^T \widetilde{\mathbf{N}}\right)^{-1} \widetilde{\mathbf{N}}^T \mathbf{x}
    \nonumber \\
    &= \left(\Phi^{-1 / 2}\mathbf{N}^T\mathbf{N} \Phi^{-1 / 2}\right)^{-1}\Phi^{-1 / 2} \mathbf{N}^T \mathbf{x}
    \nonumber \\
    &= \Phi^{1 / 2}\left(\mathbf{N}^T\mathbf{N}\right)^{-1}\mathbf{N}^T \mathbf{x}
    \nonumber \\
    &= \Phi^{1 / 2}\mathbf{x}_0.
\end{align}
Therefore, we can bound the $\norm{\boldsymbol{\xi}}_\Phi:=\sqrt{\boldsymbol{\xi}^T\Phi\boldsymbol{\xi}}$:
\begin{align}
    \norm{\boldsymbol{\xi}}_\Phi = \norm{\Phi^{1/2}\boldsymbol{\xi}}_2 \leqslant \norm{\mathbf{z}}_2 =  \norm{\mathbf{x}_0}_\Phi.
\end{align}
Since $\Phi$ is symmetric, by Rayleigh-Ritz theorem we know that:
\begin{align}
    0 \lec{(a)}{<} \lambda^{\Phi}_\textrm{min} \leqslant \frac{\boldsymbol{\xi}^T\Phi\boldsymbol{\xi}}{\boldsymbol{\xi}^T\boldsymbol{\xi}} = \frac{\norm{\boldsymbol{\xi}}^2_\Phi}{\norm{\boldsymbol{\xi}}^2_2} \leqslant 
    \lambda^{\Phi}_\textrm{max},
\end{align}
where $\lambda^{\Phi}_\textrm{min}, \lambda^{\Phi}_\textrm{max}$ are minimum and maximum eigenvalues of $\Phi$, and (a) is due to the fact that since $\Phi$ is kernel matrix of RKHS space and $\{b_k(
\cdot)\}^K_{k=1}$ are basis functions, it is positive definite. Thus, we have:
\begin{align}
    \norm{\boldsymbol{\xi}}^2_2 \leqslant \frac{1}{\lambda^{\Phi}_\textrm{min}}\norm{\boldsymbol{\xi}}^2_\Phi \leqslant \frac{1}{\lambda^{\Phi}_\textrm{min}}\cdot\norm{\mathbf{x}_0}^2_\Phi \leqslant \frac{\lambda^{\Phi}_\textrm{max}}{\lambda^{\Phi}_\textrm{min}}\norm{\mathbf{x}_0}^2_2=\kappa(\Phi)\norm{\mathbf{x}_0}^2_2.
\end{align}
Applying the bound for condition number introduce in \citep{guggenheimer1995simple}, we can complete the proof:
\begin{align}
    \kappa(\Phi)<\frac{2}{|\operatorname{det} \Phi|}\left(\frac{\|\Phi\|_F^2}{{K+4}}\right)^{\frac{{K+4}}{2}},
\end{align}
where $\norm{\cdot}_F$ is the Frobenius norm.
\end{proof}
Using Lemma~\ref{lem:ridge_reg_weight_upperbound}, \eqref{eq:spline_l_inf_2}, and \eqref{eq:spline_l_inf_p_2} we have:
\begin{align}\label{eq:spline_htild_bound}
    \norm{u^*}_{\hiltilde{2}} &\leqslant \norm{\boldsymbol{\xi}}^2_2\cdot\left(\sum^{K+4}_{k=1}\norm{b_k}^2_{\lp{\infty}} + \sum^{K+4}_{k=1}\norm{b'_k}^2_{\lp{\infty}}\right) + \int_\Omega |\widetilde{\encs}''(t)|^2\,dt \nonumber \\
    &\lec{(a)}{\leqslant} \frac{2}{|\operatorname{det} \Phi|}\left(\frac{\|\Phi\|_F^2}{{K+4}}\right)^{\frac{{K+4}}{2}} \norm{\boldsymbol{\widetilde{\xi}}}^2_2\left(\sum^{K+4}_{k=1}\norm{b_k}^2_{\lp{\infty}} + \sum^{K+4}_{k=1}\norm{b'_k}^2_{\lp{\infty}}\right) \nonumber \\ &+ \int_\Omega |\widetilde{\encs}''(t)|^2,
\end{align}
where (a) follows by applying Lemma~\ref{lem:ridge_reg_weight_upperbound}. Setting $M$ equal to the right-hand side of Equation \eqref{eq:spline_htild_bound} completes the proof.
\subsection{Proof of Theorem~\ref{th:conv_rate}}\label{app:proof_convrate}
Consider a natural spline $\widetilde{\encs}(t)$ as the encoder function fitted on the data points $\{(\alpha_k, x_k)\}^K_{k=1}$. Let $\encs^*(t)$ denote the optimal encoder minimizing the upper bound in \eqref{eq:optimcal_encdoer}. Then, we have:
\begin{align}
    \mathcal{R}(\hat{f}) &\lec{}{\leqslant} \frac{2q^2}{K} \sum^K_{k=1} (\encs^*(\alpha_k) - x_k)^2 + \enclamb \cdot g\big(\norm{\encs^*}^2_{\hiltilde{2}}\big) \nonumber \\
    &\lec{(a)}{\leqslant} \frac{2q^2}{K} \sum^K_{k=1} (\widetilde{\encs}(\alpha_k) - x_k)^2 + \enclamb \cdot g\big(\norm{\widetilde{\encs}}^2_{\hiltilde{2}}\big) \nonumber \\
    &\lec{(b)}{=} \enclamb \cdot g\big(\norm{\widetilde{\encs}}^2_{\hiltilde{2}}\big),
\end{align}
where (a) follows from the optimality of $\encs^*(t)$, and (b) is due to the fact that $\widetilde{\encs}(\alpha_k) = x_k$, since $\widetilde{\encs}(\cdot)$ is a natural spline. Note that $g\big(\norm{\widetilde{\encs}}^2_{\hiltilde{2}}\big)$ is independent of $N$ and $S$, and depends only on $\alpha_k$ and $x_k$ for $k \in [K]$. Additionally, based on the Theorem~\ref{th:optimcal_encdoer} and \eqref{eq:optim_lambda_rate}, $\enclamb = \mathcal{O}(S^3N^{-3})$ and $\enclamb = \mathcal{O}(S^\frac{8}{5}N^{-\frac{3}{5}})$ for the noiseless and noisy cases, respectively. Thus, the upper bound provided in \eqref{eq:optimcal_encdoer} converges at most at the rate of $\mathcal{O}(S^3N^{-3})$ for the noiseless case and $\mathcal{O}(S^\frac{8}{5}N^{-\frac{3}{5}})$ for the noisy case.

% Theorems~\ref{th:conv_rate_noiseless} and \ref{th:conv_rate_noisy} establish that $\ldec$ has the convergence rate of $\mathcal{O}\left(S^{4}N^{-3}\right)$ and 
% $\mathcal{O}\left(S^{\frac{8}{5}}N^{-\frac{3}{5}}\right)$, respectively. By designing the encoder such that $\encs(\alpha_k) = x_k$ for $k \in [K]$, we can achieve $\lenc = 0$ and consequently, the rate of convergence of the total loss, $\mathcal{R}(\hat{f})$, becomes $\mathcal{O}\left(S^{4}N^{-3}\right)$ and 
% $\mathcal{O}\left(S^{\frac{8}{5}}N^{-\frac{3}{5}}\right)$ for noiseless and noisy setting, respectively.
% \section{Performance Analysis}
% Figure~\ref{fig:comp_all} present the performance comparison between $\nprcc$ and \texttt{BACC} in a highly redundant regime, when the system is over-designed with a larger values of $\frac{N}{K}$. As depicted in Figure~\ref{fig:comp_all} the performance gain is relatively smaller compare to the case shown in Figure~\ref{fig:comp_all_better} specifically in the RelAcc metric.
% \begin{figure}[h]
%      \centering
%      \begin{subfigure}[b]{0.32\textwidth}
%          \centering
%          \captionsetup{justification=centering} 
%          \includegraphics[width=\textwidth]{figs/perf_comp/mse_vit_20_8_conf_1e-06_0.001_95.png}
%          \label{fig:comp_all_mse_vit}
%      \end{subfigure}
%      \hfill
%      \begin{subfigure}[b]{0.32\textwidth}
%          \centering
%          \captionsetup{justification=centering} 
%          \includegraphics[width=\textwidth]{figs/perf_comp/mse_repvgg_conf_60_20.png}
%          \label{fig:comp_all_mse_cifar}
%      \end{subfigure}
%      \hfill
%      \begin{subfigure}[b]{0.32\textwidth}
%          \centering
%          \captionsetup{justification=centering} 
%          \includegraphics[width=\textwidth]{figs/perf_comp/mse_lenet_comf_100_20.png}
%          \label{fig:comp_all_mse_mnist}
%      \end{subfigure}
%      \begin{subfigure}[b]{0.32\textwidth}
%          \centering
%          \captionsetup{justification=centering} 
%          \includegraphics[width=\textwidth]{figs/perf_comp/racc_vit_20_8_conf_1e-06_0.001_95.png}
%          \label{fig:comp_all_racc_vit}
%      \end{subfigure}
%      \hfill
%      \begin{subfigure}[b]{0.32\textwidth}
%          \centering
%          \captionsetup{justification=centering} 
%          \includegraphics[width=\textwidth]{figs/perf_comp/racc_repvgg_conf_60_20_png.png}
%          \label{fig:comp_all_racc_cifar}
%      \end{subfigure}
%      \hfill
%      \begin{subfigure}[b]{0.32\textwidth}
%          \centering
%          \captionsetup{justification=centering} 
%          \includegraphics[width=\textwidth]{figs/perf_comp/racc_lenet_comf_100_20.png}
%          \label{fig:comp_all_racc_mnist}
%      \end{subfigure}
%      %\captionsetup{font=small}
%      \caption{Performance comparison of $\nprcc$ and \texttt{BACC} with a $95\%$ confidence interval across a diverse range of stragglers for different models in a high redundant regime (larger $\frac{N}{K}$).}
%      \label{fig:comp_all}
% \end{figure}

\section{Comparison with Berrut Coded Computing}\label{app:comp_berrut}


\subsection{Convergence rate} \label{app:comp_berrut_conv_rate}
The upper bound of the infinity norm for the estimation provided in \citep{jahani2022berrut} for the coded computing scheme with $N$ workers and maximum of $S$ stragglers is as follows:
\begin{align}
\norm{\hat{f}_{\texttt{BACC}}(t)-f\circ\encs(t)}_{\lp{\infty}} \leqslant 2(1+R) \sin \left(\frac{(S+1) \pi}{2 N}\right)\norm{f\circ\encs^{\prime \prime}(t)}_{\lp{\infty}},
\end{align}
if $N-s$ is odd, and
\begin{align}
\norm{\hat{f}_{\texttt{BACC}}(t)-f\circ\encs(t)}_{\lp{\infty}} \leqslant 2(1+R) \sin \left(\frac{(S+1) \pi}{2 N}\right)&\Bigl(\norm{f\circ\encs^{\prime \prime}(t)}_{\lp{\infty}} \nonumber \\ &+\norm{f\circ\encs^{\prime}(t)}_{\lp{\infty}}\Bigr),
\end{align}
if $N-s$ is even, where $R=\frac{(s+1)(s+3) \pi^2}{4}$. Since $\norm{\cdot}_{\lp{2}}$ is upper bounded by $\norm{\cdot}_{\lp{\infty}}$, we can directly derive a convergence rate for the squared $\lp{2}$-norm of the error as $N$ increases:
\begin{align}
    \norm{\hat{f}_{\texttt{BACC}}(t)-f\circ\encs(t)}^2_{\lp{2}} \leqslant \mathcal{L}(\Omega) \cdot \norm{\hat{f}_{\texttt{BACC}}(t)-f\circ\encs(t)}^2_{\lp{\infty}} \leqslant \mathcal{O}(S^4N^{-2}).
\end{align}
Compared to our results, the upper bound for $\nprcc$ provided in Theorem~\ref{th:conv_rate_noiseless}, $\mathcal{O}(S^3N^{-3})$, is less sensitive to the number of stragglers and converges faster with increasing $N$.
Note that, since the $\norm{\cdot}_{\lp{2}}$ is upper bounded by $\norm{\cdot}_{\lp{\infty}}$, the statement above does not guarantee faster convergence of the proposed scheme compared to Berrut approach.

It should be noted that \citep{jahani2022berrut} does not analyze the noisy setting. 


\subsection{Computational complexity} \label{app:comp_berrut_complexity}
From the experimental view, we compare the whole encoding and decoding time (on a single CPU-only machine) for $\nprcc$ and \texttt{BACC} frameworks, as shown in the following table:

\setlength{\tabcolsep}{5pt}
\setlength{\extrarowheight}{2pt}
\begin{table*}[h]
%\captionsetup{font=small}
\caption{Average and std of end-to-end processing time of $\nprcc$ and \texttt{BACC} for different architectures}
\label{tab:complexity}
\begin{center}
\begin{small}
\begin{sc}
\begin{tabular}[h]{lll} 
& \texttt{BACC} & $\nprcc$ \\
\hline LeNet5, $(N, K)=(100,20)$ & $0.013 s \pm 0.002$ & $0.007 s \pm 0.001$ \\
 \hline RepVGG, $(N, K)=(60,20)$ & $1.62 s \pm 0.18$ & $1.59 s \pm 0.14$ \\
 \hline ViT, $(N, K)=(20,8)$ & $1.60 s \pm 0.28$ & $1.74 s \pm 0.29$
\end{tabular}
\end{sc}
\end{small}
\end{center}
\end{table*}

As shown in Table~\ref{tab:complexity}, the end-to-end processing time of the proposed framework is on par with \texttt{BACC}.

\section{Comparison with Lagrange Coded Computing}\label{app:comp_lagrange}
Although the {\bf only} existing coded computing scheme for general functions is Berrut coded computing \citep{jahani2022berrut}, with which we have compared our proposed scheme, other schemes are designed for specific computations, such as polynomial functions \citep{yu2019lagrange} and matrix multiplication \citep{yu2017polynomial}.
To provide further comparison, we evaluate our proposed scheme against Lagrange coded computing (\texttt{LCC}) \citep{yu2019lagrange}, which is specifically designed for polynomial computations, as follows:


\subsection{Accuracy of function approximation} \texttt{LCC} is applicable only to polynomial computing functions \citep{yu2019lagrange}. Additionally, to enable recovery, the number of servers required must be at least $(K-1) \times \text{deg}(f) + S + 1$ worker nodes \citep{yu2019lagrange, jahani2022berrut}; otherwise, the master node cannot recover any results. Moreover, \texttt{LCC} is designed for computation over finite fields and encounters serious instability when computing over real numbers, particularly when $(K-1) \times \text{deg}(f)$ is around $25$ or higher \citep{jahani2022berrut,gautschi1987lower}.

We compare the proposed framework with \texttt{LCC} in Figure~\ref{fig:comp_all_conf_poly}. Note that if $N < (K-1) \times \text{deg}(f) + S + 1$, \texttt{LCC} cannot operate effectively. To adapt \texttt{LCC} for such cases, we approximate results by fitting a lower-degree polynomial to the available workers' outputs. We run $\nprcc$ and \texttt{LCC} on the same set of input data and a fixed polynomial function for 20 trials, plotting the average performance and corresponding $95\%$ confidence intervals in Figure~\ref{fig:comp_all_conf_poly}. Figures~\ref{fig:comp_all_conf_poly_a} and \ref{fig:comp_all_conf_poly_b} illustrate the performances of \texttt{LCC} and $\nprcc$ for a low-degree polynomial and a small number of data points ($\text{deg}(f) = 3$ and $K = 5$). In contrast, Figures~\ref{fig:comp_all_conf_poly_c} and \ref{fig:comp_all_conf_poly_d} show performance for a higher-degree polynomial and a larger dataset ($\text{deg}(f) = 15$ and $K = 10$).
As shown in Figures~\ref{fig:comp_all_conf_poly_a} and \ref{fig:comp_all_conf_poly_b}, \texttt{LCC} achieves exact results for $S \leq 7$. However, at larger values of $S$, as well as larger polynomial degree (as in Figures~\ref{fig:comp_all_conf_poly_c} and \ref{fig:comp_all_conf_poly_d}), the proposed approach, without any parameter tuning, outperforms \texttt{LCC} in terms of both computational stability (lower variance) and recovery accuracy.

\begin{figure}[h]
     \begin{subfigure}[h]{0.24\textwidth}
         \centering 
         \includegraphics[width=\textwidth]{figs/rmse_lagrange_poly_20_5.png}
         \caption{}
         \label{fig:comp_all_conf_poly_a}
     \end{subfigure}
     \hfill
     \begin{subfigure}[h]{0.24\textwidth}
         \centering
         % \captionsetup{justification=centering} 
         \includegraphics[width=\textwidth]{figs/rmse_lagrange_poly_20_5_2.png}
         \caption{}
         \label{fig:comp_all_conf_poly_b}
     \end{subfigure}
     \hfill
     \begin{subfigure}[h]{0.24\textwidth}
         \centering
         % \captionsetup{justification=centering} 
         \includegraphics[width=\textwidth]
         {figs/rmse_lagrange_poly_20_10.png}
         \caption{}
         \label{fig:comp_all_conf_poly_c}
     \end{subfigure}
     \hfill
     \begin{subfigure}[h]{0.24\textwidth}
         \centering
         % \captionsetup{justification=centering} 
         \includegraphics[width=\textwidth]{figs/rmse_lagrange_poly_20_10_2.png}
         \caption{}
         \label{fig:comp_all_conf_poly_d}
     \end{subfigure}
     \caption{Average performance of $\nprcc$ and Lagrange Coded Computing, with a 95\% confidence interval. Plots (a) and (d) show the overall performance, while the zoomed-in subplots (b) and (c) highlight the performance for smaller range of stragglers.}
     \label{fig:comp_all_conf_poly}
\end{figure}

\subsection{Computational complexity} Encoding and decoding complexities in \texttt{LCC} are $\mathcal{O}(N\cdot \log^2(K) \cdot \log\log (K)\cdot d)$ and $\mathcal{O}((N-S)\cdot \log^2((N-S)) \cdot \log\log ((N-S))\cdot m)$, respectively, where $d$ and $m$ are input and output dimensions of the computing function $f(\cdot)$, respectively \citep{yu2019lagrange}. In contrast, as mentioned before, for smoothing splines, the encoding and decoding process, which involves evaluation on new points and calculating the fitted coefficients, have the computational complexity of $\mathcal{O}(K.d)$ and $\mathcal{O}((N-s).m)$. Consequently, the computational complexity of the proposed scheme is less than \texttt{LCC}.
\section{Sensitivity Analysis}\label{app:sens_analysis}
\subsection{Sensitivity to number of stragglers}
The smoothing parameters for each model show low sensitivity to the number of stragglers (or worker nodes). To find the optimal smoothing parameter, we use cross-validation across different $\frac{S}{N}$ values. The following table presents the optimal smoothing parameters for selected numbers of stragglers for LeNet5 with $(N, K) = (100, 60)$ and RepVGG with $(N, K) = (60, 20)$, respectively. As shown in Table~\ref{tab:sensitivity_straggler}, the optimal values of $\lambda_\textrm{e}$ and $\lambda_\textrm{d}$ exhibit low sensitivity to the number of stragglers.
\setlength{\tabcolsep}{4pt}
\begin{table*}[!htb]
%\captionsetup{font=small}
\caption{Optimal smoothing parameters for different number of stragglers for LeNet and RepVGG architectures.}
\label{tab:sensitivity_straggler}
% \vskip 0.15in
\begin{center}
\begin{small}
\begin{sc}
\begin{tabular}{ccc|cc}
    \hline  & \\[-2.7ex]
    ~ %& 
    & \multicolumn{2}{c}{LeNet5} & \multicolumn{2}{c}{RepVGG}  \\
    $(N, K)$
    &  \multicolumn{2}{c}{$(100, 60)$}  & \multicolumn{2}{c}{$(60, 20)$} \\
    \cline{1-3}
    \cline{4-5}
     & \\[-2.7ex]
     S & \multicolumn{1}{c}{$\lambda^*_{e}$} & \multicolumn{1}{c}{$\lambda^*_{d}$} & \multicolumn{1}{c}{$\lambda^*_{e}$} & \multicolumn{1}{c}{$\lambda^*_{d}$} \\
    \hline  & \\[-1.5ex]
    0 & $10^{-13}$ & $10^{-6}$ & $10^{-6}$ & $10^{-4}$ \\
    5 & $10^{-13}$ & $10^{-6}$ & $10^{-6}$ & $10^{-4}$ \\
    10 & $10^{-13}$ & $10^{-6}$  & $10^{-5}$ & $10^{-4}$\\
    15 & $10^{-13}$ & $10^{-6}$  & $10^{-5}$ & $10^{-4}$\\
    20 & $10^{-13}$ & $10^{-6}$ & $10^{-5}$ & $10^{-4}$ \\
    25 & $10^{-8}$ & $10^{-5}$ & $10^{-5}$ & $10^{-4}$\\
    30 & $10^{-8}$ & $10^{-4}$ & $10^{-5}$ & $10^{-3}$ \\
    35 & $10^{-8}$ & $10^{-4}$ & $10^{-5}$ & $10^{-3}$\\
    \bottomrule
\end{tabular}
\end{sc}
\end{small}
\end{center}
% \vskip -0.1in
\end{table*}
\subsection{Sensitivity to smoothing parameters}
To assess the performance of the proposed scheme with respect to the smoothing parameters, we vary each parameter individually around its optimal point while holding the other parameter fixed at their optimal value. We then record the average percentage increase in RMSE relative to the RMSE at the optimal point. Figure~\ref{fig:sm_sens} presents these results for LeNet with $(N, K, S) = (100, 60, 20)$ (Figures~\ref{fig:dec_sens_mnist} and ~\ref{fig:enc_sens_mnist}) and for RepVGG with $(N, K, S) = (60, 20, 35)$ (Figures~\ref{fig:dec_sens_rep} and ~\ref{fig:enc_sens_rep}).

\begin{figure}[h]
     \begin{subfigure}[h]{0.49\textwidth}
         \centering 
         \includegraphics[width=\textwidth]{figs/perf_comp/decoder_sense_mnist.png}
         \caption{}
         \label{fig:dec_sens_mnist}
     \end{subfigure}
     \hfill
     \begin{subfigure}[h]{0.49\textwidth}
         \centering
         % \captionsetup{justification=centering} 
         \includegraphics[width=\textwidth]{figs/perf_comp/encoder_sens_mnist.png}
         \caption{}
         \label{fig:enc_sens_mnist}
     \end{subfigure}
     \hfill
     \begin{subfigure}[h]{0.49\textwidth}
         \centering
         % \captionsetup{justification=centering} 
         \includegraphics[width=\textwidth]{figs/perf_comp/decoder_sense_rep.png}
         \caption{}
         \label{fig:dec_sens_rep}
     \end{subfigure}
     \hfill
     \begin{subfigure}[h]{0.49\textwidth}
         \centering
         % \captionsetup{justification=centering} 
         \includegraphics[width=\textwidth]{figs/perf_comp/encoder_sens_rep.png}
         \caption{}
         \label{fig:enc_sens_rep}
     \end{subfigure}
     \caption{Sensitivity of $\nprcc$ performance with respect to $\log_{10}(\declamb)$ and $\log_{10}(\enclamb)$. The yellow line represents the performance when the variable smoothing parameter is set to zero.}
     \label{fig:sm_sens}
\end{figure}
As shown in Figure~\ref{fig:sm_sens}, the presence of more stragglers increases the sensitivity of $\nprcc$ with respect to its smoothing parameter. However, even in a high-straggler regime, the RMSE increases by only around $3\%$ when the smoothing parameter deviates from its optimal value by a scale of $10$.
\section{High-dimensional computing function}\label{app:hdim_f}
Let us consider more general cases where $f = [f_1,\dots,f_m]$ is a vector-valued function, where each component function $f_j:\mathbb{R} \to \mathbb{R}$ is $q_j$-Lipschitz continuous. Based on \eqref{eq:decompose}, we have:
\begin{align}
    \mathcal{R}(\fhat)&\leqslant \mathop{\mathbb{E}}_{\bm{\epsilon}, \stset \sim F_{S,N}} \left[\frac{2}{K} \sum^K_{k=1} \norm{\dec(\alpha_k) - \func(\encs(\alpha_k))}^2_2\right] +  \frac{2}{K} \sum^K_{k=1} \norm{\func(\encs(\alpha_k)) - \func(x_k)}^2_2.\nonumber \\
    &\leqslant \mathop{\mathbb{E}}_{\bm{\epsilon}, \stset \sim F_{S,N}}\left[\frac{2}{K} \sum^K_{k=1} \sum^m_{j=1} \left(\decsj(\alpha_k) - f_j(\encs(\alpha_k))\right)^2_2\right] + \frac{2\sum^m_{j=1}q_j^2}{K} \sum^K_{k=1} \norm{\encs(\alpha_k) - x_k}^2_2 \nonumber \\ \nonumber
    &= \sum^m_{j=1} \mathop{\mathbb{E}}_{\bm{\epsilon}, \stset \sim F_{S,N}}\left[\frac{2}{K} \sum^K_{k=1}  \left(\decsj(\alpha_k) - f_j(\encs(\alpha_k))\right)^2_2\right] + \frac{2\sum^m_{j=1}q_j^2}{K} \sum^K_{k=1} \norm{\encs(\alpha_k) - x_k}^2_2 
\end{align}
Let us define the following objective for the decoder function:
\begin{align}\label{eq:decoder_opt_highdim}
    \dec^\star=\underset{\mathbf{u} \in \hilm{2}}{\operatorname{argmin}} \frac{1}{|\stset|} \sum_{v \in \stset}\norm{\mathbf{u}\left(\beta_v\right)-\func\left(\encs\left(\beta_v\right)\right)}^2_2 + \sum^m_{j=1} \declamb \int_{\Omega} \left(u_j''(t)\right)^2\,dt.
\end{align}
The solution to \eqref{eq:decoder_opt_highdim}, denoted as $\dec^\star$, is a vector-valued function, where each component $\decsj(\cdot)$ is a smoothing spline function fitted to the data points $\left\{\left(\beta_v, f_j\left(\encs\left(\beta_v\right)\right)\right)\right\}_{v \in \stset}$. As a result, By defining $q=\sqrt{\sum^m_{j=1}q_j^2}$ and scaling up all upper bounds for $\ldec$ by a factor of $m$, all previous results and theorems seamlessly extend to high-dimensional computing functions.
\section{Coded data points}\label{sec:apx_coded_sample}
% \subsection{Performance comparison}\label{sec:apx_perf}

% As depicted in Figure~\ref{fig:comp_berrut_all}, the $\nprcc$ scheme consistently outperforms the \texttt{BACC} scheme across a wide range of straggler numbers, for different model architecture, and in terms of both MSE and RelAcc metrics.

% \subsection{Smoothing Effect}\label{sec:apx_smoothing_effect}
% Figure~\ref{fig:comp_ns} illustrates the impact of using smoothing parameters in $\nprcc$ and presents the overall performance comparison between $\nprcc$ and its model-agnostic version, $\nprcc^{Ag}$, along with the \texttt{BACC} framework. It is evident that for less complex functions like LeNet, both $\nprcc$ and $\nprcc^{Ag}$ yield the same results. However, for more complex functions such as ViT and RepVGG, while $\nprcc^{Ag}$ performs better in scenarios with a lower number of stragglers (benefiting from data point redundancy), it notably falls behind $\nprcc$ in higher straggler regimes.


% \begin{figure}[h]
%      \centering
%      \begin{subfigure}[b]{0.3\textwidth}
%          \centering
%          \captionsetup{justification=centering} 
%          \includegraphics[width=\textwidth]{figs/mse_mnist_comp_ns.png}
%          \caption{$(N, K)=(200, 50)$,\\ LeNet model}
%          \label{fig:comp_ns_mse_mnist}
%      \end{subfigure}
%      \hfill
%      \begin{subfigure}[b]{0.3\textwidth}
%          \centering
%          \captionsetup{justification=centering} 
%          \includegraphics[width=\textwidth]{figs/mse_cifar_comp_ns.png}
%          \caption{$(N, K)=(200, 50)$,\\ RepVGG model}
%          \label{fig:comp_ns_mse_cifar}
%      \end{subfigure}
%      \hfill
%      \begin{subfigure}[b]{0.3\textwidth}
%          \centering
%          \captionsetup{justification=centering} 
%          \includegraphics[width=\textwidth]{figs/mse_vit_comp_ns.png}
%          \caption{$(N, K)=(100, 25)$,\\ ViT model}
%          \label{fig:comp_ns_mse_vit}
%      \end{subfigure}
%         \caption{Comparison of Mean square error for $\nprcc^{Ag}$ and $\nprcc, \texttt{BACC}$ over different number of stragglers for \Cref{fig:comp_ns_mse_mnist} LeNet, \Cref{fig:comp_ns_mse_cifar} RepVGG, and \Cref{fig:comp_ns_mse_vit} ViT networks as the machine learning models.}
%         \label{fig:comp_ns}
% \end{figure}
% \subsection{Coded samples}\label{sec:apx_coded_sample}
Figures~\ref{fig:coded_samples_b} and \ref{fig:coded_samples_c} display coded samples generated by \texttt{BACC} and $\nprcc$, respectively, derived from the same initial data points depicted in Figure~\ref{fig:coded_samples_a}. These samples are presented for the MNIST dataset with parameters $(N, K) = (70, 30)$. From the figures, it is apparent (Specifically in paired ones that are shown with the same color) that while both schemes' coded samples are a weighted combination of multiple initial samples, \texttt{BACC}'s coded samples exhibit high-frequency noise. This observation suggests that $\nprcc$ regression functions produce more refined coded samples without any disruptive noise.
\begin{figure}[h]
     \centering
     \begin{subfigure}[b]{\textwidth}
         \centering
         \includegraphics[width=0.31\textwidth]{figs/original_samples_mnist.png}
         \caption{Initial inputs}
         \label{fig:coded_samples_a}
     \end{subfigure}
     \hfill
     \begin{subfigure}[b]{0.42\textwidth}
         \centering
         \includegraphics[width=\textwidth]{figs/berrut_mnist_coded_samples_high_contrast.png}
         \caption{\texttt{BACC} coded samples}
         \label{fig:coded_samples_b}
     \end{subfigure}
     \hfill
     \begin{subfigure}[b]{0.42\textwidth}
         \centering
         \includegraphics[width=\textwidth]{figs/spline_mnist_coded_samples_high_contrast.png}
         \caption{$\nprcc$ coded samples}
         \label{fig:coded_samples_c}
     \end{subfigure}
     %\captionsetup{font=small}
        \caption{Comparison of coded samples between \texttt{BACC} and $\nprcc$ frameworks. Figure~\ref{fig:coded_samples_a} represents the initial data points $\{\mathbf{x}_k\}^K_{k=1}$ for $K=30$. Figures~\ref{fig:coded_samples_b} and \ref{fig:coded_samples_c} display $N=70$ coded samples $\{\mathbf{\Tilde{x}}_n\}^N_{n=1}$ from \texttt{BACC} and $\nprcc$, respectively. Samples with clear differences are highlighted with the same color.}
        \label{fig:coded_samples_all}
\end{figure}

% \section{Submission of papers to NeurIPS 2023}


% Please read the instructions below carefully and follow them faithfully. \textbf{Important:} This year the checklist will be submitted separately from the main paper in OpenReview, please review it well ahead of the submission deadline: \url{https://neurips.cc/public/guides/PaperChecklist}.


% \subsection{Style}


% Papers to be submitted to NeurIPS 2023 must be prepared according to the
% instructions presented here. Papers may only be up to {\bf nine} pages long,
% including figures. Additional pages \emph{containing only acknowledgments and
% references} are allowed. Papers that exceed the page limit will not be
% reviewed, or in any other way considered for presentation at the conference.


% The margins in 2023 are the same as those in previous years.


% Authors are required to use the NeurIPS \LaTeX{} style files obtainable at the
% NeurIPS website as indicated below. Please make sure you use the current files
% and not previous versions. Tweaking the style files may be grounds for
% rejection.


% \subsection{Retrieval of style files}


% The style files for NeurIPS and other conference information are available on
% the website at
% \begin{center}
%   \url{http://www.neurips.cc/}
% \end{center}
% The file \verb+neurips_2023.pdf+ contains these instructions and illustrates the
% various formatting requirements your NeurIPS paper must satisfy.


% The only supported style file for NeurIPS 2023 is \verb+neurips_2023.sty+,
% rewritten for \LaTeXe{}.  \textbf{Previous style files for \LaTeX{} 2.09,
%   Microsoft Word, and RTF are no longer supported!}


% The \LaTeX{} style file contains three optional arguments: \verb+final+, which
% creates a camera-ready copy, \verb+preprint+, which creates a preprint for
% submission to, e.g., arXiv, and \verb+nonatbib+, which will not load the
% \verb+natbib+ package for you in case of package clash.


% \paragraph{Preprint option}
% If you wish to post a preprint of your work online, e.g., on arXiv, using the
% NeurIPS style, please use the \verb+preprint+ option. This will create a
% nonanonymized version of your work with the text ``Preprint. Work in progress.''
% in the footer. This version may be distributed as you see fit, as long as you do not say which conference it was submitted to. Please \textbf{do
%   not} use the \verb+final+ option, which should \textbf{only} be used for
% papers accepted to NeurIPS. 


% At submission time, please omit the \verb+final+ and \verb+preprint+
% options. This will anonymize your submission and add line numbers to aid
% review. Please do \emph{not} refer to these line numbers in your paper as they
% will be removed during generation of camera-ready copies.


% The file \verb+neurips_2023.tex+ may be used as a ``shell'' for writing your
% paper. All you have to do is replace the author, title, abstract, and text of
% the paper with your own.


% The formatting instructions contained in these style files are summarized in
% Sections \ref{gen_inst}, \ref{headings}, and \ref{others} below.


% \section{General formatting instructions}
% \label{gen_inst}


% The text must be confined within a rectangle 5.5~inches (33~picas) wide and
% 9~inches (54~picas) long. The left margin is 1.5~inch (9~picas).  Use 10~point
% type with a vertical spacing (leading) of 11~points.  Times New Roman is the
% preferred typeface throughout, and will be selected for you by default.
% Paragraphs are separated by \nicefrac{1}{2}~line space (5.5 points), with no
% indentation.


% The paper title should be 17~point, initial caps/lower case, bold, centered
% between two horizontal rules. The top rule should be 4~points thick and the
% bottom rule should be 1~point thick. Allow \nicefrac{1}{4}~inch space above and
% below the title to rules. All pages should start at 1~inch (6~picas) from the
% top of the page.


% For the final version, authors' names are set in boldface, and each name is
% centered above the corresponding address. The lead author's name is to be listed
% first (left-most), and the co-authors' names (if different address) are set to
% follow. If there is only one co-author, list both author and co-author side by
% side.


% Please pay special attention to the instructions in Section \ref{others}
% regarding figures, tables, acknowledgments, and references.


% \section{Headings: first level}
% \label{headings}


% All headings should be lower case (except for first word and proper nouns),
% flush left, and bold.


% First-level headings should be in 12-point type.


% \subsection{Headings: second level}


% Second-level headings should be in 10-point type.


% \subsubsection{Headings: third level}


% Third-level headings should be in 10-point type.


% \paragraph{Paragraphs}


% There is also a \verb+\paragraph+ command available, which sets the heading in
% bold, flush left, and inline with the text, with the heading followed by 1\,em
% of space.


% \section{Citations, figures, tables, references}
% \label{others}


% These instructions apply to everyone.


% \subsection{Citations within the text}


% The \verb+natbib+ package will be loaded for you by default.  Citations may be
% author/year or numeric, as long as you maintain internal consistency.  As to the
% format of the references themselves, any style is acceptable as long as it is
% used consistently.


% The documentation for \verb+natbib+ may be found at
% \begin{center}
%   \url{http://mirrors.ctan.org/macros/latex/contrib/natbib/natnotes.pdf}
% \end{center}
% Of note is the command \verb+\citet+, which produces citations appropriate for
% use in inline text.  For example,
% \begin{verbatim}
%    \citep{hasselmo} investigated\dots
% \end{verbatim}
% produces
% \begin{quote}
%   Hasselmo, et al.\ (1995) investigated\dots
% \end{quote}


% If you wish to load the \verb+natbib+ package with options, you may add the
% following before loading the \verb+neurips_2023+ package:
% \begin{verbatim}
%    \PassOptionsToPackage{options}{natbib}
% \end{verbatim}


% If \verb+natbib+ clashes with another package you load, you can add the optional
% argument \verb+nonatbib+ when loading the style file:
% \begin{verbatim}
%    \usepackage[nonatbib]{neurips_2023}
% \end{verbatim}


% As submission is double blind, refer to your own published work in the third
% person. That is, use ``In the previous work of Jones et al.\ [4],'' not ``In our
% previous work [4].'' If you cite your other papers that are not widely available
% (e.g., a journal paper under review), use anonymous author names in the
% citation, e.g., an author of the form ``A.\ Anonymous'' and include a copy of the anonymized paper in the supplementary material.


% \subsection{Footnotes}


% Footnotes should be used sparingly.  If you do require a footnote, indicate
% footnotes with a number\footnote{Sample of the first footnote.} in the
% text. Place the footnotes at the bottom of the page on which they appear.
% Precede the footnote with a horizontal rule of 2~inches (12~picas).


% Note that footnotes are properly typeset \emph{after} punctuation
% marks.\footnote{As in this example.}


% \subsection{Figures}


% \begin{figure}
%   \centering
%   \fbox{\rule[-.5cm]{0cm}{4cm} \rule[-.5cm]{4cm}{0cm}}
%   \caption{Sample figure caption.}
% \end{figure}


% All artwork must be neat, clean, and legible. Lines should be dark enough for
% purposes of reproduction. The figure number and caption always appear after the
% figure. Place one line space before the figure caption and one line space after
% the figure. The figure caption should be lower case (except for first word and
% proper nouns); figures are numbered consecutively.


% You may use color figures.  However, it is best for the figure captions and the
% paper body to be legible if the paper is printed in either black/white or in
% color.


% \subsection{Tables}


% All tables must be centered, neat, clean and legible.  The table number and
% title always appear before the table.  See Table~\ref{sample-table}.


% Place one line space before the table title, one line space after the
% table title, and one line space after the table. The table title must
% be lower case (except for first word and proper nouns); tables are
% numbered consecutively.


% Note that publication-quality tables \emph{do not contain vertical rules.} We
% strongly suggest the use of the \verb+booktabs+ package, which allows for
% typesetting high-quality, professional tables:
% \begin{center}
%   \url{https://www.ctan.org/pkg/booktabs}
% \end{center}
% This package was used to typeset Table~\ref{sample-table}.


% \begin{table}
%   \caption{Sample table title}
%   \label{sample-table}
%   \centering
%   \begin{tabular}{lll}
%     \toprule
%     \multicolumn{2}{c}{Part}                   \\
%     \cmidrule(r){1-2}
%     Name     & Description     & Size ($\mu$m) \\
%     \midrule
%     Dendrite & Input terminal  & $\sim$100     \\
%     Axon     & Output terminal & $\sim$10      \\
%     Soma     & Cell body       & up to $10^6$  \\
%     \bottomrule
%   \end{tabular}
% \end{table}

% \subsection{Math}
% Note that display math in bare TeX commands will not create correct line numbers for submission. Please use LaTeX (or AMSTeX) commands for unnumbered display math. (You really shouldn't be using \$\$ anyway; see \url{https://tex.stackexchange.com/questions/503/why-is-preferable-to} and \url{https://tex.stackexchange.com/questions/40492/what-are-the-differences-between-align-equation-and-displaymath} for more information.)

% \subsection{Final instructions}

% Do not change any aspects of the formatting parameters in the style files.  In
% particular, do not modify the width or length of the rectangle the text should
% fit into, and do not change font sizes (except perhaps in the
% \textbf{References} section; see below). Please note that pages should be
% numbered.


% \section{Preparing PDF files}


% Please prepare submission files with paper size ``US Letter,'' and not, for
% example, ``A4.''


% Fonts were the main cause of problems in the past years. Your PDF file must only
% contain Type 1 or Embedded TrueType fonts. Here are a few instructions to
% achieve this.


% \begin{itemize}


% \item You should directly generate PDF files using \verb+pdflatex+.


% \item You can check which fonts a PDF files uses.  In Acrobat Reader, select the
%   menu Files$>$Document Properties$>$Fonts and select Show All Fonts. You can
%   also use the program \verb+pdffonts+ which comes with \verb+xpdf+ and is
%   available out-of-the-box on most Linux machines.


% \item \verb+xfig+ "patterned" shapes are implemented with bitmap fonts.  Use
%   "solid" shapes instead.


% \item The \verb+\bbold+ package almost always uses bitmap fonts.  You should use
%   the equivalent AMS Fonts:
% \begin{verbatim}
%    \usepackage{amsfonts}
% \end{verbatim}
% followed by, e.g., \verb+\mathbb{R}+, \verb+\mathbb{N}+, or \verb+\mathbb{C}+
% for $\mathbb{R}$, $\mathbb{N}$ or $\mathbb{C}$.  You can also use the following
% workaround for reals, natural and complex:
% \begin{verbatim}
%    \newcommand{\RR}{I\!\!R} %real numbers
%    \newcommand{\Nat}{I\!\!N} %natural numbers
%    \newcommand{\CC}{I\!\!\!\!C} %complex numbers
% \end{verbatim}
% Note that \verb+amsfonts+ is automatically loaded by the \verb+amssymb+ package.


% \end{itemize}


% If your file contains type 3 fonts or non embedded TrueType fonts, we will ask
% you to fix it.


% \subsection{Margins in \LaTeX{}}


% Most of the margin problems come from figures positioned by hand using
% \verb+\special+ or other commands. We suggest using the command
% \verb+\includegraphics+ from the \verb+graphicx+ package. Always specify the
% figure width as a multiple of the line width as in the example below:
% \begin{verbatim}
%    \usepackage[pdftex]{graphicx} ...
%    \includegraphics[width=0.8\linewidth]{myfile.pdf}
% \end{verbatim}
% See Section 4.4 in the graphics bundle documentation
% (\url{http://mirrors.ctan.org/macros/latex/required/graphics/grfguide.pdf})


% A number of width problems arise when \LaTeX{} cannot properly hyphenate a
% line. Please give LaTeX hyphenation hints using the \verb+\-+ command when
% necessary.


% \begin{ack}
% Use unnumbered first level headings for the acknowledgments. All acknowledgments
% go at the end of the paper before the list of references. Moreover, you are required to declare
% funding (financial activities supporting the submitted work) and competing interests (related financial activities outside the submitted work).
% More information about this disclosure can be found at: \url{https://neurips.cc/Conferences/2023/PaperInformation/FundingDisclosure}.


% Do {\bf not} include this section in the anonymized submission, only in the final paper. You can use the \texttt{ack} environment provided in the style file to autmoatically hide this section in the anonymized submission.
% \end{ack}



% \section{Supplementary Material}

% Authors may wish to optionally include extra information (complete proofs, additional experiments and plots) in the appendix. All such materials should be part of the supplemental material (submitted separately) and should NOT be included in the main submission.


% \section*{References}


% References follow the acknowledgments in the camera-ready paper. Use unnumbered first-level heading for
% the references. Any choice of citation style is acceptable as long as you are
% consistent. It is permissible to reduce the font size to \verb+small+ (9 point)
% when listing the references.
% Note that the Reference section does not count towards the page limit.
% \medskip


% {
% \small


% [1] Alexander, J.A.\ \& Mozer, M.C.\ (1995) Template-based algorithms for
% connectionist rule extraction. In G.\ Tesauro, D.S.\ Touretzky and T.K.\ Leen
% (eds.), {\it Advances in Neural Information Processing Systems 7},
% pp.\ 609--616. Cambridge, MA: MIT Press.


% [2] Bower, J.M.\ \& Beeman, D.\ (1995) {\it The Book of GENESIS: Exploring
%   Realistic Neural Models with the GEneral NEural SImulation System.}  New York:
% TELOS/Springer--Verlag.


% [3] Hasselmo, M.E., Schnell, E.\ \& Barkai, E.\ (1995) Dynamics of learning and
% recall at excitatory recurrent synapses and cholinergic modulation in rat
% hippocampal region CA3. {\it Journal of Neuroscience} {\bf 15}(7):5249-5262.
% }

%%%%%%%%%%%%%%%%%%%%%%%%%%%%%%%%%%%%%%%%%%%%%%%%%%%%%%%%%%%%





%%%%%%%%%%%%%%%%%%%%%%%%%%%%%%%%%%%%%%%%%%%%%%%%%%%%%%%%%%%%
\newpage
\section*{NeurIPS Paper Checklist}
% %%% BEGIN INSTRUCTIONS %%%
% The checklist is designed to encourage best practices for responsible machine learning research, addressing issues of reproducibility, transparency, research ethics, and societal impact. Do not remove the checklist: {\bf The papers not including the checklist will be desk rejected.} The checklist should follow the references and follow the (optional) supplemental material.  The checklist does NOT count towards the page
% limit. 

% Please read the checklist guidelines carefully for information on how to answer these questions. For each question in the checklist:
% \begin{itemize}
%     \item You should answer \answerYes{}, \answerNo{}, or \answerNA{}.
%     \item \answerNA{} means either that the question is Not Applicable for that particular paper or the relevant information is Not Available.
%     \item Please provide a short (1–2 sentence) justification right after your answer (even for NA). 
%    % \item {\bf The papers not including the checklist will be desk rejected.}
% \end{itemize}

% {\bf The checklist answers are an integral part of your paper submission.} They are visible to the reviewers, area chairs, senior area chairs, and ethics reviewers. You will be asked to also include it (after eventual revisions) with the final version of your paper, and its final version will be published with the paper.

% The reviewers of your paper will be asked to use the checklist as one of the factors in their evaluation. While "\answerYes{}" is generally preferable to "\answerNo{}", it is perfectly acceptable to answer "\answerNo{}" provided a proper justification is given (e.g., "error bars are not reported because it would be too computationally expensive" or "we were unable to find the license for the dataset we used"). In general, answering "\answerNo{}" or "\answerNA{}" is not grounds for rejection. While the questions are phrased in a binary way, we acknowledge that the true answer is often more nuanced, so please just use your best judgment and write a justification to elaborate. All supporting evidence can appear either in the main paper or the supplemental material, provided in appendix. If you answer \answerYes{} to a question, in the justification please point to the section(s) where related material for the question can be found.

% IMPORTANT, please:
% \begin{itemize}
%     \item {\bf Delete this instruction block, but keep the section heading ``NeurIPS paper checklist"},
%     \item  {\bf Keep the checklist subsection headings, questions/answers and guidelines below.}
%     \item {\bf Do not modify the questions and only use the provided macros for your answers}.
% \end{itemize} 
 

% %%% END INSTRUCTIONS %%%


\begin{enumerate}

\item {\bf Claims}
    \item[] Question: Do the main claims made in the abstract and introduction accurately reflect the paper's contributions and scope?
    \item[] Answer: \answerYes{} % Replace by \answerYes{}, \answerNo{}, or \answerNA{}.
    \item[] Justification: 
    We detailed our contributions clearly in the abstract and the introduction sections of the paper.
    \item[] Guidelines:
    \begin{itemize}
        \item The answer NA means that the abstract and introduction do not include the claims made in the paper.
        \item The abstract and/or introduction should clearly state the claims made, including the contributions made in the paper and important assumptions and limitations. A No or NA answer to this question will not be perceived well by the reviewers. 
        \item The claims made should match theoretical and experimental results, and reflect how much the results can be expected to generalize to other settings. 
        \item It is fine to include aspirational goals as motivation as long as it is clear that these goals are not attained by the paper. 
    \end{itemize}

\item {\bf Limitations}
    \item[] Question: Does the paper discuss the limitations of the work performed by the authors?
    \item[] Answer: \answerYes{} % Replace by \answerYes{}, \answerNo{}, or \answerNA{}.
    \item[] Justification:We provided all the details regarding the assumptions, conditions, and limitations in the framework explanations (Section~\ref{sec:framework}), theorems (Section~\ref{sec:mainreults}), as well as the experiments section (Section~\ref{sec:exp_result}) in the paper.
    \item[] Guidelines:
    \begin{itemize}
        \item The answer NA means that the paper has no limitation while the answer No means that the paper has limitations, but those are not discussed in the paper. 
        \item The authors are encouraged to create a separate "Limitations" section in their paper.
        \item The paper should point out any strong assumptions and how robust the results are to violations of these assumptions (e.g., independence assumptions, noiseless settings, model well-specification, asymptotic approximations only holding locally). The authors should reflect on how these assumptions might be violated in practice and what the implications would be.
        \item The authors should reflect on the scope of the claims made, e.g., if the approach was only tested on a few datasets or with a few runs. In general, empirical results often depend on implicit assumptions, which should be articulated.
        \item The authors should reflect on the factors that influence the performance of the approach. For example, a facial recognition algorithm may perform poorly when image resolution is low or images are taken in low lighting. Or a speech-to-text system might not be used reliably to provide closed captions for online lectures because it fails to handle technical jargon.
        \item The authors should discuss the computational efficiency of the proposed algorithms and how they scale with dataset size.
        \item If applicable, the authors should discuss possible limitations of their approach to address problems of privacy and fairness.
        \item While the authors might fear that complete honesty about limitations might be used by reviewers as grounds for rejection, a worse outcome might be that reviewers discover limitations that aren't acknowledged in the paper. The authors should use their best judgment and recognize that individual actions in favor of transparency play an important role in developing norms that preserve the integrity of the community. Reviewers will be specifically instructed to not penalize honesty concerning limitations.
    \end{itemize}

\item {\bf Theory Assumptions and Proofs}
    \item[] Question: For each theoretical result, does the paper provide the full set of assumptions and a complete (and correct) proof?
    \item[] Answer: \answerYes{} % Replace by \answerYes{}, \answerNo{}, or \answerNA{}.
    \item[] Justification: Our paper includes theoretical results in Section \ref{sec:mainreults}. In our theorems, we clearly mentioned all the required assumptions, and a complete (and correct) proof of them is available in appendices (e.g., see Appendix \ref{app:theorem_proofs}). Please see Section \ref{sec:framework} for a full definition of the problem and introduction to the notations used in the paper. 
    \item[] Guidelines:
    \begin{itemize}
        \item The answer NA means that the paper does not include theoretical results. 
        \item All the theorems, formulas, and proofs in the paper should be numbered and cross-referenced.
        \item All assumptions should be clearly stated or referenced in the statement of any theorems.
        \item The proofs can either appear in the main paper or the supplemental material, but if they appear in the supplemental material, the authors are encouraged to provide a short proof sketch to provide intuition. 
        \item Inversely, any informal proof provided in the core of the paper should be complemented by formal proofs provided in appendix or supplemental material.
        \item Theorems and Lemmas that the proof relies upon should be properly referenced. 
    \end{itemize}

    \item {\bf Experimental Result Reproducibility}
    \item[] Question: Does the paper fully disclose all the information needed to reproduce the main experimental results of the paper to the extent that it affects the main claims and/or conclusions of the paper (regardless of whether the code and data are provided or not)?
    \item[] Answer: \answerYes{} % Replace by \answerYes{}, \answerNo{}, or \answerNA{}.
    \item[] Justification:
    We completely explained our proposed framework in Section \ref{sec:framework} and we also provided details regarding our empirical evaluations in Section \ref{sec:framework} in the paper. 
    \item[] Guidelines:
    \begin{itemize}
        \item The answer NA means that the paper does not include experiments.
        \item If the paper includes experiments, a No answer to this question will not be perceived well by the reviewers: Making the paper reproducible is important, regardless of whether the code and data are provided or not.
        \item If the contribution is a dataset and/or model, the authors should describe the steps taken to make their results reproducible or verifiable. 
        \item Depending on the contribution, reproducibility can be accomplished in various ways. For example, if the contribution is a novel architecture, describing the architecture fully might suffice, or if the contribution is a specific model and empirical evaluation, it may be necessary to either make it possible for others to replicate the model with the same dataset, or provide access to the model. In general. releasing code and data is often one good way to accomplish this, but reproducibility can also be provided via detailed instructions for how to replicate the results, access to a hosted model (e.g., in the case of a large language model), releasing of a model checkpoint, or other means that are appropriate to the research performed.
        \item While NeurIPS does not require releasing code, the conference does require all submissions to provide some reasonable avenue for reproducibility, which may depend on the nature of the contribution. For example
        \begin{enumerate}
            \item If the contribution is primarily a new algorithm, the paper should make it clear how to reproduce that algorithm.
            \item If the contribution is primarily a new model architecture, the paper should describe the architecture clearly and fully.
            \item If the contribution is a new model (e.g., a large language model), then there should either be a way to access this model for reproducing the results or a way to reproduce the model (e.g., with an open-source dataset or instructions for how to construct the dataset).
            \item We recognize that reproducibility may be tricky in some cases, in which case authors are welcome to describe the particular way they provide for reproducibility. In the case of closed-source models, it may be that access to the model is limited in some way (e.g., to registered users), but it should be possible for other researchers to have some path to reproducing or verifying the results.
        \end{enumerate}
    \end{itemize}


\item {\bf Open access to data and code}
    \item[] Question: Does the paper provide open access to the data and code, with sufficient instructions to faithfully reproduce the main experimental results, as described in supplemental material?
    \item[] Answer: \answerNo{} % Replace by \answerYes{}, \answerNo{}, or \answerNA{}.
    \item[] Justification: We provided references to all the open datasets that we used in the paper. Regarding the code, we are happy to share it later if required.
    \item[] Guidelines:
    \begin{itemize}
        \item The answer NA means that paper does not include experiments requiring code.
        \item Please see the NeurIPS code and data submission guidelines (\url{https://nips.cc/public/guides/CodeSubmissionPolicy}) for more details.
        \item While we encourage the release of code and data, we understand that this might not be possible, so “No” is an acceptable answer. Papers cannot be rejected simply for not including code, unless this is central to the contribution (e.g., for a new open-source benchmark).
        \item The instructions should contain the exact command and environment needed to run to reproduce the results. See the NeurIPS code and data submission guidelines (\url{https://nips.cc/public/guides/CodeSubmissionPolicy}) for more details.
        \item The authors should provide instructions on data access and preparation, including how to access the raw data, preprocessed data, intermediate data, and generated data, etc.
        \item The authors should provide scripts to reproduce all experimental results for the new proposed method and baselines. If only a subset of experiments are reproducible, they should state which ones are omitted from the script and why.
        \item At submission time, to preserve anonymity, the authors should release anonymized versions (if applicable).
        \item Providing as much information as possible in supplemental material (appended to the paper) is recommended, but including URLs to data and code is permitted.
    \end{itemize}


\item {\bf Experimental Setting/Details}
    \item[] Question: Does the paper specify all the training and test details (e.g., data splits, hyperparameters, how they were chosen, type of optimizer, etc.) necessary to understand the results?
    \item[] Answer: \answerYes{} % Replace by \answerYes{}, \answerNo{}, or \answerNA{}.
    \item[] Justification: We provided full experimental details in the paper (see Section \ref{sec:exp_result}).
    \item[] Guidelines:
    \begin{itemize}
        \item The answer NA means that the paper does not include experiments.
        \item The experimental setting should be presented in the core of the paper to a level of detail that is necessary to appreciate the results and make sense of them.
        \item The full details can be provided either with the code, in appendix, or as supplemental material.
    \end{itemize}

\item {\bf Experiment Statistical Significance}
    \item[] Question: Does the paper report error bars suitably and correctly defined or other appropriate information about the statistical significance of the experiments?
    \item[] Answer: \answerYes{}  % Replace by \answerYes{}, \answerNo{}, or \answerNA{}.
    \item[] Justification: The details are provided in Section~\ref{sec:exp_result}.
    \item[] Guidelines:
    \begin{itemize}
        \item The answer NA means that the paper does not include experiments.
        \item The authors should answer "Yes" if the results are accompanied by error bars, confidence intervals, or statistical significance tests, at least for the experiments that support the main claims of the paper.
        \item The factors of variability that the error bars are capturing should be clearly stated (for example, train/test split, initialization, random drawing of some parameter, or overall run with given experimental conditions).
        \item The method for calculating the error bars should be explained (closed form formula, call to a library function, bootstrap, etc.)
        \item The assumptions made should be given (e.g., Normally distributed errors).
        \item It should be clear whether the error bar is the standard deviation or the standard error of the mean.
        \item It is OK to report 1-sigma error bars, but one should state it. The authors should preferably report a 2-sigma error bar than state that they have a 96\% CI, if the hypothesis of Normality of errors is not verified.
        \item For asymmetric distributions, the authors should be careful not to show in tables or figures symmetric error bars that would yield results that are out of range (e.g. negative error rates).
        \item If error bars are reported in tables or plots, The authors should explain in the text how they were calculated and reference the corresponding figures or tables in the text.
    \end{itemize}

\item {\bf Experiments Compute Resources}
    \item[] Question: For each experiment, does the paper provide sufficient information on the computer resources (type of compute workers, memory, time of execution) needed to reproduce the experiments?
    \item[] Answer: \answerYes{} % Replace by \answerYes{}, \answerNo{}, or \answerNA{}.
    \item[] Justification: We provided all the details regarding our experiments in Section~\ref{sec:exp_result}. 
    \item[] Guidelines:
    \begin{itemize}
        \item The answer NA means that the paper does not include experiments.
        \item The paper should indicate the type of compute workers CPU or GPU, internal cluster, or cloud provider, including relevant memory and storage.
        \item The paper should provide the amount of compute required for each of the individual experimental runs as well as estimate the total compute. 
        \item The paper should disclose whether the full research project required more compute than the experiments reported in the paper (e.g., preliminary or failed experiments that didn't make it into the paper). 
    \end{itemize}
    
\item {\bf Code Of Ethics}
    \item[] Question: Does the research conducted in the paper conform, in every respect, with the NeurIPS Code of Ethics \url{https://neurips.cc/public/EthicsGuidelines}?
    \item[] Answer: \answerYes{} % Replace by \answerYes{}, \answerNo{}, or \answerNA{}.
    \item[] Justification: We followed the NeurIPS code of ethics in our paper. 
    \item[] Guidelines:
    \begin{itemize}
        \item The answer NA means that the authors have not reviewed the NeurIPS Code of Ethics.
        \item If the authors answer No, they should explain the special circumstances that require a deviation from the Code of Ethics.
        \item The authors should make sure to preserve anonymity (e.g., if there is a special consideration due to laws or regulations in their jurisdiction).
    \end{itemize}


\item {\bf Broader Impacts}
    \item[] Question: Does the paper discuss both potential positive societal impacts and negative societal impacts of the work performed?
    \item[] Answer: \answerNA{} % Replace by \answerYes{}, \answerNo{}, or \answerNA{}.
    \item[] Justification: Our paper is focused on developing a new framework for coded distributed computing and should be categorized as foundational research. We believe this work has no direct societal impact that should be explained in the paper.
    \item[] Guidelines:
    \begin{itemize}
        \item The answer NA means that there is no societal impact of the work performed.
        \item If the authors answer NA or No, they should explain why their work has no societal impact or why the paper does not address societal impact.
        \item Examples of negative societal impacts include potential malicious or unintended uses (e.g., disinformation, generating fake profiles, surveillance), fairness considerations (e.g., deployment of technologies that could make decisions that unfairly impact specific groups), privacy considerations, and security considerations.
        \item The conference expects that many papers will be foundational research and not tied to particular applications, let alone deployments. However, if there is a direct path to any negative applications, the authors should point it out. For example, it is legitimate to point out that an improvement in the quality of generative models could be used to generate deepfakes for disinformation. On the other hand, it is not needed to point out that a generic algorithm for optimizing neural networks could enable people to train models that generate Deepfakes faster.
        \item The authors should consider possible harms that could arise when the technology is being used as intended and functioning correctly, harms that could arise when the technology is being used as intended but gives incorrect results, and harms following from (intentional or unintentional) misuse of the technology.
        \item If there are negative societal impacts, the authors could also discuss possible mitigation strategies (e.g., gated release of models, providing defenses in addition to attacks, mechanisms for monitoring misuse, mechanisms to monitor how a system learns from feedback over time, improving the efficiency and accessibility of ML).
    \end{itemize}
    
\item {\bf Safeguards}
    \item[] Question: Does the paper describe safeguards that have been put in place for responsible release of data or models that have a high risk for misuse (e.g., pretrained language models, image generators, or scraped datasets)?
    \item[] Answer: \answerNA{} % Replace by \answerYes{}, \answerNo{}, or \answerNA{}.
    \item[] Justification: This is not applicable to our work and this paper poses no such risks.
    \item[] Guidelines:
    \begin{itemize}
        \item The answer NA means that the paper poses no such risks.
        \item Released models that have a high risk for misuse or dual-use should be released with necessary safeguards to allow for controlled use of the model, for example by requiring that users adhere to usage guidelines or restrictions to access the model or implementing safety filters. 
        \item Datasets that have been scraped from the Internet could pose safety risks. The authors should describe how they avoided releasing unsafe images.
        \item We recognize that providing effective safeguards is challenging, and many papers do not require this, but we encourage authors to take this into account and make a best faith effort.
    \end{itemize}

\item {\bf Licenses for existing assets}
    \item[] Question: Are the creators or original owners of assets (e.g., code, data, models), used in the paper, properly credited and are the license and terms of use explicitly mentioned and properly respected?
    \item[] Answer:  \answerNA{} % Replace by \answerYes{}, \answerNo{}, or \answerNA{}.
    \item[] Justification: This paper does not use existing assets.
    \item[] Guidelines:
    \begin{itemize}
        \item The answer NA means that the paper does not use existing assets.
        \item The authors should cite the original paper that produced the code package or dataset.
        \item The authors should state which version of the asset is used and, if possible, include a URL.
        \item The name of the license (e.g., CC-BY 4.0) should be included for each asset.
        \item For scraped data from a particular source (e.g., website), the copyright and terms of service of that source should be provided.
        \item If assets are released, the license, copyright information, and terms of use in the package should be provided. For popular datasets, \url{paperswithcode.com/datasets} has curated licenses for some datasets. Their licensing guide can help determine the license of a dataset.
        \item For existing datasets that are re-packaged, both the original license and the license of the derived asset (if it has changed) should be provided.
        \item If this information is not available online, the authors are encouraged to reach out to the asset's creators.
    \end{itemize}

\item {\bf New Assets}
    \item[] Question: Are new assets introduced in the paper well documented and is the documentation provided alongside the assets?
    \item[] Answer:  \answerNA{} % Replace by \answerYes{}, \answerNo{}, or \answerNA{}.
    \item[] Justification: This paper does not release new assets.
    \item[] Guidelines:
    \begin{itemize}
        \item The answer NA means that the paper does not release new assets.
        \item Researchers should communicate the details of the dataset/code/model as part of their submissions via structured templates. This includes details about training, license, limitations, etc. 
        \item The paper should discuss whether and how consent was obtained from people whose asset is used.
        \item At submission time, remember to anonymize your assets (if applicable). You can either create an anonymized URL or include an anonymized zip file.
    \end{itemize}

\item {\bf Crowdsourcing and Research with Human Subjects}
    \item[] Question: For crowdsourcing experiments and research with human subjects, does the paper include the full text of instructions given to participants and screenshots, if applicable, as well as details about compensation (if any)? 
    \item[] Answer: \answerNA{} % Replace by \answerYes{}, \answerNo{}, or \answerNA{}.
    \item[] Justification: Our paper does not involve these.
    \item[] Guidelines:
    \begin{itemize}
        \item The answer NA means that the paper does not involve crowdsourcing nor research with human subjects.
        \item Including this information in the supplemental material is fine, but if the main contribution of the paper involves human subjects, then as much detail as possible should be included in the main paper. 
        \item According to the NeurIPS Code of Ethics, workers involved in data collection, curation, or other labor should be paid at least the minimum wage in the country of the data collector. 
    \end{itemize}

\item {\bf Institutional Review Board (IRB) Approvals or Equivalent for Research with Human Subjects}
    \item[] Question: Does the paper describe potential risks incurred by study participants, whether such risks were disclosed to the subjects, and whether Institutional Review Board (IRB) approvals (or an equivalent approval/review based on the requirements of your country or institution) were obtained?
    \item[] Answer: \answerNA{} % Replace by \answerYes{}, \answerNo{}, or \answerNA{}.
    \item[] Justification: Our paper does not involve these.
    \item[] Guidelines:
    \begin{itemize}
        \item The answer NA means that the paper does not involve crowdsourcing nor research with human subjects.
        \item Depending on the country in which research is conducted, IRB approval (or equivalent) may be required for any human subjects research. If you obtained IRB approval, you should clearly state this in the paper. 
        \item We recognize that the procedures for this may vary significantly between institutions and locations, and we expect authors to adhere to the NeurIPS Code of Ethics and the guidelines for their institution. 
        \item For initial submissions, do not include any information that would break anonymity (if applicable), such as the institution conducting the review.
    \end{itemize}

\end{enumerate}






\end{document}
